% Alternating Current — Physics Upper Sixth — Cours signé OnBuch+
% Official MINESEC syllabus. Compiler avec : tectonic PHY09-alternating-current.tex
\newcommand{\DOCMATIERE}{Physics}
\newcommand{\DOCNIVEAU}{Upper Sixth}
\newcommand{\DOCMODULE}{Upper Sixth — Physics}
\newcommand{\DOCLECON}{9}
\newcommand{\DOCTITRE}{Alternating Current}
% OnBuch+ — préambule « cours signés » ENRICHI (compilé avec tectonic / XeLaTeX)
% Système visuel « Pop » d'OnBuch+ : fond crème (jamais blanc), bordures encre
% épaisses, ombres « dures » décalées, titres Archivo Black en capitales.
% Version MATHÉMATIQUES : graphes (pgfplots/tikz), boîtes pédagogiques
% (définition, propriété, méthode, attention, le savais-tu, exercice résolu,
% exercice, TP, bilan), page de garde et signature OnBuch+ ancrée en bas.
%
% Macros attendues dans main.tex AVANT \input{preamble} :
%   \DOCMATIERE  \DOCNIVEAU  \DOCMODULE  \DOCLECON (numéro)  \DOCTITRE
\documentclass[11pt]{article}

\usepackage{fontspec}
\usepackage[english]{babel}
\usepackage[a4paper,margin=2cm,top=3.4cm,bottom=2.8cm,headheight=1.4cm,footskip=1.3cm]{geometry}
\usepackage{amsmath,amssymb,amsfonts}
\usepackage{mathtools} % \xmapsto, \coloneqq... souvent utilisés par les modèles
\usepackage{cancel} % simplifications barrées (bilans de forces)
% \abs{z} (module, valeur absolue) et \norm{u} (norme), utilisés par les modèles
\providecommand{\abs}{}\let\abs\relax\DeclarePairedDelimiter{\abs}{\lvert}{\rvert}
\providecommand{\norm}{}\let\norm\relax\DeclarePairedDelimiter{\norm}{\lVert}{\rVert}
\usepackage[table]{xcolor} % [table] : fournit aussi \rowcolors (alternance de lignes)
\usepackage{pagecolor}
\usepackage{tikz}
\usetikzlibrary{arrows.meta,positioning,calc,shapes.geometric,shapes.misc,shapes.arrows,decorations.pathmorphing,decorations.pathreplacing,decorations.markings,patterns,shadows,mindmap,trees,fit,backgrounds,matrix,intersections,angles,quotes}
\usepackage{pgfplots}
\usepackage[version=4]{mhchem} % équations nucléaires \ce{^{235}_{92}U}
\usepackage[european,siunitx]{circuitikz} % schémas électriques
\pgfplotsset{compat=1.18}
\usepgfplotslibrary{fillbetween,groupplots} % aires entre courbes (calcul intégral)
\usepackage{enumitem}
\usepackage{fancyhdr}
% [most] charge toutes les bibliothèques tcolorbox (listings, minted, raster,
% documentation...) et gonfle inutilement la mémoire TeX sur un document long
% (~300 boîtes) : on ne charge que ce qui est réellement utilisé.
\usepackage[skins,breakable]{tcolorbox}
\usepackage{graphicx}
\usepackage{colortbl}
\usepackage{booktabs}
\usepackage{tabularx}
\usepackage{diagbox} % case d'en-tête diagonale des tableaux à double entrée
% Colonnes extensibles centrée / gauche / droite (C, L, R), souvent utilisées par les modèles
\newcolumntype{C}{>{\centering\arraybackslash}X}
\newcolumntype{L}{>{\raggedright\arraybackslash}X}
\newcolumntype{R}{>{\raggedleft\arraybackslash}X}
\usepackage{array}
\usepackage{multicol}
\usepackage{siunitx}
% parse-numbers=false : accepte aussi \SI{2,5 \times 10^{-3}}{...}
\sisetup{locale=UK,output-decimal-marker={.},group-minimum-digits=4,parse-numbers=false}
\DeclareSIUnit\ohm{\ensuremath{\symup{\Omega}}} % U+2126 absent de la police
\DeclareSIUnit\division{div} % réglages d'oscilloscope (ms/div, V/div)
\DeclareSIUnit\tr{tr} % tours (vitesse de rotation, ex. \SI{1500}{\tr\per\minute})
\usepackage{qrcode}
% \degree : commande fréquemment utilisée par les modèles pour un angle ou une
% température en dehors de \SI{}{} (non fournie par les paquets chargés).
\providecommand{\degree}{^{\circ}}
% \qed (fin de démonstration), utilisé par les modèles sans amsthm
\providecommand{\qed}{\hfill$\square$}
% \barre : double barre des valeurs interdites dans un tableau de variations (tkz-tab)
\providecommand{\barre}{\ensuremath{\|}}
% environnement proof (amsthm non chargé) utilisé par les modèles
\ifdefined\proof\else\newenvironment{proof}[1][Proof]{\par\noindent\textit{#1.}\ }{\qed\par}\fi
\usepackage{lastpage}
\usepackage{hyperref}
\hypersetup{hidelinks}

% ============================================================
% Palette « Pop » OnBuch+ — mêmes hex que la classe P (plus_ui.dart)
% ============================================================
\definecolor{poporange}{HTML}{F59321}
\definecolor{popdark}{HTML}{E07A0C}
\definecolor{popcream}{HTML}{FFF6E8}
\definecolor{popink}{HTML}{1C1714}
\definecolor{poppurple}{HTML}{7A5AE0}
\definecolor{popgreen}{HTML}{2E9E6A}
\definecolor{poppink}{HTML}{E0457B}
\definecolor{popblue}{HTML}{2F7DE1}
\definecolor{popgold}{HTML}{C98A12}
\definecolor{popmuted}{HTML}{8A7F76}
\definecolor{popline}{HTML}{EADFD2}
\definecolor{popgray}{HTML}{8A7F76}
\colorlet{popgrey}{popgray}
% Alias tolérants (le modèle nomme parfois les couleurs différemment)
\colorlet{popwhite}{white}
\colorlet{popblack}{popink}
\colorlet{popred}{poppink}
\colorlet{popyellow}{popgold}
% Teintes claires pour fonds de tableaux / zones
\colorlet{popblueL}{popblue!10!white}
\colorlet{poporangeL}{poporange!14!white}
\colorlet{popgreenL}{popgreen!12!white}
\colorlet{poppurpleL}{poppurple!12!white}
\colorlet{poppinkL}{poppink!12!white}
\colorlet{popgoldL}{popgold!14!white}

\pagecolor{popcream}

% ============================================================
% Typographie
% ============================================================
\defaultfontfeatures{Ligatures=TeX}
\setmainfont{PlusJakartaSans-Medium.ttf}[Path=../../../fonts/,BoldFont=PlusJakartaSans-Bold.ttf]
% Maths en sans-serif (Fira Math) avec chiffres et lettres droites de Plus
% Jakarta, pour que les formules se fondent dans le texte.
\usepackage[math-style=ISO,bold-style=ISO]{unicode-math}
\setmathfont{FiraMath-Regular.otf}[Path=../../../fonts/]
\setmathfont{PlusJakartaSans-Medium.ttf}[Path=../../../fonts/,range={up/num,up/{Latin,latin},\mathup}]
% (icomma retiré : décimales avec point en anglais)
% Caractères mathématiques Unicode absents des polices de texte (Plus Jakarta,
% Archivo Black) : rendus par leur équivalent mathématique, en texte comme en
% formule — sinon ils s'affichent en carré vide (ex. « Arithmétique dans ℤ »).
\usepackage{newunicodechar}
\newunicodechar{ℝ}{\ensuremath{\mathbb{R}}}
\newunicodechar{ℤ}{\ensuremath{\mathbb{Z}}}
\newunicodechar{ℂ}{\ensuremath{\mathbb{C}}}
\newunicodechar{ℕ}{\ensuremath{\mathbb{N}}}
\newunicodechar{ℚ}{\ensuremath{\mathbb{Q}}}
\newunicodechar{𝔻}{\ensuremath{\mathbb{D}}}
\newunicodechar{α}{\ensuremath{\alpha}}
\newunicodechar{β}{\ensuremath{\beta}}
\newunicodechar{γ}{\ensuremath{\gamma}}
\newunicodechar{δ}{\ensuremath{\delta}}
\newunicodechar{ε}{\ensuremath{\varepsilon}}
\newunicodechar{θ}{\ensuremath{\theta}}
\newunicodechar{λ}{\ensuremath{\lambda}}
\newunicodechar{μ}{\ensuremath{\mu}}
\newunicodechar{σ}{\ensuremath{\sigma}}
\newunicodechar{τ}{\ensuremath{\tau}}
\newunicodechar{φ}{\ensuremath{\varphi}}
\newunicodechar{ψ}{\ensuremath{\psi}}
\newunicodechar{ω}{\ensuremath{\omega}}
\newunicodechar{Σ}{\ensuremath{\Sigma}}
% majuscules produites par \MakeUppercase dans les titres (α-aminés -> Α)
\newunicodechar{Α}{\ensuremath{\alpha}}
\newunicodechar{Β}{\ensuremath{\beta}}
\newunicodechar{Γ}{\ensuremath{\Gamma}}
\newunicodechar{Δ}{\ensuremath{\Delta}}
\newunicodechar{Ε}{\ensuremath{\varepsilon}}
\newunicodechar{Θ}{\ensuremath{\Theta}}
\newunicodechar{Λ}{\ensuremath{\Lambda}}
\newunicodechar{Μ}{\ensuremath{\mu}}
\newunicodechar{Τ}{\ensuremath{\tau}}
\newunicodechar{Φ}{\ensuremath{\Phi}}
\newunicodechar{Ψ}{\ensuremath{\Psi}}
\newunicodechar{Ω}{\ensuremath{\Omega}}
\newunicodechar{↦}{\ensuremath{\mapsto}}
\newunicodechar{∈}{\ensuremath{\in}}
\newunicodechar{∉}{\ensuremath{\notin}}
\newunicodechar{∧}{\ensuremath{\wedge}}
\newunicodechar{⇔}{\ensuremath{\Leftrightarrow}}
\newunicodechar{⟺}{\ensuremath{\Longleftrightarrow}}
\newunicodechar{⇒}{\ensuremath{\Rightarrow}}
\newunicodechar{⟹}{\ensuremath{\Longrightarrow}}
\newunicodechar{∘}{\ensuremath{\circ}}
\newunicodechar{∪}{\ensuremath{\cup}}
\newunicodechar{∩}{\ensuremath{\cap}}
\newunicodechar{≡}{\ensuremath{\equiv}}
\newunicodechar{⊂}{\ensuremath{\subset}}
\newunicodechar{⊥}{\ensuremath{\perp}}
\newunicodechar{∃}{\ensuremath{\exists}}
\newunicodechar{∀}{\ensuremath{\forall}}
\newunicodechar{∛}{\ensuremath{\sqrt[3]{\,}}}
\newunicodechar{∜}{\ensuremath{\sqrt[4]{\,}}}
\newunicodechar{⃗}{\ensuremath{{}^{\to}}}
\newunicodechar{ⁿ}{\ifmmode^{n}\else\textsuperscript{\ensuremath{n}}\fi}
\newunicodechar{ˣ}{\ifmmode^{x}\else\textsuperscript{\ensuremath{x}}\fi}
\newunicodechar{ᵇ}{\ifmmode^{b}\else\textsuperscript{\ensuremath{b}}\fi}
\newunicodechar{ʳ}{\ifmmode^{r}\else\textsuperscript{\ensuremath{r}}\fi}
\newunicodechar{ᵘ}{\ifmmode^{u}\else\textsuperscript{\ensuremath{u}}\fi}
\newunicodechar{ʸ}{\ifmmode^{y}\else\textsuperscript{\ensuremath{y}}\fi}
\newunicodechar{ᵃ}{\ifmmode^{a}\else\textsuperscript{\ensuremath{a}}\fi}
\newunicodechar{ᵉ}{\ifmmode^{e}\else\textsuperscript{\ensuremath{e}}\fi}
\newunicodechar{ᵏ}{\ifmmode^{k}\else\textsuperscript{\ensuremath{k}}\fi}
\newunicodechar{ᵖ}{\ifmmode^{p}\else\textsuperscript{\ensuremath{p}}\fi}
\newunicodechar{ᴺ}{\ifmmode^{N}\else\textsuperscript{\ensuremath{N}}\fi}
\newunicodechar{ᶜ}{\ifmmode^{c}\else\textsuperscript{\ensuremath{c}}\fi}
\newunicodechar{ʰ}{\ifmmode^{h}\else\textsuperscript{\ensuremath{h}}\fi}
\newunicodechar{ᵠ}{\ifmmode^{q}\else\textsuperscript{\ensuremath{q}}\fi}
\newunicodechar{ₙ}{\ifmmode_{n}\else\textsubscript{\ensuremath{n}}\fi}
\newunicodechar{ₐ}{\ifmmode_{a}\else\textsubscript{\ensuremath{a}}\fi}
\newunicodechar{ᵢ}{\ifmmode_{i}\else\textsubscript{\ensuremath{i}}\fi}
\newunicodechar{ᵣ}{\ifmmode_{r}\else\textsubscript{\ensuremath{r}}\fi}
\newunicodechar{ₖ}{\ifmmode_{k}\else\textsubscript{\ensuremath{k}}\fi}
\newunicodechar{ᵦ}{\ifmmode_{\beta}\else\textsubscript{\ensuremath{\beta}}\fi}
\newunicodechar{ₓ}{\ifmmode_{x}\else\textsubscript{\ensuremath{x}}\fi}
\newfontfamily\popdisplay{ArchivoBlack-Regular.ttf}[Path=../../../fonts/]
\newfontfamily\poplogo{SpaceGrotesk-Bold.ttf}[Path=../../../fonts/]
\newfontfamily\popsemi{PlusJakartaSans-SemiBold.ttf}[Path=../../../fonts/]
\newfontfamily\popxbold{PlusJakartaSans-ExtraBold.ttf}[Path=../../../fonts/]
\newfontfamily\popxboldls{PlusJakartaSans-ExtraBold.ttf}[Path=../../../fonts/,LetterSpace=3.0]

\setlist[enumerate]{leftmargin=1.7em,itemsep=5pt,topsep=4pt}
\setlist[itemize]{leftmargin=1.7em,itemsep=4pt,topsep=4pt,label={\color{poporange}\textbullet}}
\setlength{\parindent}{0pt}
\setlength{\parskip}{5pt}
\renewcommand{\arraystretch}{1.25}

% pgfplots : style Pop par défaut
\pgfplotsset{
  every axis/.append style={axis line style={popink,line width=1pt},tick style={popink},
    grid style={popline,line width=0.5pt},label style={font=\small},tick label style={font=\footnotesize}},
  popcurve/.style={poporange,line width=1.8pt,no markers,smooth},
}
% Même style disponible dans un \draw TikZ simple (hors pgfplots)
\tikzset{popcurve/.style={poporange,line width=1.8pt}}
\tikzset{popgrid/.style={popline,very thin}}
\colorlet{popcurve}{poporange} % le nom du style est parfois utilisé comme couleur

% ============================================================
% Pill / squiggle
% ============================================================
\newcommand{\poppill}[3]{%
  \tikz[baseline=-0.55ex]\node[draw=popink,line width=1.2pt,rounded corners=2.6pt,%
    fill=#2,inner xsep=6.5pt,inner ysep=3pt]{%
    \popxboldls\fontsize{7}{7}\selectfont\color{#3}\MakeUppercase{#1}};%
}
\newcommand{\popsquiggle}[1]{%
  \tikz[baseline=-0.4ex]{
    \draw[#1,line width=2.6pt,line cap=round]
      plot [smooth] coordinates {(0,0.09) (0.34,0.01) (0.68,0.21) (1.02,0.01) (1.36,0.10)};
  }%
}
% Mot-clé surligné (marqueur orange sous le texte)
\newcommand{\cle}[1]{{\popxbold\color{popdark}#1}}

% ============================================================
% En-tête / pied
% ============================================================
\pagestyle{fancy}
\fancyhf{}
\renewcommand{\headrulewidth}{0pt}
\renewcommand{\footrulewidth}{0pt}
\lhead{\raisebox{-0.15ex}{\poplogo\fontsize{13}{13}\selectfont\color{popink}OnBuch\color{poporange}+}}
\rhead{\poppill{\DOCMATIERE\ \textbullet\ \DOCNIVEAU\ \textbullet\ Chapter \DOCLECON}{white}{popink}}
\lfoot{\popxbold\fontsize{7.6}{7.6}\selectfont\color{popmuted}\MakeUppercase{Signed course OnBuch+}\ \textbullet\ {\popsemi\DOCTITRE}}
\rfoot{\tikz[baseline=-0.6ex]\node[rounded corners=8.5pt,fill=popink,minimum height=17pt,minimum width=17pt,inner xsep=5pt,inner sep=0]{\popxbold\fontsize{8}{8}\selectfont\color{popcream}\thepage\,/\,\pageref*{LastPage}};}

% ============================================================
% Titres
% ============================================================
\newcommand{\chaptitre}[2]{%
  \begin{tcolorbox}[enhanced,colback=poporange,colframe=popink,boxrule=2.6pt,arc=11pt,
      drop shadow=popink,left=15pt,right=15pt,top=12pt,bottom=11pt,before skip=3pt,after skip=18pt]
    \poppill{#2}{popink}{popcream}\par\vspace{9pt}
    {\raggedright\hyphenpenalty=10000\exhyphenpenalty=10000\popdisplay\fontsize{22}{24}\selectfont\color{popcream}\MakeUppercase{#1}\par}
  \end{tcolorbox}
}
% Section numérotée : pastille orange avec le numéro + titre Archivo
\newcounter{popsec}
\newcommand{\coursec}[1]{%
  \stepcounter{popsec}\setcounter{popsub}{0}%
  \par\vspace{16pt}\noindent
  \begin{minipage}{\linewidth}
  \tikz[baseline=-0.7ex]\node[draw=popink,line width=1.6pt,fill=poporange,rounded corners=4pt,minimum size=22pt,inner sep=0]{\popdisplay\fontsize{12}{12}\selectfont\color{popcream}\thepopsec};%
  \hspace{8pt}{\popdisplay\fontsize{15}{17}\selectfont\color{popink}\MakeUppercase{#1}}\par
  \vspace{4pt}\noindent{\color{poporange}\rule{36pt}{3.6pt}}
  \end{minipage}\par\nopagebreak\vspace{8pt}
}
\newcounter{popsub}
\newcommand{\courssub}[1]{%
  \stepcounter{popsub}%
  \par\vspace{9pt}\noindent{\popxbold\fontsize{11.5}{13}\selectfont\color{popink}{\color{poporange}\thepopsec.\thepopsub}\hspace{6pt}#1}\par\nopagebreak\vspace{4pt}
}

% ============================================================
% Boîtes pédagogiques — même recette : fond blanc, bordure épaisse colorée,
% ombre dure encre, badge plein. Titre optionnel [#1].
% ============================================================
\tcbset{popbase/.style={breakable,enhanced,colback=white,boxrule=2.3pt,arc=9pt,
  shadow={3pt}{-3pt}{0pt}{popink},left=14pt,right=14pt,top=11pt,bottom=11pt,before skip=11pt,after skip=15pt,
  fontupper=\popsemi\fontsize{10.6}{14.5}\selectfont\color{popink},
  fontlower=\popsemi\fontsize{10.6}{14.5}\selectfont\color{popink}}}
% Coche dessinée (le glyphe ✓ n'existe pas dans les polices du document)
\newcommand{\popcheck}[1]{\tikz[baseline=-0.5ex]\draw[#1,line width=1.6pt,line cap=round,line join=round](-0.13,0.0)--(-0.04,-0.09)--(0.13,0.1);}
\providecommand{\coche}{\popcheck{popgreen}}
\providecommand{\resultat}[1]{\par\smallskip\noindent\fcolorbox{popgreen}{popgreenL}{\parbox{\dimexpr\linewidth-2\fboxsep-2\fboxrule\relax}{\textbf{Result.} #1}}\par\smallskip}
\newcommand{\popboxhead}[3]{\poppill{#1}{#2}{white}\ifx\relax#3\relax\else\hspace{6pt}{\popxbold\fontsize{10.6}{12}\selectfont\color{popink}#3}\fi\par\vspace{7pt}}

\newtcolorbox{aretenir}[1][]{popbase,colframe=popgreen,before upper={\popboxhead{Key points}{popgreen}{#1}}}
\newtcolorbox{exemplebox}[1][]{popbase,colframe=poporange,before upper={\popboxhead{Exemple}{poporange}{#1}}}
\newtcolorbox{definition}[1][]{popbase,colframe=popblue,before upper={\popboxhead{Definition}{popblue}{#1}}}
\newtcolorbox{propriete}[1][]{popbase,colframe=poppurple,before upper={\popboxhead{Property}{poppurple}{#1}}}
\newtcolorbox{methode}[1][]{popbase,colframe=popgold,before upper={\popboxhead{Method}{popgold}{#1}}}
\newtcolorbox{attention}[1][]{popbase,colframe=poppink,before upper={\popboxhead{Warning}{poppink}{#1}}}
\newtcolorbox{experience}[1][]{popbase,colframe=popblue,colback=popblueL,before upper={\popboxhead{Experiment}{popblue}{#1}}}
% « Le savais-tu ? » : carte violette pleine (contexte, culture, Cameroun)
\newtcolorbox{savaistu}[1][]{popbase,colback=poppurple,colframe=popink,
  fontupper=\popsemi\fontsize{10.4}{14.2}\selectfont\color{white},
  before upper={\poppill{Did you know?}{white}{poppurple}\ifx\relax#1\relax\else\hspace{6pt}{\popxbold\color{white}#1}\fi\par\vspace{7pt}}}
% Exercice résolu : énoncé en haut, \tcblower puis la solution sur fond crème
\newcounter{popexo}\newcounter{popexr}
\newtcolorbox{exoresolu}[1][]{popbase,colframe=poporange,colbacklower=popcream,
  segmentation style={popink,line width=1.2pt,dashed},
  before upper={\stepcounter{popexr}\popboxhead{Worked example \thepopexr}{poporange}{#1}},
  before lower={\poppill{Solution}{popink}{popcream}\par\vspace{6pt}}}
% Exercice (non résolu en place) — la correction est dans la partie Corrigés
\newtcolorbox{exercice}[1][]{popbase,colframe=popink,
  before upper={\stepcounter{popexo}\popboxhead{Exercise \thepopexo}{popink}{#1}}}
% Corrigé
\newtcolorbox{corrige}[1][]{popbase,colframe=popgreen,colback=popgreenL,
  before upper={\popboxhead{Solution}{popgreen}{#1}}}
% Objectifs / prérequis / bilan
\newtcolorbox{objectifs}{popbase,colframe=popink,colback=poporangeL,
  before upper={\popboxhead{Chapter objectives}{popink}{}}}
\newtcolorbox{prerequis}{popbase,colframe=popink,colback=popblueL,
  before upper={\popboxhead{Prerequisites}{popblue}{}}}
\newtcolorbox{bilan}{popbase,colframe=popink,colback=white,boxrule=2.8pt,
  before upper={\popboxhead{Fiche bilan}{poporange}{L'essentiel à connaître pour l'examen}}}
% Figure encadrée avec légende
\newtcolorbox{popfigure}[1][]{enhanced,breakable=false,colback=white,colframe=popline,boxrule=1.4pt,arc=8pt,
  left=10pt,right=10pt,top=10pt,bottom=8pt,before skip=10pt,after skip=12pt,halign=center,
  lower separated=false,halign lower=center,
  fontlower=\popsemi\fontsize{9}{11.5}\selectfont\color{popmuted},
  #1}
\newcommand{\legende}[1]{\par\vspace{6pt}{\popsemi\fontsize{9}{11.5}\selectfont\color{popmuted}#1}}

% ============================================================
% Signature OnBuch+
% ============================================================
\newcommand{\poplogoinline}[1]{{\poplogo\fontsize{#1}{#1}\selectfont\color{popink}OnBuch\color{poporange}+}}
\newcommand{\signatureonbuch}{%
  \begin{tcolorbox}[enhanced,colback=popink,colframe=popink,boxrule=2.2pt,arc=10pt,
      drop shadow=poporange,left=14pt,right=14pt,top=10pt,bottom=10pt,before skip=0pt,after skip=0pt]
    \tikz[baseline=-0.7ex]\node[circle,fill=popgreen,draw=popcream,line width=1.3pt,minimum size=22pt,inner sep=0]{\popcheck{white}};%
    \hspace{9pt}\begin{minipage}[c]{0.62\linewidth}
      {\popxbold\fontsize{12}{13}\selectfont\color{popcream}Signed\ \textbullet\ {\poplogo OnBuch\color{poporange}+}}\par\vspace{2pt}
      {\raggedright\popsemi\fontsize{8}{10.5}\selectfont\color{popline}\MakeUppercase{OnBuch+ teaching team}\\\MakeUppercase{Follows the official MINESEC syllabus}\par}
    \end{minipage}\hfill
    {\poplogo\fontsize{22}{22}\selectfont\color{popcream}OnBuch\color{poporange}+}
  \end{tcolorbox}%
}

% ============================================================
% Page de garde
% #1 titre · #2 matière · #3 niveau · #4 module · #5 n° leçon · #6 durée ·
% #7 description / objectifs (texte ou itemize)
% ============================================================
\newcommand{\pagegarde}[7]{%
  \thispagestyle{empty}
  \newgeometry{a4paper,margin=1.7cm,top=1.6cm,bottom=1.4cm}%
  \begin{tikzpicture}[remember picture,overlay]
    \fill[poporange!18] ([xshift=-3.2cm,yshift=-3.6cm]current page.north east) circle (4.6cm);
    \fill[poppurple!14] ([xshift=2.4cm,yshift=7.6cm]current page.south west) circle (3.8cm);
  \end{tikzpicture}%
  {\poplogo\fontsize{30}{30}\selectfont\color{popink}OnBuch\color{poporange}+}\hfill
  \raisebox{0.6ex}{\poppill{Signed course \textbullet\ Premium}{popink}{popcream}}\par
  \vspace{1pt}\popsquiggle{poppurple}\par
  \vspace{8pt}
  \begin{tcolorbox}[enhanced,colback=poporange,colframe=popink,boxrule=2.8pt,arc=13pt,
      drop shadow=popink,left=18pt,right=18pt,top=12pt,bottom=13pt,after skip=10pt]
    \poppill{#2 \textbullet\ #3}{popink}{popcream}\hspace{5pt}\poppill{#4}{white}{popink}\par\vspace{8pt}
    {\popdisplay\fontsize{14}{14}\selectfont\color{popink}CHAPTER #5}\par\vspace{4pt}
    {\raggedright\hyphenpenalty=10000\exhyphenpenalty=10000\popdisplay\fontsize{25}{27}\selectfont\color{popcream}\MakeUppercase{#1}\par}
    \vspace{8pt}
    {\popxbold\fontsize{9.5}{9.5}\selectfont\color{popink}\MakeUppercase{Suggested time: #6}}
  \end{tcolorbox}
  \begin{tcolorbox}[enhanced,colback=white,colframe=popink,boxrule=2.2pt,arc=10pt,
      drop shadow=popink,left=16pt,right=16pt,top=10pt,bottom=10pt,after skip=8pt]
    \poppill{In this chapter}{poppurple}{white}\par\vspace{5pt}
    \popsemi\fontsize{9.3}{12.6}\selectfont\color{popink}#7
  \end{tcolorbox}
  \noindent\begin{minipage}[t]{0.487\linewidth}
    \begin{tcolorbox}[enhanced,colback=popgreen,colframe=popink,boxrule=2.2pt,arc=10pt,
        drop shadow=popink,left=11pt,right=11pt,top=10pt,bottom=10pt,height=3.1cm,valign=center]
      \tcbox[colback=white,colframe=popink,boxrule=1.3pt,arc=5pt,boxsep=0pt,nobeforeafter,
        left=3pt,right=3pt,top=3pt,bottom=3pt,valign=center]{\qrcode[height=1.9cm]{https://play.google.com/store/apps/details?id=com.luvvix.onbuch&hl=en}}%
      \hspace{8pt}\begin{minipage}[c]{0.5\linewidth}
        {\popdisplay\fontsize{11}{12.5}\selectfont\color{white}\MakeUppercase{Download OnBuch}}\par\vspace{4pt}
        {\raggedright\popsemi\fontsize{8.4}{11}\selectfont\color{white}Free \textbullet\ Google Play\\AI tutor, past papers, results\par}
      \end{minipage}
    \end{tcolorbox}
  \end{minipage}\hfill
  \begin{minipage}[t]{0.487\linewidth}
    \begin{tcolorbox}[enhanced,colback=poppurple,colframe=popink,boxrule=2.2pt,arc=10pt,
        drop shadow=popink,left=11pt,right=11pt,top=10pt,bottom=10pt,height=3.1cm,valign=center]
      \tikz[baseline=-0.7ex]\node[circle,fill=white,draw=popink,line width=1.3pt,minimum size=1.6cm,inner sep=0]{\popdisplay\fontsize{24}{24}\selectfont\color{poppurple}+};%
      \hspace{8pt}\begin{minipage}[c]{0.6\linewidth}
        {\popdisplay\fontsize{11}{12.5}\selectfont\color{white}\MakeUppercase{Go OnBuch+}}\par\vspace{4pt}
        {\raggedright\popsemi\fontsize{8.4}{11}\selectfont\color{white}All signed courses, unlimited AI tutor, PDF sheets — OnBuch+ tab in the app\par}
      \end{minipage}
    \end{tcolorbox}
  \end{minipage}
  \vspace{8pt}
  \popcontenu
  \vfill
  \signatureonbuch
  \restoregeometry
  \newpage
}

% Rangée « Ce cours contient » (page de garde)
\newcommand{\poptile}[3]{% #1 couleur #2 gros texte #3 légende
  \begin{tcolorbox}[enhanced,colback=#1,colframe=popink,boxrule=2pt,arc=9pt,shadow={2.5pt}{-2.5pt}{0pt}{popink},
      left=6pt,right=6pt,top=9pt,bottom=9pt,halign=center,valign=center,height=1.75cm,width=\linewidth,nobeforeafter]
    {\popdisplay\fontsize{12.5}{13}\selectfont\color{white}\MakeUppercase{#2}}\par\vspace{3pt}
    {\centering\popsemi\fontsize{7.8}{9.5}\selectfont\color{white}#3\par}
  \end{tcolorbox}}
\newcommand{\popcontenu}{%
  \par\noindent{\popxboldls\fontsize{7.5}{7.5}\selectfont\color{popmuted}THIS SIGNED COURSE CONTAINS}\par\vspace{7pt}
  \noindent\begin{minipage}[t]{0.235\linewidth}\poptile{popblue}{Course}{complete, illustrated, step by step}\end{minipage}\hfill
  \begin{minipage}[t]{0.235\linewidth}\poptile{poporange}{Methods}{and worked examples}\end{minipage}\hfill
  \begin{minipage}[t]{0.235\linewidth}\poptile{poppink}{Exercises}{exam style + solutions}\end{minipage}\hfill
  \begin{minipage}[t]{0.235\linewidth}\poptile{popgreen}{Summary sheet}{the essentials to remember}\end{minipage}\par}

% Sommaire
\newtcolorbox{sommaire}{enhanced,colback=white,colframe=popink,boxrule=2.2pt,arc=10pt,
  drop shadow=popink,left=18pt,right=18pt,top=13pt,bottom=13pt,before skip=6pt,after skip=20pt,
  before upper={\poppill{Contents}{poppurple}{white}\par\vspace{9pt}\popsemi\fontsize{10.6}{14.5}\selectfont\color{popink}}}

% Signature de fin de cours — poussée tout en bas de la dernière page
\newcommand{\findecours}{%
  \par\vfill\nopagebreak
  \begin{center}\popsquiggle{poporange}\end{center}\vspace{4pt}
  \signatureonbuch
}
\AtBeginDocument{\def\llbracket{\mathopen{[\mkern-3.5mu[}}\def\rrbracket{\mathclose{]\mkern-3.5mu]}}} % intervalles d'entiers (glyphe ⟦ absent de la police)
% Glyphes absents de Fira Math : redéfinis à partir de glyphes disponibles
\AtBeginDocument{%
  \def\setminus{\mathbin{\backslash}}\let\smallsetminus\setminus
  \def\vdots{\mathord{\vbox{\baselineskip3.2pt\lineskiplimit0pt\kern4pt\hbox{.}\hbox{.}\hbox{.}}}}%
  \def\ddots{\mathinner{\mkern1mu\raise6.5pt\hbox{.}\mkern2mu\raise3.6pt\hbox{.}\mkern2mu\raise0.7pt\hbox{.}\mkern1mu}}%
  \def\triangle{\mathord{\scalebox{0.9}{$\symup{\Delta}$}}}%
  \def\nabla{\mathord{\raisebox{\depth}{\scalebox{1}[-1]{$\symup{\Delta}$}}}}%
  \def\checkmark{\popcheck{popdark}}%
  \def\star{\mathbin{\ast}}\def\bigstar{\mathord{\ast}}%
  \def\diamond{\mathbin{\scalebox{0.6}{$\Diamond$}}}%
  \def\bigcup{\mathop{\mathchoice{\raisebox{-0.3em}{\scalebox{1.7}{$\cup$}}}{\scalebox{1.25}{$\cup$}}{\cup}{\cup}}\displaylimits}%
  \def\bigcap{\mathop{\mathchoice{\raisebox{-0.3em}{\scalebox{1.7}{$\cap$}}}{\scalebox{1.25}{$\cap$}}{\cap}{\cap}}\displaylimits}%
  \def\lesseqqgtr{\mathrel{\lessgtr}}\def\bigoplus{\mathop{\oplus}\displaylimits}%
}
\providecommand{\ssi}{if and only if} % abréviation fréquente
\AtBeginDocument{\providecommand{\wideparen}{\overparen}} % arc de cercle

%% FIGFIX-BEGIN v4 — correctifs globaux des figures (docs/onbuch-generation/tools/figfix)
\usepackage{adjustbox}\usepackage{etoolbox}
% toute figure TikZ/pgfplots plus large que la ligne est réduite à la largeur disponible
\BeforeBeginEnvironment{tikzpicture}{\begin{adjustbox}{max width=\linewidth}}
\AfterEndEnvironment{tikzpicture}{\end{adjustbox}}
% légendes pgfplots lisibles par-dessus les courbes
\pgfplotsset{every axis/.append style={legend style={fill=white,fill opacity=0.92,text opacity=1,draw=black!15,inner sep=3pt}}}
% étiquettes d'arêtes (syntaxe edge["..."]) avec fond blanc
\tikzset{every edge quotes/.append style={fill=white,inner sep=1.2pt}}
%% FIGFIX-END
\begin{document}

\pagegarde{Alternating Current}{Physics}{Upper Sixth}{Upper Sixth — Physics}{9}{3 periods}{Ce chapitre explore les grandeurs efficaces du courant alternatif, l'impédance et la résonance des circuits RLC, la théorie des transformateurs idéaux et réels, ainsi que les techniques de redressage et de lissage pour obtenir une tension continue. Il prépare l'élève à analyser, dimensionner et dépanner les systèmes d'alimentation électrique courants.
\vspace{4pt}
\begin{multicols}{2}\small
\begin{itemize}
\item Calculer la valeur efficace (RMS) d'une tension ou d'un courant sinusoïdal et non sinusoïdal.
\item Déterminer l'impédance complexe d'un circuit RLC série ou parallèle et tracer son diagramme de phase.
\item Appliquer la condition de résonance f₀ = 1/(2π√LC) et évaluer le facteur de qualité Q et la bande passante.
\item Calculer les puissances active, réactive et apparente, le facteur de puissance et proposer une correction.
\item Énoncer et utiliser les relations du transformateur idéal Vp/Vs = Np/Ns = Is/Ip.
\item Identifier les pertes dans un transformateur réel (fer, cuivre, dispersées) et expliquer leur minimisation.
\end{itemize}\end{multicols}}

\begin{sommaire}
\begin{multicols}{2}
\begin{enumerate}
\item Grandeurs efficaces et impédance en régime sinusoïdal
\item Résonance et puissance dans les circuits AC
\item Transformateurs : modèle idéal et pertes réelles
\item Redressement et lissage des signaux alternatifs
\item Integration activity
\item Methods and worked examples
\item Exercises
\item Solutions to the exercises
\item Summary sheet
\end{enumerate}
\end{multicols}
\end{sommaire}

\begin{objectifs}
\begin{itemize}
\item Calculer la valeur efficace (RMS) d'une tension ou d'un courant sinusoïdal et non sinusoïdal.
\item Déterminer l'impédance complexe d'un circuit RLC série ou parallèle et tracer son diagramme de phase.
\item Appliquer la condition de résonance f₀ = 1/(2π√LC) et évaluer le facteur de qualité Q et la bande passante.
\item Calculer les puissances active, réactive et apparente, le facteur de puissance et proposer une correction.
\item Énoncer et utiliser les relations du transformateur idéal Vp/Vs = Np/Ns = Is/Ip.
\item Identifier les pertes dans un transformateur réel (fer, cuivre, dispersées) et expliquer leur minimisation.
\item Analyser les redresseurs à demi‑onde, à double alternance (centre‑tapé et pont) et calculer leur facteur de ripple.
\item Dimensionner un condensateur de lissage pour obtenir une tension de sortie avec un ripple spécifié.
\item Évaluer le rendement et la régulation de tension d'un transformateur sous charge.
\item Concevoir une alimentation complète (transformateur + redresseur + filtre) pour une charge donnée.
\end{itemize}
\end{objectifs}

\begin{prerequis}
\begin{itemize}
\item Fonctions sinusoïdales : amplitude, période, fréquence, phase (Terminale).
\item Loi d'Ohm et lois de Kirchhoff en régime continu (Première).
\item Notions de base d'électromagnétisme : induction mutuelle, flux magnétique (Première).
\item Manipulation des nombres complexes : forme algébrique, exponentielle, diagramme d'Argand (Terminale).
\item Énergie et puissance en régime continu : P = UI, effet Joule (Première).
\end{itemize}
\end{prerequis}

\newpage

\coursec{Grandeurs efficaces et impédance en régime sinusoïdal}

Dans les quartiers de Yaoundé, les poteaux électriques ploient sous le poids des transformateurs qui abaissent la tension du réseau pour alimenter nos foyers. Chaque fois que vous branchez votre téléphone, un petit transformateur, un redresseur et un condensateur travaillent de concert pour convertir la tension alternative de \SI{230}{V} en une tension continue stable de \SI{5}{V}. Pourquoi le courant alternatif est-il privilégié pour le transport, et comment les composants électroniques le transforment-ils en tension stable pour nos appareils ? Avant de répondre, nous devons apprendre le langage du courant alternatif : comment mesurer une grandeur qui change de signe cinquante fois par seconde, et comment résistances, bobines et condensateurs y répondent chacun. C'est l'objet de cette première section.

\courssub{Valeur instantanée, amplitude et valeur efficace}

\textbf{Observation.}\par
Reliez un générateur basse tension alternatif à un oscilloscope. Le spot dessine une onde qui monte, descend, s'inverse et se répète. Le secteur au Cameroun le fait à une fréquence $f = \SI{50}{Hz}$ : une période complète dure
\[
T=\frac{1}{f}=\frac{1}{50}=\SI{0.020}{s}=\SI{20}{ms}.
\]

\begin{definition}[Valeur instantanée, amplitude, période]
Une tension alternative sinusoïdale s'écrit
\[
u(t)=V_{0}\cos(\omega t+\varphi),
\]
où $u(t)$ est la \cle{valeur instantanée} (volts), $V_{0}$ est l'\cle{amplitude} ou \emph{valeur de crête} (volts), $\omega = 2\pi f$ est la pulsation (rad\,s$^{-1}$) et $\varphi$ est la phase à l'origine. De même pour le courant : $i(t)=I_{0}\cos(\omega t+\varphi)$.
\end{definition}

L'amplitude $V_{0}$ donne l'\emph{extrême} de la tension, mais la tension est la plupart du temps plus petite que $V_{0}$, et la moitié du temps elle est négative. Si l'on veut un seul nombre qui dise ce que le courant alternatif peut \emph{réellement faire} — par exemple, combien de chaleur il produit dans une bouilloire — il faut une meilleure mesure.

\textbf{L'idée thermique derrière la valeur efficace.}\par
La puissance instantanée dissipée dans une résistance $R$ est $p(t)=R\,i^{2}(t)$. Comme $i^{2}$ est toujours positif, le courant alternatif chauffe la résistance pendant les deux demi-périodes. La puissance moyenne sur une période est
\[
P_{\text{moy}}=\frac{1}{T}\int_{0}^{T}R\,i^{2}(t)\,dt = R\,\langle i^{2}\rangle .
\]
On définit la \cle{valeur efficace} (RMS : \emph{root mean square}) comme le courant continu qui dissiperait la \emph{même} puissance moyenne dans la même résistance : $P_{\text{moy}} = R\,I_{\mathrm{eff}}^{2}$.

\begin{definition}[Valeur efficace (RMS)]
La valeur efficace d'un courant périodique est la racine carrée de la moyenne de son carré sur une période :
\[
I_{\mathrm{eff}}=\sqrt{\dfrac{1}{T}\displaystyle\int_{0}^{T} i^{2}(t)\,dt},
\qquad
V_{\mathrm{eff}}=\sqrt{\dfrac{1}{T}\displaystyle\int_{0}^{T} u^{2}(t)\,dt}.
\]
C'est la valeur continue qui délivre la même puissance moyenne à une résistance.
\end{definition}

\begin{propriete}[Valeur efficace d'un sinusoïde pur]
Pour $i(t)=I_{0}\cos(\omega t)$ et $u(t)=V_{0}\cos(\omega t)$,
\[
I_{\mathrm{eff}}=\frac{I_{0}}{\sqrt{2}},\qquad V_{\mathrm{eff}}=\frac{V_{0}}{\sqrt{2}}.
\]
\emph{Preuve (examinable).} En utilisant $\cos^{2}(\omega t)=\tfrac{1}{2}\bigl[1+\cos(2\omega t)\bigr]$,
\[
I_{\mathrm{eff}}^{2}=\frac{1}{T}\int_{0}^{T}I_{0}^{2}\cos^{2}(\omega t)\,dt
=\frac{I_{0}^{2}}{2T}\int_{0}^{T}\bigl[1+\cos(2\omega t)\bigr]dt
=\frac{I_{0}^{2}}{2T}\left[T+0\right]=\frac{I_{0}^{2}}{2},
\]
car l'intégrale de $\cos(2\omega t)$ sur une période entière est nulle. D'où $I_{\mathrm{eff}}=I_{0}/\sqrt{2}$. Même calcul pour la tension.
\end{propriete}

\begin{attention}[Crête versus efficace]
La tension du secteur ``\SI{230}{V}'' est une valeur \emph{efficace}. Sa crête est $V_{0}=230\sqrt{2}\approx \SI{325}{V}$. Ne mélangez jamais valeurs de crête et efficaces dans une même formule : les formules de puissance $P=VI=RI^{2}$ ne valent qu'avec des valeurs \textbf{efficaces}. Utiliser $V_{0}$ dans $P=V^{2}/R$ double la puissance réelle.
\end{attention}

\textbf{Extension aux signaux non sinusoïdaux.}\par
La définition intégrale fonctionne pour \emph{toute} forme d'onde. Pour un signal carré prenant les valeurs $+V_{0}$ et $-V_{0}$, $u^{2}=V_{0}^{2}$ en tout instant, donc $V_{\mathrm{eff}}=V_{0}$. Pour une sinusoïde redressée à demi-onde (zéro la moitié du temps, $V_{0}\cos\omega t$ l'autre moitié), la moyenne de $u^{2}$ est divisée par deux, donnant $V_{\mathrm{eff}}=V_{0}/2$. Ces résultats seront utiles pour l'étude du redressement en section 4.

\begin{popfigure}
\begin{center}
\begin{tikzpicture}[scale=1.0]
% axes
\draw[->,thick,popink] (0,0) -- (7.2,0) node[below left]{$t$};
\draw[->,thick,popink] (0,-1.8) -- (0,2.0) node[left]{$u(t)$};
% sine wave: period 6, amplitude 1.5
\draw[very thick,popblue,domain=0:6.6,samples=120] plot (\x,{1.5*cos(60*\x)});
% peak line
\draw[dashed,poporange] (0,1.5) -- (1.5,1.5);
\draw[<->,poporange,thick] (1.5,0) -- (1.5,1.5) node[midway,right]{$V_{0}$};
% rms line
\draw[dashed,popgreen] (0,1.06) -- (4.6,1.06);
\node[popgreen,anchor=west] at (4.6,1.06) {$V_{\mathrm{eff}}=V_{0}/\sqrt{2}$};
% negative peak
\draw[dashed,popmuted] (0,-1.5) -- (4.5,-1.5);
\node[popmuted,anchor=north west] at (4.6,-1.5) {$-V_{0}$};
% period arrow
\draw[<->,thick,popdark] (0,-1.75) -- (6,-1.75) node[midway,below]{$T$};
\draw[dashed,popmuted] (6,-1.5) -- (6,0);
\end{tikzpicture}
\end{center}
\legende{Tension sinusoïdale montrant la valeur de crête $V_{0}$, la valeur efficace $V_{\mathrm{eff}}=V_{0}/\sqrt{2}$ et la période $T$.}
\end{popfigure}

\begin{exemplebox}[Lecture d'un oscilloscope]
Un oscilloscope affiche un signal secteur avec une hauteur crête-à-crête de \SI{650}{V}. La valeur de crête est $V_{0}=650/2=\SI{325}{V}$, donc $V_{\mathrm{eff}}=325/\sqrt{2}\approx \SI{230}{V}$ — exactement la tension nominale du secteur camerounais.
\end{exemplebox}

\courssub{Fasors et représentation complexe}

Additionner deux sinusoïdes de même fréquence par trigonométrie est fastidieux. L'astuce consiste à représenter chaque grandeur par un \cle{fasor} (ou \emph{phasor}) : un vecteur de longueur égale à la valeur \emph{efficace} (convention usuelle en puissance), tournant à la vitesse angulaire $\omega$, dont la projection sur l'axe vertical donne la valeur instantanée. Comme tous les fasors d'un circuit tournent ensemble, on les ``gèle'' à $t=0$ et on les additionne comme des vecteurs ordinaires.

Mathématiquement, on remplace $u(t)=V_{0}\cos(\omega t+\varphi)$ par le nombre complexe
\[
\underline{V}=V_{\mathrm{eff}}\,e^{j\varphi}=V_{\mathrm{eff}}\bigl(\cos\varphi+j\sin\varphi\bigr),\qquad j^{2}=-1,
\]
et on retrouve le signal physique comme la partie réelle de $\sqrt{2}\,\underline{V}\,e^{j\omega t}$. La dérivation devient une multiplication par $j\omega$ : c'est là toute la puissance de la méthode.

\courssub{Impédance de R, L et C}

Appliquons $u(t)=V_{0}\cos\omega t$ à chaque composant et déterminons le courant.

\textbf{Résistance.} $u=Ri$ donne $i=\dfrac{V_{0}}{R}\cos\omega t$ : tension et courant sont \emph{en phase} ($\varphi=0$). Donc $\underline{Z}_{R}=R$.

\textbf{Bobine (inductance).} $u=L\dfrac{di}{dt}$. En complexe, $\underline{V}=j\omega L\,\underline{I}$, d'où
\[
\underline{Z}_{L}=j\omega L = jX_{L},\qquad X_{L}=\omega L \ \text{(réactance inductive, ohms)}.
\]
La tension \emph{devance} le courant de $+90^{\circ}$.

\textbf{Condensateur.} $i=C\dfrac{du}{dt}$. En complexe, $\underline{I}=j\omega C\,\underline{V}$, soit $\underline{V}=\dfrac{1}{j\omega C}\,\underline{I}=-j\dfrac{1}{\omega C}\,\underline{I}$, donc
\[
\underline{Z}_{C}=\frac{1}{j\omega C}=-\frac{j}{\omega C}=-jX_{C},\qquad X_{C}=\frac{1}{\omega C}\ \text{(réactance capacitive, ohms)}.
\]
Le courant \emph{devance} la tension de $90^{\circ}$ (la tension retarde de $-90^{\circ}$).

\begin{aretenir}[Relations de phase]
\begin{itemize}
\item Résistance : $V$ et $I$ en phase, $\varphi=0^{\circ}$.
\item Bobine : $V$ devance $I$ de $+90^{\circ}$ (mnémotechnique ``ELI'' : dans une inductance $L$, $E$ (tension) devance $I$).
\item Condensateur : $I$ devance $V$ de $90^{\circ}$ (mnémotechnique ``ICE'' : dans un condensateur $C$, $I$ devance $E$).
\end{itemize}
\end{aretenir}

\courssub{Circuit RLC série : impédance totale}

Pour des composants en série, le même courant les traverse et les tensions s'additionnent \emph{comme des fasors} :
\[
\underline{V}=\underline{V}_{R}+\underline{V}_{L}+\underline{V}_{C}
=\bigl[R+j\omega L-\tfrac{j}{\omega C}\bigr]\underline{I}.
\]

\begin{definition}[Impédance complexe]
L'impédance complexe d'un circuit RLC série est
\[
\underline{Z}=R+j\left(\omega L-\frac{1}{\omega C}\right)=R+j(X_{L}-X_{C}),
\]
et la loi d'Ohm généralisée s'écrit $\underline{V}=\underline{Z}\,\underline{I}$.
\end{definition}

\begin{methode}[Module et phase de Z]
\begin{enumerate}
\item Calculez $X_{L}=\omega L$ et $X_{C}=1/(\omega C)$.
\item Formez $\underline{Z}=R+j(X_{L}-X_{C})$.
\item Module : $|\underline{Z}|=\sqrt{R^{2}+(X_{L}-X_{C})^{2}}$.
\item Phase de la tension par rapport au courant : $\tan\varphi=\dfrac{X_{L}-X_{C}}{R}$.
\item Courant efficace : $I_{\mathrm{eff}}=V_{\mathrm{eff}}/|\underline{Z}|$.
\end{enumerate}
\end{methode}

\begin{attention}[Astuce du triangle d'impédance]
Tracez le \cle{triangle d'impédance} : côté horizontal $R$, côté vertical $(X_{L}-X_{C})$, hypoténuse $|\underline{Z}|$. L'angle à la base est exactement le déphasage $\varphi=\arctan\!\bigl((X_{L}-X_{C})/R\bigr)$. Si $X_{L}>X_{C}$ le circuit est inductif ($\varphi>0$) ; si $X_{C}>X_{L}$ il est capacitif ($\varphi<0$).
\end{attention}

\begin{popfigure}
\begin{center}
\begin{tikzpicture}[scale=1.1]
% axes
\draw[->,thick,popink] (0,0) -- (5.6,0) node[below left]{$I$ (référence)};
\draw[->,thick,popink] (0,-2.2) -- (0,3.0);
% V_R
\draw[->,very thick,popgreen] (0,0) -- (3,0) node[midway,below]{$V_{R}=RI$};
% V_L
\draw[->,very thick,popblue] (3,0) -- (3,2.4) node[midway,right]{$V_{L}=X_{L}I$};
% V_C (opposé à V_L)
\draw[->,very thick,poporange] (3,2.4) -- (3,1.0) node[midway,right]{$V_{C}=X_{C}I$};
% V total
\draw[->,very thick,popdark] (0,0) -- (3,1.0) node[midway,above left]{$V$};
% angle phi
\draw[popdark] (1.2,0) arc (0:18.4:1.2);
\node[popdark] at (1.55,0.22) {$\varphi$};
% dashed helper
\draw[dashed,popmuted] (3,1.0) -- (0,1.0);
\node[popmuted,left] at (0,1.0) {$(X_{L}-X_{C})I$};
\end{tikzpicture}
\end{center}
\legende{Diagramme de fasors pour un circuit RLC série avec $X_{L}>X_{C}$ : la tension totale $V$ devance le courant de l'angle $\varphi$.}
\end{popfigure}

\textbf{Association parallèle.}\par
En parallèle la tension est commune et les courants s'additionnent : il est plus commode d'utiliser l'\cle{admittance} $\underline{Y}=1/\underline{Z}$ (siemens, S). On a alors $\underline{Y}_{\text{total}}=\underline{Y}_{1}+\underline{Y}_{2}+\cdots$ Pour une résistance, une bobine et un condensateur en parallèle,
\[
\underline{Y}=\frac{1}{R}+j\left(\omega C-\frac{1}{\omega L}\right).
\]

\begin{exoresolu}[Circuit RLC série sur le secteur]
Une bobine de résistance $R=\SI{40}{\ohm}$ et d'inductance $L=\SI{0.20}{H}$ est montée en série avec un condensateur $C=\SI{10}{\micro F}$ aux bornes du secteur ($V_{\mathrm{eff}}=\SI{230}{V}$, $f=\SI{50}{Hz}$). Déterminer (a) l'impédance, (b) le courant, (c) le déphasage, (d) les tensions aux bornes de chaque élément.
\tcblower
\textbf{Solution.} Pulsation : $\omega=2\pi f=2\pi\times 50=\SI{314.2}{rad.s^{-1}}$.
\begin{align*}
X_{L}&=\omega L=314.2\times 0.20=\SI{62.8}{\ohm},\\
X_{C}&=\frac{1}{\omega C}=\frac{1}{314.2\times 10\times 10^{-6}}=\SI{318.3}{\ohm}.
\end{align*}
(a) $\underline{Z}=40+j(62.8-318.3)=40-j255.5\ \Omega$, donc
\[
|\underline{Z}|=\sqrt{40^{2}+255.5^{2}}=\sqrt{1600+65280}\approx\SI{258.4}{\ohm}.
\]
(b) $I_{\mathrm{eff}}=\dfrac{V_{\mathrm{eff}}}{|\underline{Z}|}=\dfrac{230}{258.4}\approx\SI{0.890}{A}$.
(c) $\tan\varphi=\dfrac{62.8-318.3}{40}=-6.39\Rightarrow\varphi\approx -81.1^{\circ}$ : le circuit est fortement capacitif, le courant devance la tension.
(d) $V_{R}=RI=40\times 0.890\approx\SI{35.6}{V}$ ; $V_{L}=X_{L}I\approx\SI{55.9}{V}$ ; $V_{C}=X_{C}I\approx\SI{283.3}{V}$.
Remarquez que $V_{C}$ dépasse la tension d'alimentation ! C'est normal en courant alternatif : les fasors $V_{L}$ et $V_{C}$ s'opposent partiellement, et seule la somme vectorielle $V=\sqrt{V_{R}^{2}+(V_{L}-V_{C})^{2}}=\sqrt{35.6^{2}+227.4^{2}}\approx\SI{230}{V}$ égale la tension du secteur. N'additionnez jamais les tensions individuelles arithmétiquement.
\end{exoresolu}

\begin{exercice}[Entraînement]
Une résistance $R=\SI{100}{\ohm}$, une bobine $L=\SI{0.50}{H}$ et un condensateur $C=\SI{20}{\micro F}$ sont en série aux bornes d'une alimentation \SI{240}{V}, \SI{50}{Hz}. Calculer $X_{L}$, $X_{C}$, $|\underline{Z}|$, $I_{\mathrm{eff}}$ et $\varphi$. Préciser si le circuit est inductif ou capacitif.
\end{exercice}

\begin{corrige}[Exercice 1]
$\omega=\SI{314.2}{rad.s^{-1}}$ ; $X_{L}=314.2\times 0.50=\SI{157.1}{\ohm}$ ; $X_{C}=1/(314.2\times 20\times10^{-6})=\SI{159.2}{\ohm}$. Alors $|\underline{Z}|=\sqrt{100^{2}+(157.1-159.2)^{2}}=\sqrt{10000+4.4}\approx\SI{100.0}{\ohm}$ ; $I_{\mathrm{eff}}=240/100.0=\SI{2.40}{A}$ ; $\tan\varphi=-2.1/100\Rightarrow\varphi\approx-1.2^{\circ}$. Le circuit est très légèrement capacitif — il est presque en résonance, situation que nous exploiterons en section 2.
\end{corrige}

\begin{savaistu}[Pourquoi le courant alternatif a gagné la bataille]
Dans les années 1880, Edison défendait le courant continu tandis que Westinghouse et Tesla championnaient le courant alternatif. L'alternatif a gagné précisément parce que sa tension peut être élevée ou abaissée sans effort par les transformateurs (section 3) : transporter à haute tension et faible courant réduit drastiquement les pertes Joule $RI^{2}$ dans les lignes. La valeur efficace que vous venez de rencontrer est ce qui permet une comparaison équitable entre les deux systèmes — \SI{230}{V} alternatif chauffe votre fer à repasser exactement comme \SI{230}{V} continu le ferait.
\end{savaistu}

\coursec{Resonance and power in a.c. circuits}

In the previous section we learnt how to describe a sinusoidal quantity by its \cle{rms value} and how each component (resistor, inductor, capacitor) opposes the current through its \cle{reactance}, leading to the notion of \cle{impedance} $Z$. Two questions now arise naturally. First, since $X_L = L\omega$ grows with frequency while $X_C = 1/(C\omega)$ decreases, there must exist a special frequency where the two effects cancel: what happens then? Second, because current and voltage are generally out of phase, what power is \emph{really} delivered to an a.c.\ circuit? These two questions — \cle{resonance} and \cle{power} — are the heart of this section, and they are among the most frequently examined topics at the GCE Advanced Level.

\courssub{Series resonance}

\textbf{Observation.}\par
Tune an old radio: as you turn the knob, one station suddenly becomes loud and clear while the others fade away. The radio's $RLC$ circuit is being adjusted so that its natural frequency matches the frequency of the transmitter. This selective amplification is \cle{resonance}.

\begin{popfigure}
\begin{center}
\begin{tikzpicture}
\begin{axis}[
  width=0.85\linewidth, height=6.2cm,
  xlabel={frequency $f$}, ylabel={current amplitude $I_{\mathrm{rms}}$},
  xmin=0, xmax=10, ymin=0, ymax=11,
  xtick={5}, xticklabels={$f_0$},
  ytick=\empty,
  axis lines=left, thick,
  every axis x label/.style={at={(ticklabel* cs:1.0)}, anchor=west},
]
\addplot[popblue, very thick, domain=0.6:10, samples=200, smooth] {10/(sqrt(max(0,1 + 9*(x/5 - 5/x)^2)))};
\addplot[poporange, thick, dashed, domain=0.6:10, samples=200, smooth] {10/(sqrt(max(0,1 + 36*(x/5 - 5/x)^2)))};
\draw[dashed, popmuted] (axis cs:5,0) -- (axis cs:5,10);
\draw[dashed, popmuted] (axis cs:0,7.07) -- (axis cs:5,7.07);
\node[popblue, anchor=south] at (axis cs:5,10) {$I_{\max}=\dfrac{V_{\mathrm{rms}}}{R}$};
\node[popmuted, anchor=west] at (axis cs:0.15,7.3) {$I_{\max}/\sqrt{2}$};
\draw[<->, popdark, thick] (axis cs:4.1,7.07) -- (axis cs:5.9,7.07);
\node[popdark] at (axis cs:5,6.2) {$\Delta f$};
\node[popblue, anchor=west] at (axis cs:7.2,8.6) {small $Q$};
\node[poporange, anchor=west] at (axis cs:7.2,4.2) {large $Q$};
\end{axis}
\end{tikzpicture}
\end{center}
\legende{Resonance curve of a series $RLC$ circuit: rms current versus frequency. The current peaks at $f_0$; the bandwidth $\Delta f$ is measured between the frequencies where the current falls to $I_{\max}/\sqrt{2}$. A large quality factor $Q$ gives a sharp, selective peak.}
\end{popfigure}

Consider a series $RLC$ circuit driven by a source of constant rms voltage $V_{\mathrm{rms}}$ but variable angular frequency $\omega$. The impedance is
\[
Z = \sqrt{R^2 + \left(L\omega - \frac{1}{C\omega}\right)^{\!2}} .
\]
The term in brackets can vanish. The \cle{resonance condition} is
\[
X_L = X_C \quad\Longleftrightarrow\quad L\omega_0 = \frac{1}{C\omega_0}
\quad\Longrightarrow\quad \boxed{\ \omega_0 = \frac{1}{\sqrt{LC}}\ }, \qquad
\boxed{\ f_0 = \frac{1}{2\pi\sqrt{LC}}\ } .
\]

\begin{propriete}[Series resonance (examinable)]
At series resonance, $X_L = X_C$, so the impedance is purely resistive and \emph{minimal}: $Z = R$. Consequently the current is \emph{maximal}, $I_{\max} = V_{\mathrm{rms}}/R$, and the current is \emph{in phase} with the supply voltage ($\varphi = 0$, $\cos\varphi = 1$). The voltages across the inductor and the capacitor are equal in magnitude and opposite in phase; each can individually \emph{exceed} the supply voltage (voltage magnification).
\end{propriete}

\begin{attention}[Danger at resonance]
Because $V_L = I_{\max}L\omega_0$ and $V_C = I_{\max}/(C\omega_0)$ can be much larger than the supply voltage, a poorly designed circuit can destroy its capacitor at resonance. Never assume a component voltage cannot exceed the source voltage in a.c.\ circuits!
\end{attention}

\textbf{Quality factor and bandwidth.}\par
The sharpness of the resonance is measured by the \cle{quality factor}
\[
Q = \frac{\omega_0 L}{R} = \frac{1}{\omega_0 C R} = \frac{1}{R}\sqrt{\frac{L}{C}} .
\]
$Q$ is dimensionless. It also equals the voltage magnification $V_L/V_{\mathrm{rms}}$ at resonance. The \cle{bandwidth} $\Delta f$ is the width of the resonance peak measured between the two frequencies where the current has fallen to $I_{\max}/\sqrt{2}$ (the "half-power" points), and one shows that
\[
\Delta f = \frac{f_0}{Q} .
\]

\begin{methode}[Computing $Q$ and $\Delta f$]
\begin{enumerate}
\item Compute the resonant frequency $f_0 = \dfrac{1}{2\pi\sqrt{LC}}$, then $\omega_0 = 2\pi f_0$.
\item Compute the quality factor $Q = \dfrac{\omega_0 L}{R}$ (check with $\dfrac{1}{\omega_0 C R}$).
\item Deduce the bandwidth $\Delta f = \dfrac{f_0}{Q}$.
\item Interpret: $Q \gg 1$ means a sharp, highly selective circuit (good for a radio tuner); $Q$ close to $1$ means a broad response.
\end{enumerate}
\end{methode}

\begin{exoresolu}[Series resonance of a radio tuner]
A series circuit consists of a coil of inductance $L = \SI{2.0}{mH}$ and resistance $R = \SI{10}{\ohm}$, and a capacitor $C = \SI{200}{pF}$, connected to a generator of constant rms voltage $V_{\mathrm{rms}} = \SI{5.0}{V}$. Determine (a) the resonant frequency, (b) the maximum current, (c) the quality factor, (d) the bandwidth, (e) the voltage across the capacitor at resonance.
\tcblower
\textbf{(a)} Resonant frequency:
\[
f_0 = \frac{1}{2\pi\sqrt{LC}} = \frac{1}{2\pi\sqrt{2.0\times10^{-3}\times 200\times10^{-12}}}
= \frac{1}{2\pi\sqrt{4.0\times10^{-13}}} \approx \SI{2.5e5}{Hz} = \SI{250}{kHz}.
\]
\textbf{(b)} At resonance $Z = R$, so
\[
I_{\max} = \frac{V_{\mathrm{rms}}}{R} = \frac{5.0}{10} = \SI{0.50}{A}.
\]
\textbf{(c)} With $\omega_0 = 2\pi f_0 \approx \SI{1.58e6}{rad.s^{-1}}$:
\[
Q = \frac{\omega_0 L}{R} = \frac{1.58\times10^{6}\times 2.0\times10^{-3}}{10} \approx 316.
\]
\textbf{(d)} Bandwidth:
\[
\Delta f = \frac{f_0}{Q} = \frac{2.5\times10^{5}}{316} \approx \SI{7.9e2}{Hz} \approx \SI{0.8}{kHz}.
\]
\textbf{(e)} Capacitor voltage at resonance:
\[
V_C = \frac{I_{\max}}{C\omega_0} = Q\,V_{\mathrm{rms}} = 316 \times 5.0 \approx \SI{1.6e3}{V}.
\]
The capacitor must withstand about $\SI{1.6}{kV}$ although the source supplies only $\SI{5}{V}$: a spectacular illustration of voltage magnification.
\end{exoresolu}

\courssub{Parallel resonance}

\begin{propriete}[Parallel resonance]
For a parallel $LC$ (or $RLC$) circuit, resonance occurs when the susceptances cancel. The \cle{admittance} $Y = 1/Z$ is then \emph{minimal}, so the impedance is \emph{maximal} and the total current drawn from the source is minimal. A parallel resonant circuit therefore \emph{blocks} the frequency $f_0$ — the opposite behaviour of the series circuit, which \emph{selects} it.
\end{propriete}

For parallel circuits it is convenient to work with \cle{admittance} $Y$ (unit: siemens, \si{S}), whose real part is the \cle{conductance} $G = 1/R$ and whose imaginary part is the \cle{susceptance} $B = \omega C - 1/(\omega L)$. Admittances in parallel simply add, exactly as impedances add in series.

\courssub{Power in a.c. circuits}

\textbf{Observation.}\par
An electric iron (almost a pure resistor) heats strongly, while an ideal inductor connected to the same mains becomes only slightly warm in its winding. Yet an ammeter shows a current in both cases. Where does the energy go? The answer lies in the phase difference $\varphi$.

\textbf{Instantaneous and average power.}\par
Let $v(t) = V_{\max}\cos(\omega t)$ and $i(t) = I_{\max}\cos(\omega t - \varphi)$. The \cle{instantaneous power} is
\[
p(t) = v(t)\,i(t) = V_{\max}I_{\max}\cos(\omega t)\cos(\omega t - \varphi).
\]
Using $2\cos A\cos B = \cos(A-B) + \cos(A+B)$:
\[
p(t) = \frac{V_{\max}I_{\max}}{2}\big[\cos\varphi + \cos(2\omega t - \varphi)\big].
\]
The second term oscillates at twice the supply frequency and averages to zero over one period. The average (\cle{active}) power is therefore
\[
\boxed{\ P = \frac{V_{\max}I_{\max}}{2}\cos\varphi = V_{\mathrm{rms}}\,I_{\mathrm{rms}}\cos\varphi\ } \qquad (\text{watt, W}).
\]

\begin{definition}[The three powers and the power factor]
For a load with rms voltage $V_{\mathrm{rms}}$, rms current $I_{\mathrm{rms}}$ and phase angle $\varphi$:
\begin{itemize}
\item \cle{Active power} $P = V_{\mathrm{rms}}I_{\mathrm{rms}}\cos\varphi$, in watts (\si{W}): the power actually converted into heat, light or mechanical work.
\item \cle{Reactive power} $Q = V_{\mathrm{rms}}I_{\mathrm{rms}}\sin\varphi$, in volt-ampere reactive (\si{var}): power exchanged back and forth with magnetic and electric fields; it does no net work.
\item \cle{Apparent power} $S = V_{\mathrm{rms}}I_{\mathrm{rms}}$, in volt-ampere (\si{VA}): the power the source must be able to supply.
\item \cle{Power factor} $\cos\varphi = \dfrac{P}{S}$.
\end{itemize}
\end{definition}

\begin{popfigure}
\begin{center}
\begin{tikzpicture}[scale=1.15, thick]
\coordinate (O) at (0,0);
\coordinate (A) at (4.6,0);
\coordinate (B) at (4.6,2.6);
\draw[->, popblue, very thick] (O) -- (A) node[midway, below] {$P = V_{\mathrm{rms}}I_{\mathrm{rms}}\cos\varphi$ (W)};
\draw[->, poporange, very thick] (A) -- (B) node[midway, right] {$Q = V_{\mathrm{rms}}I_{\mathrm{rms}}\sin\varphi$ (var)};
\draw[->, poppurple, very thick] (O) -- (B) node[midway, above left] {$S = V_{\mathrm{rms}}I_{\mathrm{rms}}$ (VA)};
\draw (1.1,0) arc (0:29.5:1.1);
\node at (1.55,0.32) {$\varphi$};
\draw[dashed, popmuted] (4.6,0) -- (4.6,-0.001);
\draw (4.6,0.28) -- (4.32,0.28) -- (4.32,0);
\end{tikzpicture}
\end{center}
\legende{The power triangle: $S^2 = P^2 + Q^2$ and $\cos\varphi = P/S$. An inductive load gives $Q > 0$; a capacitive load gives $Q < 0$ (drawn downward).}
\end{popfigure}

\begin{aretenir}[Power triangle]
The three powers form a right-angled triangle:
\[
S^2 = P^2 + Q^2, \qquad \tan\varphi = \frac{Q}{P}, \qquad \cos\varphi = \frac{P}{S}.
\]
A pure resistor has $\varphi = 0$: $P = S$, $Q = 0$. A pure inductor or capacitor has $|\varphi| = 90^{\circ}$: $P = 0$ — it consumes no average power at all.
\end{aretenir}

\begin{attention}[Reactive power is not free]
Although $Q$ performs no useful work, the current associated with it is real and flows through the transmission lines, where it causes Joule heating ($P_{\text{line}} = RI_{\mathrm{rms}}^2$) and extra voltage drops. This is why the electricity utility penalises industrial consumers with a poor (low) power factor, and why improving $\cos\varphi$ reduces line losses.
\end{attention}

\begin{exoresolu}[Power in a workshop motor]
A single-phase motor connected to the mains ($V_{\mathrm{rms}} = \SI{220}{V}$, $f = \SI{50}{Hz}$) draws $I_{\mathrm{rms}} = \SI{8.0}{A}$ with a power factor $\cos\varphi = 0.75$ (inductive). Find (a) the apparent, active and reactive powers; (b) the energy consumed in 6 hours of operation.
\tcblower
\textbf{(a)} Apparent power:
\[
S = V_{\mathrm{rms}}I_{\mathrm{rms}} = 220 \times 8.0 = \SI{1.76e3}{VA} = \SI{1.76}{kVA}.
\]
Active power:
\[
P = S\cos\varphi = 1760 \times 0.75 = \SI{1.32e3}{W} = \SI{1.32}{kW}.
\]
Since $\cos\varphi = 0.75$, $\sin\varphi = \sqrt{1 - 0.75^2} \approx 0.661$, so
\[
Q = S\sin\varphi = 1760 \times 0.661 \approx \SI{1.16e3}{var} = \SI{1.16}{kvar}.
\]
Check: $\sqrt{P^2 + Q^2} = \sqrt{1.32^2 + 1.16^2} \approx 1.76\ \mathrm{kVA} = S$. \checkmark

\textbf{(b)} Only the active power does work:
\[
E = P\,t = 1.32\ \mathrm{kW} \times 6\ \mathrm{h} = \SI{7.9}{kWh}.
\]
The consumer is billed for about $\SI{7.9}{kWh}$, not for the reactive energy.
\end{exoresolu}

\courssub{Power factor correction}

Most industrial loads (motors, transformers, fluorescent lamps) are \emph{inductive}: they draw a lagging current and $Q_L > 0$. A capacitor draws a \emph{leading} current, i.e.\ a negative reactive power $Q_C = -V_{\mathrm{rms}}^2\,C\omega$. By connecting a capacitor \emph{in parallel} with the load, the reactive power circulates locally between the coil and the capacitor instead of travelling along the whole line.

\begin{methode}[Choosing the correction capacitor]
\begin{enumerate}
\item Note the initial state: $P$, $\cos\varphi_1$, hence $\tan\varphi_1$ and $Q_1 = P\tan\varphi_1$.
\item Fix the target $\cos\varphi_2$ (often $0.9$ or $1$); compute $Q_2 = P\tan\varphi_2$ (the active power $P$ is unchanged by the parallel capacitor, since an ideal capacitor consumes no average power).
\item The capacitor must supply $Q_C = Q_2 - Q_1 = -V_{\mathrm{rms}}^2 C\omega$, hence
\[
C = \frac{P\left(\tan\varphi_1 - \tan\varphi_2\right)}{V_{\mathrm{rms}}^2\,\omega}.
\]
\item For \emph{full} compensation ($\cos\varphi_2 = 1$), require $Q_C = -Q_L$, i.e.\ at the operating point $\omega C = 1/(\omega L_{\text{eq}})$.
\end{enumerate}
\end{methode}

\begin{exemplebox}[Correcting the workshop motor]
Take the motor above ($P = \SI{1.32}{kW}$, $\cos\varphi_1 = 0.75$, so $\tan\varphi_1 = 0.882$) and correct to $\cos\varphi_2 = 1$ ($\tan\varphi_2 = 0$):
\[
C = \frac{1320 \times 0.882}{220^2 \times 2\pi\times 50} \approx \SI{7.7e-5}{F} \approx \SI{77}{\micro F}.
\]
The line current drops from $\SI{8.0}{A}$ to $I' = P/V_{\mathrm{rms}} = 1320/220 = \SI{6.0}{A}$, cutting the line losses by a factor $(6/8)^2 = 0.56$, i.e.\ a $44\ \%$ reduction.
\end{exemplebox}

\begin{attention}[Common mistakes with a.c. power]
\begin{itemize}
\item Do \emph{not} write $P = V_{\mathrm{rms}}I_{\mathrm{rms}}$ for a general load: the factor $\cos\varphi$ is essential.
\item Do not add powers of different natures arithmetically: use $S^2 = P^2 + Q^2$.
\item The correction capacitor is placed \emph{in parallel}, never in series (a series capacitor would change the voltage applied to the load).
\item Units: $P$ in \si{W}, $Q$ in \si{var}, $S$ in \si{VA} — mixing them costs marks at the GCE.
\end{itemize}
\end{attention}

\begin{exercice}[Resonance and power]
A series $RLC$ circuit has $R = \SI{40}{\ohm}$, $L = \SI{0.10}{H}$, $C = \SI{10}{\micro F}$ and is fed by $V_{\mathrm{rms}} = \SI{120}{V}$.
\begin{enumerate}
\item Calculate $f_0$, $Q$ and $\Delta f$.
\item Calculate the current and the power consumed at resonance.
\item At a frequency where $\cos\varphi = 0.60$, calculate $P$, $Q$ and $S$ given that the current is then $\SI{2.0}{A}$.
\end{enumerate}
\end{exercice}

\begin{corrige}[Exercise 1]
\textbf{1.} $f_0 = \dfrac{1}{2\pi\sqrt{0.10\times 10^{-5}}} = \dfrac{1}{2\pi\times 10^{-3}} \approx \SI{159}{Hz}$; $\omega_0 \approx \SI{1000}{rad.s^{-1}}$; $Q = \dfrac{\omega_0 L}{R} = \dfrac{1000\times0.10}{40} = 2.5$; $\Delta f = f_0/Q \approx \SI{64}{Hz}$.

\textbf{2.} At resonance $I_{\max} = 120/40 = \SI{3.0}{A}$ and $\cos\varphi = 1$, so $P = 120\times3.0 = \SI{360}{W}$.

\textbf{3.} $S = 120\times2.0 = \SI{240}{VA}$; $P = S\cos\varphi = \SI{144}{W}$; $Q = S\sin\varphi = 240\times0.80 = \SI{192}{var}$.
\end{corrige}

\begin{savaistu}[Why ENEO cares about your power factor]
Large factories in Douala and Yaoundé install banks of capacitors next to their motor drives. By raising $\cos\varphi$ towards $1$, they draw less current for the same useful power: the cables heat less, the transformers are relieved, and the utility avoids building extra generating capacity. Resonance, too, is everywhere in daily life: every radio receiver and every mobile phone antenna is an $RLC$ circuit tuned to $f_0 = 1/(2\pi\sqrt{LC})$.
\end{savaistu}

We now understand how a.c.\ circuits select frequencies and how energy really flows in them. The next section uses these ideas — in particular the rms voltage and the notion of an ideal, lossless transfer of power — to study the \cle{transformer}, the device that makes the national high-voltage grid possible.

\coursec{Transformateurs : modèle idéal et pertes réelles}

Dans la section précédente, nous avons étudié la résonance et la puissance en régime sinusoïdal, et nous avons vu que la puissance moyenne fournie à un récepteur est $P = V_{\mathrm{rms}} I_{\mathrm{rms}} \cos\varphi$. Une question se pose naturellement : \emph{comment l'énergie électrique est-elle acheminée économiquement d'une centrale vers un foyer à Bamenda ou Douala ?} La réponse est : à \cle{haute tension} et faible courant, de sorte que les pertes par effet Joule $RI^{2}$ dans les lignes de transport soient faibles. Mais une haute tension est dangereuse et inutilisable dans une habitation. Il faut donc un dispositif capable d'\emph{élever} ou d'\emph{abaisser} une tension alternative sans (idéalement) perdre de puissance. Ce dispositif est le \cle{transformateur}, et il ne fonctionne \emph{qu'avec le courant alternatif} --- conséquence directe des lois de l'induction.

\courssub{Principe : l'induction mutuelle}

\begin{experience}[Observation : deux bobines qui ``dialoguent'' sans contact]
Placez deux bobines face à face. Reliez la bobine 1 à un générateur alternatif et la bobine 2 à une petite lampe. La lampe s'allume bien qu'il n'y ait \emph{aucun contact électrique} entre les bobines. Si l'on insère un noyau en fer à travers les deux bobines, la lampe brille beaucoup plus fort. Si l'on remplace le générateur par une pile (courant continu), la lampe ne s'allume qu'aux instants de fermeture et d'ouverture de l'interrupteur.
\end{experience}

\textbf{Explication.} Le courant alternatif $i_p$ dans la \cle{bobine primaire} crée un flux magnétique alternatif $\Phi$. Le noyau en fer \emph{guide} ce flux de sorte que (presque) la totalité traverse la \cle{bobine secondaire}. Comme le flux à travers la secondaire varie dans le temps, la loi de Faraday y induit une f.é.m. :
\[ e_s = -N_s \frac{\mathrm{d}\Phi}{\mathrm{d}t}, \qquad e_p = -N_p \frac{\mathrm{d}\Phi}{\mathrm{d}t}, \]
où $N_p$ et $N_s$ sont les nombres de spires du primaire et du secondaire. Les deux bobines sont liées par un \emph{flux commun} : c'est l'\cle{induction mutuelle}. On comprend immédiatement pourquoi le courant alternatif est indispensable : avec un courant continu constant, $\mathrm{d}\Phi/\mathrm{d}t = 0$ et aucune f.é.m. n'est induite.

En divisant les deux équations membre à membre (même flux $\Phi$ par spire, le noyau le canalisant) :
\[ \frac{e_p}{e_s} = \frac{N_p}{N_s}. \]
Ce rapport simple est le cœur du transformateur.

\begin{popfigure}
\begin{center}
\begin{tikzpicture}[scale=1.0, every node/.style={font=\small}]
% Noyau
\draw[very thick, popmuted] (0.4,0) rectangle (7.6,4.4);
\draw[very thick, popmuted, fill=white] (1.7,0.9) rectangle (6.3,3.5);
% Enroulement primaire (jambe gauche)
\foreach \y in {1.1,1.5,1.9,2.3,2.7,3.1}{
  \draw[thick, popblue] (0.4,\y) arc (180:0:0.65 and 0.18);
}
\draw[thick, popblue] (0.4,1.1) -- (-1.2,1.1);
\draw[thick, popblue] (0.4,3.28) -- (-1.2,3.28);
\node[popblue, left] at (-1.2,3.28) {$V_p$};
\node[popblue, left] at (-1.2,1.1) {};
\node[popblue, align=center] at (-0.2,2.2) {primaire\\ $N_p$ spires};
\draw[->, thick, popblue] (-1.2,2.9) -- (-0.5,2.9) node[midway, above]{$I_p$};
% Enroulement secondaire (jambe droite)
\foreach \y in {1.5,1.9,2.3,2.7}{
  \draw[thick, poporange] (6.3,\y) arc (0:-180:0.65 and 0.18);
}
\draw[thick, poporange] (7.6,1.5) -- (9.0,1.5);
\draw[thick, poporange] (7.6,2.68) -- (9.0,2.68);
\node[poporange, right] at (9.0,2.68) {$V_s$};
\node[poporange, align=center] at (8.3,3.3) {secondaire\\ $N_s$ spires};
\draw[->, thick, poporange] (8.0,2.1) -- (9.0,2.1) node[midway, above]{$I_s$};
% Flux
\draw[->, thick, poppurple, dashed] (2.6,4.0) -- (5.6,4.0) node[midway, above]{$\Phi$};
\draw[->, thick, poppurple, dashed] (5.6,0.4) -- (2.6,0.4);
\node[popmuted] at (4.0,2.2) {noyau en fer doux};
\end{tikzpicture}
\end{center}
\legende{Transformateur idéal : le flux alternatif $\Phi$ produit par le primaire est guidé par le noyau en fer laminé à travers le secondaire, induisant la f.é.m.\ $V_s$.}
\end{popfigure}

\courssub{Le transformateur idéal}

\begin{definition}[Transformateur idéal]
Un \cle{transformateur idéal} est un transformateur dans lequel :
\begin{itemize}
\item il n'y a \textbf{aucune perte de puissance} (pas d'effet Joule dans les enroulements, pas de pertes dans le noyau) ;
\item le couplage est \textbf{parfait} ($k = 1$) : toute ligne de champ produite par le primaire traverse le secondaire, donc \textbf{pas de flux de fuite} ;
\item les enroulements ont une \textbf{résistance nulle} et le noyau a une perméabilité infinie (courant d'aimantation négligeable).
\end{itemize}
Sous ces hypothèses, les relations du rapport de transformation sont \emph{exactes}.
\end{definition}

\begin{propriete}[Rapport de transformation]
Pour un transformateur idéal, avec des valeurs efficaces (ou crêtes) :
\[ \boxed{\;a = \frac{N_p}{N_s} = \frac{V_p}{V_s} = \frac{I_s}{I_p}\;} \]
La quantité $a$ est le \cle{rapport de transformation}. Si $N_s > N_p$, le transformateur est \textbf{élévateur} ($V_s > V_p$) ; si $N_s < N_p$, il est \textbf{abaisseur}.
\end{propriete}

\textbf{Justification.} Le rapport de tensions $V_p/V_s = N_p/N_s$ découle de la loi de Faraday appliquée au flux commun, comme montré ci-dessus. Pour les courants, on invoque la \emph{conservation de l'énergie} : comme le transformateur idéal n'a pas de pertes, la puissance absorbée par le primaire égale la puissance fournie par le secondaire :
\[ V_p I_p = V_s I_s \quad\Longrightarrow\quad \frac{I_s}{I_p} = \frac{V_p}{V_s} = \frac{N_p}{N_s}. \]
Remarquez l'\emph{inversion} : l'enroulement à la \textbf{plus haute tension} porte le \textbf{plus faible courant}. C'est exactement pourquoi l'élévation de la tension réduit les pertes de transport.

\begin{attention}[Un transformateur ne crée pas de puissance]
Un transformateur élévateur augmente la tension mais \emph{diminue} le courant dans le même rapport : $V_p I_p = V_s I_s$. La puissance (et l'énergie) est conservée. De plus, un transformateur \textbf{ne fonctionne pas en courant continu} : avec un flux constant, il n'y a pas de f.é.m. induite, et la faible résistance du primaire ferait simplement fondre ce dernier.
\end{attention}

\begin{exoresolu}[Transformateur abaisseur pour alimentation de laboratoire]
Un transformateur idéal est raccordé au secteur $\SI{220}{V}$ efficace. Le primaire a $N_p = 1100$ spires. (a) Combien de spires au secondaire pour obtenir $V_s = \SI{12}{V}$ ? (b) Le secondaire alimente une lampe tirant $I_s = \SI{2,0}{A}$. Trouver le courant primaire et la puissance transférée.
\tcblower
(a) $\dfrac{V_p}{V_s} = \dfrac{N_p}{N_s} \Rightarrow N_s = N_p \dfrac{V_s}{V_p} = 1100 \times \dfrac{12}{220} = 60$ spires.

(b) $I_p = I_s \dfrac{N_s}{N_p} = 2,0 \times \dfrac{60}{1100} = \SI{0,109}{A} \approx \SI{0,11}{A}$.

Puissance : $P = V_s I_s = 12 \times 2,0 = \SI{24}{W}$ ; vérification : $V_p I_p = 220 \times 0,109 \approx \SI{24}{W}$. \checkmark
\end{exoresolu}

\courssub{Pertes dans un transformateur réel}

Un transformateur réel n'est jamais parfait. Sa puissance d'entrée se partage entre la puissance utile de sortie et plusieurs mécanismes de pertes.

\textbf{1. Pertes fer (noyau), $P_{\mathrm{fer}}$.} Elles se produisent dans le noyau et sont (presque) indépendantes du courant de charge :
\begin{itemize}
\item \cle{Pertes par hystérésis} : le noyau est aimanté et désaimanté 50 fois par seconde (secteur 50~Hz) ; chaque cycle dissipe une énergie proportionnelle à la surface du cycle $B$--$H$.
\item \cle{Pertes par courants de Foucault} : le flux variable induit des courants circulaires (``courants de Foucault'') dans la masse du fer, qui l'échauffent par effet Joule.
\end{itemize}

\textbf{2. Pertes cuivre, $P_{\mathrm{cu}}$.} Les enroulements primaire et secondaire ont des résistances $R_1$ et $R_2$, si bien que l'effet Joule dissipe
\[ P_{\mathrm{cu}} = R_1 I_p^{2} + R_2 I_s^{2}, \]
ce qui croît avec le \emph{carré} du courant de charge.

\textbf{3. Flux de fuite (dispersion).} Certaines lignes de champ produites par le primaire se referment par l'air sans atteindre le secondaire (et vice versa). Le flux de fuite se comporte comme une inductance série (la \emph{réactance de fuite} $X$), provoquant une chute de tension interne et dégradant à la fois le rendement et la régulation de tension à forte charge.

\begin{attention}[Ne pas négliger le flux de fuite]
Les élèves traitent souvent un transformateur ``bon'' comme idéal même à pleine charge. La réactance de fuite et la résistance des enroulements font \emph{baisser} la tension secondaire quand le courant de charge augmente, et ils réduisent le rendement précisément là où cela compte le plus --- à forte charge.
\end{attention}

\courssub{Comment minimiser les pertes}

\begin{center}
\begin{tabularx}{\linewidth}{|l|X|X|}
\hline
\rowcolor{popblueL} \textbf{Perte} & \textbf{Cause} & \textbf{Remède} \\
\hline
Hystérésis & cycles d'aimantation répétés & utiliser de l'\textbf{acier au silicium} (haute perméabilité, \textbf{faible coercivité}, cycle $B$--$H$ étroit) \\
\hline
Courants de Foucault & courants induits dans le noyau massif & construire le noyau avec de fines \textbf{tôles laminées} (isolées), perpendiculaires au chemin des courants \\
\hline
Pertes cuivre $RI^{2}$ & résistance des enroulements & utiliser du fil de \textbf{section adéquate}, \textbf{cuivre de haute pureté} ; enroulement bifilaire pour réduire résistance et fuite \\
\hline
Flux de fuite & couplage imparfait & enrouler primaire et secondaire sur la \textbf{même jambe}, l'un sur l'autre ; utiliser un noyau fermé à haute perméabilité \\
\hline
\end{tabularx}
\end{center}

\begin{savaistu}[Pourquoi les lames fonctionnent]
Les courants de Foucault circulent en boucles dans le plan perpendiculaire au flux. En tranchant le noyau en fines feuilles vernies, on force les boucles dans des chemins très étroits : la f.é.m. induite par boucle et la section de la boucle s'effondrent toutes deux, et les pertes de Foucault chutent dramatiquement. C'est pour cela que tout noyau de transformateur qu'on ouvre ressemble à une pile de fines plaques en E et en I.
\end{savaistu}

\courssub{Rendement d'un transformateur}

\begin{definition}[Rendement]
Le \cle{rendement} d'un transformateur est
\[ \eta = \frac{P_{\mathrm{out}}}{P_{\mathrm{in}}} = \frac{P_{\mathrm{out}}}{P_{\mathrm{out}} + P_{\mathrm{fer}} + P_{\mathrm{cu}}}, \qquad \eta\,(\%) = \frac{P_{\mathrm{out}}}{P_{\mathrm{in}}} \times 100\,\%. \]
\end{definition}

\begin{methode}[Calcul du rendement à une charge donnée]
\begin{enumerate}
\item Calculer la puissance utile de sortie : $P_{\mathrm{out}} = V_s I_s \cos\varphi$ (avec $\cos\varphi$ le facteur de puissance de la charge).
\item Relever les pertes fer $P_{\mathrm{fer}}$ (constantes, mesurées par l'essai à vide).
\item Calculer les pertes cuivre à cette charge : $P_{\mathrm{cu}} = R_{\mathrm{eq}} I_s^{2}$.
\item Appliquer $\eta = \dfrac{V_s I_s \cos\varphi}{V_s I_s \cos\varphi + P_{\mathrm{fer}} + P_{\mathrm{cu}}} \times 100\,\%$.
\end{enumerate}
\end{methode}

Comme $P_{\mathrm{fer}}$ est constant tandis que $P_{\mathrm{cu}} \propto I_s^{2}$, le rendement est nul à vide, augmente fortement, passe par un \textbf{maximum} (lorsque $P_{\mathrm{cu}} = P_{\mathrm{fer}}$), puis diminue lentement en surcharge. Les bons transformateurs de puissance atteignent $\eta \approx 98\,\%$ près de leur charge nominale.

\begin{popfigure}
\begin{center}
\begin{tikzpicture}
\begin{axis}[
  width=0.72\linewidth, height=6cm,
  xlabel={Courant de charge $I_s$ (A)}, ylabel={$\eta$ (\%)},
  xmin=0, xmax=12, ymin=0, ymax=100,
  xtick={0,2,4,6,8,10,12}, ytick={0,20,40,60,80,100},
  grid=both, grid style={popline, very thin},
  axis line style={popmuted},
]
\addplot[very thick, popblue, smooth, domain=0.15:12, samples=80]
  {100*(9.5*x)/(9.5*x + 4 + 0.09*x*x)};
\addplot[dashed, poporange] coordinates {(6.67,0) (6.67,100)};
\node[poporange, anchor=west, font=\small] at (axis cs:6.9,30) {$\eta_{\max}$ quand $P_{\mathrm{cu}} = P_{\mathrm{fer}}$};
\end{axis}
\end{tikzpicture}
\end{center}
\legende{Courbe typique de rendement d'un transformateur : $\eta$ culmine à la charge pour laquelle pertes cuivre égalent pertes fer.}
\end{popfigure}

\begin{exoresolu}[Rendement d'un transformateur de distribution]
Un transformateur monophasé délivre $V_s = \SI{220}{V}$ sous $I_s = \SI{20}{A}$ à une charge de facteur de puissance $\cos\varphi = 0,85$. Ses pertes fer sont $P_{\mathrm{fer}} = \SI{120}{W}$ et la résistance équivalente d'enroulement rapportée au secondaire est $R_{\mathrm{eq}} = \SI{0,15}{\Omega}$. Calculer le rendement.
\tcblower
Puissance de sortie : $P_{\mathrm{out}} = V_s I_s \cos\varphi = 220 \times 20 \times 0,85 = \SI{3740}{W}$.

Pertes cuivre : $P_{\mathrm{cu}} = R_{\mathrm{eq}} I_s^{2} = 0,15 \times 20^{2} = \SI{60}{W}$.

Rendement :
\[ \eta = \frac{3740}{3740 + 120 + 60} \times 100\,\% = \frac{3740}{3920} \times 100\,\% \approx 95,4\,\%. \]
\end{exoresolu}

\courssub{Régulation de tension}

\begin{propriete}[Chute de tension en charge]
À cause de l'impédance interne des enroulements (résistance équivalente $R_{\mathrm{eq}}$ et réactance de fuite $X_{\mathrm{eq}}$, modélisées en série), la tension aux bornes du secondaire \emph{baisse} quand le courant de charge augmente :
\[ V_{s,\mathrm{fl}} \approx V_{s,\mathrm{nl}} - I_s\,(R_{\mathrm{eq}}\cos\varphi + X_{\mathrm{eq}}\sin\varphi). \]
La \cle{régulation de tension} quantifie cette baisse :
\[ \boxed{\;\mathrm{Régulation} = \frac{V_{\mathrm{nl}} - V_{\mathrm{fl}}}{V_{\mathrm{fl}}} \times 100\,\%\;}
\]
Une faible régulation (quelques pour cent) signifie un transformateur ``raide'' ; une forte réactance de fuite la dégrade.
\end{propriete}

\begin{exemplebox}[Illustration numérique]
Un transformateur donne $\SI{240}{V}$ à vide et $\SI{228}{V}$ à pleine charge. Régulation $= \dfrac{240 - 228}{228} \times 100\,\% \approx 5,3\,\%$.
\end{exemplebox}

\courssub{Essais à vide et en court-circuit}

Pour prévoir le rendement et la régulation, il faut les paramètres du \emph{schéma équivalent}. Deux essais normalisés les fournissent sans charger le transformateur à fond.

\begin{experience}[Essai à vide (circuit ouvert)]
Le secondaire est laissé ouvert ($I_s = 0$) et le primaire est alimenté à sa tension nominale. Un wattmètre sur le primaire lit une faible puissance. Comme $I_s = 0$, les pertes cuivre sont négligeables : le wattmètre lit essentiellement les \textbf{pertes fer} $P_{\mathrm{fer}}$. Cet essai donne aussi les paramètres de la branche d'aimantation et le vrai rapport $V_p/V_s$.
\end{experience}

\begin{experience}[Essai en court-circuit]
Le secondaire est court-circuité et la tension primaire est \emph{réduite} (quelques pour cent de la nominale) jusqu'à ce que le courant atteigne sa valeur nominale. Le flux est alors très faible, donc les pertes fer sont négligeables : le wattmètre lit les \textbf{pertes cuivre} au courant nominal, $P_{\mathrm{cu}} = R_{\mathrm{eq}} I^{2}$. De la tension appliquée $V_{cc}$ et du courant $I_{cc}$, on obtient l'impédance de fuite $Z_{\mathrm{eq}} = V_{cc}/I_{cc}$, d'où $R_{\mathrm{eq}}$ et $X_{\mathrm{eq}}$.
\end{experience}

\begin{propriete}[Ce que mesure chaque essai]
\begin{itemize}
\item \textbf{Essai à vide} $\rightarrow$ pertes fer $P_{\mathrm{fer}}$ et branche d'aimantation.
\item \textbf{Essai en court-circuit} $\rightarrow$ pertes cuivre au courant nominal et l'impédance série (fuite) $R_{\mathrm{eq}} + jX_{\mathrm{eq}}$.
\end{itemize}
Avec ces données, rendement et régulation se calculent pour \emph{toute} charge et facteur de puissance.
\end{propriete}

\begin{aretenir}[Points clés]
\begin{itemize}
\item Un transformateur fonctionne par induction mutuelle à travers un flux commun : il nécessite le \textbf{courant alternatif}.
\item Transformateur idéal : $\dfrac{V_p}{V_s} = \dfrac{N_p}{N_s} = \dfrac{I_s}{I_p}$, avec $V_p I_p = V_s I_s$ (puissance conservée).
\item Pertes réelles : pertes fer (hystérésis + Foucault, constantes), pertes cuivre ($RI^{2}$, variables avec la charge), flux de fuite.
\item Remèdes : noyau en acier au silicium laminé, enroulements en cuivre pur de forte section, couplage serré.
\item $\eta = \dfrac{P_{\mathrm{out}}}{P_{\mathrm{out}} + P_{\mathrm{fer}} + P_{\mathrm{cu}}} \times 100\,\%$ ; maximum quand $P_{\mathrm{cu}} = P_{\mathrm{fer}}$.
\item Régulation $= \dfrac{V_{\mathrm{nl}} - V_{\mathrm{fl}}}{V_{\mathrm{fl}}} \times 100\,\%$, due à $R_{\mathrm{eq}}$ et $X_{\mathrm{eq}}$.
\item Essai à vide $\rightarrow$ $P_{\mathrm{fer}}$ ; essai en court-circuit $\rightarrow$ $P_{\mathrm{cu}}$ et impédance de fuite.
\end{itemize}
\end{aretenir}

\begin{exercice}[Rendement et régulation]
Un transformateur de $\SI{5}{kVA}$ a des pertes fer de $\SI{60}{W}$ et, d'après l'essai en court-circuit, des pertes cuivre à pleine charge de $\SI{140}{W}$. (a) Trouver son rendement à pleine charge avec $\cos\varphi = 1$ et avec $\cos\varphi = 0,8$. (b) À quelle fraction de la charge nominale le rendement est-il maximal ? (c) Sa tension secondaire à vide est $\SI{235}{V}$ et à pleine charge $\SI{225}{V}$ : calculer la régulation.
\end{exercice}

\begin{corrige}[Exercice 1]
(a) $\cos\varphi = 1$ : $P_{\mathrm{out}} = \SI{5000}{W}$ ; $\eta = \dfrac{5000}{5000 + 60 + 140} \times 100\,\% = \dfrac{5000}{5200} \times 100\,\% \approx 96,2\,\%$.

$\cos\varphi = 0,8$ : $P_{\mathrm{out}} = \SI{4000}{W}$ ; $\eta = \dfrac{4000}{4200} \times 100\,\% \approx 95,2\,\%$.

(b) Rendement maximal quand $P_{\mathrm{cu}} = P_{\mathrm{fer}}$ : les pertes cuivre varient comme $x^{2}$, donc $x^{2} \times 140 = 60 \Rightarrow x = \sqrt{60/140} \approx 0,65$, soit environ $65\,\%$ de la charge nominale.

(c) Régulation $= \dfrac{235 - 225}{225} \times 100\,\% \approx 4,4\,\%$.
\end{corrige}

\coursec{Redressement et lissage des signaux alternatifs}

Dans les sections précédentes, nous avons vu comment un transformateur abaisse la tension secteur \SI{220}{V} vers une tension de sécurité utilisable. Mais le secondaire d'un transformateur délivre encore une tension \emph{alternative}, alors que pratiquement tous les appareils électroniques autour de nous --- le chargeur de téléphone dans votre poche, la radio d'un taxi à Douala, l'ordinateur du laboratoire --- ont besoin d'une tension \emph{continue et stable} pour fonctionner. L'étage final de toute alimentation est donc la conversion du courant alternatif en courant continu : c'est le \cle{redressement}, suivi du \cle{lissage} et de la \cle{régulation}. Cette section étudie chacun de ces étages et se termine par la conception complète d'une petite alimentation, exercice classique de l'examen GCE.

\courssub{Redressement : principe et redresseur à demi-onde}

\begin{definition}[Redressement]
Le \cle{redressement} est la conversion d'une tension alternative (qui change de polarité périodiquement) en une tension \cle{unilatérale} (pulsatoire), c'est-à-dire une tension qui conserve la même polarité en tout instant, même si son amplitude varie.
\end{definition}

Le composant clé est la \cle{diode} : un dispositif semi-conducteur qui ne laisse passer le courant que dans un seul sens (de l'anode vers la cathode) lorsqu'elle est polarisée en direct, et le bloque en inverse. Une diode idéale se comporte comme un interrupteur fermé si $V_{\text{anode}} > V_{\text{cathode}}$ et ouvert sinon. Une diode silicium réelle a besoin d'environ \SI{0.7}{V} de tension directe pour conduire.

\textbf{Le redresseur à demi-onde.}\par
Le redresseur le plus simple utilise une seule diode en série avec la résistance de charge $R$. Pendant la demi-période positive de l'entrée $v(t) = V_p\sin(\omega t)$, la diode conduit et la tension aux bornes de la charge suit l'entrée ; pendant la demi-période négative, la diode bloque et la tension de charge est nulle. La sortie est donc une succession de demi-arches sinusoïdales positives séparées par des intervalles nuls.

La valeur moyenne (composante continue) de cette forme d'onde s'obtient en moyennant sur une période complète $T$ :
\begin{align*}
V_{dc} &= \frac{1}{T}\int_0^T v(t)\,dt = \frac{1}{T}\int_0^{T/2} V_p\sin(\omega t)\,dt = \frac{V_p}{T}\left[-\frac{\cos(\omega t)}{\omega}\right]_0^{T/2} = \frac{V_p}{\pi}.
\end{align*}

\begin{propriete}[Redresseur à demi-onde]
Pour un redresseur à demi-onde alimenté par $v = V_p\sin(\omega t)$ :
\[
V_{dc} = \frac{V_p}{\pi} \approx 0.318\,V_p, \qquad r = 1.21,
\]
où $r$ est le facteur de ripple (ondulation). La moitié de l'énergie d'entrée est gaspillée et la sortie est nulle la moitié du temps : ce circuit est peu performant et rarement utilisé seul.
\end{propriete}

\courssub{Redresseurs à double alternance : à prise médiane et en pont}

\textbf{Redresseur à prise médiane.}\par
Un transformateur à prise médiane au secondaire fournit deux tensions $v_1$ et $v_2$ de même amplitude mais de phases opposées. Deux diodes relient les deux extrémités du secondaire à la charge, dont l'autre borne revient à la prise médiane. Chaque diode conduit pendant une demi-période, si bien que la charge reçoit du courant pendant \emph{les deux} demi-périodes, toujours dans le même sens. La sortie est une succession d'arches positives identiques à une fréquence double de celle du secteur.

\begin{propriete}[Redressement à double alternance]
Pour les deux redresseurs à double alternance (prise médiane et pont) :
\[
V_{dc} = \frac{2V_p}{\pi} \approx 0.637\,V_p, \qquad r = 0.48.
\]
La fréquence de l'ondulation est $2f$ (par exemple \SI{100}{Hz} pour un secteur \SI{50}{Hz}), ce qui facilite grandement le lissage par rapport au cas demi-onde.
\end{propriete}

\textbf{Le redresseur en pont (pont de Graetz).}\par
Le redresseur en pont utilise quatre diodes disposées en losange et ne nécessite \emph{pas} de transformateur à prise médiane --- avantage économique décisif. Pendant la demi-période positive, les diodes $D_1$ et $D_3$ conduisent ; pendant la demi-période négative, $D_2$ et $D_4$ conduisent. Le courant dans la charge circule toujours dans le même sens.

\begin{popfigure}
\begin{center}
\begin{tikzpicture}[european, scale=1.0, transform shape]
% AC source
\draw (0,0) to[sV, l={$v_s$}] (0,3);
\draw (0,3) -- (1.5,3);
\draw (0,0) -- (1.5,0);
% Bridge diamond
\coordinate (A) at (1.5,1.5);
\coordinate (B) at (3,3);
\coordinate (C) at (4.5,1.5);
\coordinate (D) at (3,0);
\draw (1.5,3) -- (B);
\draw (1.5,0) -- (D);
\draw (B) to[D, l_={$D_1$}] (C);
\draw (D) to[D, l_={$D_2$}] (C);
\draw (A) to[D, l^={$D_4$}] (B);
\draw (A) to[D, l_={$D_3$}] (D);
\draw (1.5,3) -- (A) -- (1.5,0);
% Output with smoothing capacitor and load
\draw (C) -- (6.5,1.5);
\draw (6.5,3) -- (C);
\draw (C) -- (6.5,0);
\draw (6.5,3) to[C, l={$C$}] (6.5,0);
\draw (6.5,3) -- (8.5,3) to[R, l={$R$}] (8.5,0) -- (6.5,0);
\node[popblue] at (8.5,1.5) [right=0.15cm, align=center] {$u_R$};
\node[popmuted] at (3,-0.7) {Pont redresseur + condensateur de lissage};
\end{tikzpicture}
\end{center}
\legende{Redresseur en pont à quatre diodes, condensateur de lissage $C$ en parallèle avec la charge résistive $R$.}
\end{popfigure}

\begin{propriete}[Tension inverse de crête (PIC) des diodes]
La \cle{tension inverse de crête} (PIC ou PIV en anglais) est la tension inverse maximale que chaque diode doit supporter sans claquage :
\begin{center}
\begin{tabular}{|l|c|}
\hline
\rowcolor{popblueL} Circuit & PIC par diode \\
\hline
Demi-onde & $V_p$ \\
Pont (Graetz) & $V_p$ \\
Prise médiane & $2V_p$ \\
\hline
\end{tabular}
\end{center}
Noter que dans le pont, \emph{deux} diodes conduisent en série à chaque instant, donc la tension de crête en sortie est réduite de deux chutes de diode : $V_{p,\text{out}} \approx V_p - 1.4\ \mathrm{V}$ pour des diodes silicium.
\end{propriete}

\begin{attention}[Chutes de tension des diodes non négligeables à basse tension]
Oublier la chute de \SI{0.7}{V} par diode conductrice est une erreur classique. Si la tension de crête secondaire n'est que \SI{9}{V}, un pont perd \SI{1.4}{V}, soit plus de $15\,\%$ de la sortie. Soustrayez toujours les chutes de diodes avant de calculer $V_{dc}$, surtout dans les exercices de dimensionnement.
\end{attention}

\begin{popfigure}
\begin{center}
\begin{tikzpicture}[scale=0.9]
% Axes
\draw[->, popline] (0,0) -- (10.6,0) node[below left]{$t$};
\draw[->, popline] (0,-2.6) -- (0,2.9) node[left]{$u$};
% AC input (thin gray)
\draw[popmuted, domain=0:10, samples=200, smooth] plot (\x, {2.2*sin(180*\x/2.5)});
\node[popmuted] at (9.6,-2.3) {entrée c.a.};
% Rectified (bridge)
\draw[popblue, thick, domain=0:10, samples=300, smooth] plot (\x, {abs(2.2*sin(180*\x/2.5))});
\node[popblue] at (5.0,2.55) {après pont};
% Smoothed (approximate exponential discharge)
\draw[poporange, very thick] (0,0)
 \foreach \k in {0,1,2,3} {
  -- plot[domain=0:0.625, samples=40] ({\k*2.5+\x}, {2.2*sin(180*(\k*2.5+\x)/2.5)})
  -- plot[domain=0.625:2.5, samples=40] ({\k*2.5+\x}, {2.2*exp(-(\x-0.625)/1.1)})
 };
\node[poporange] at (8.3,1.15) {après lissage};
% Ripple annotation
\draw[<->, poppink, thick] (7.9,{2.2*exp(-(7.9-7.5-0.625+2.5)/1.1)}) -- (7.9,2.2) node[midway, right]{$\Delta V$};
\end{tikzpicture}
\end{center}
\legende{Formes d'onde : entrée sinusoïdale (gris), sortie redressée double alternance (bleu), et tension lissée avec ripple $\Delta V$ (orange).}
\end{popfigure}

\courssub{Lissage par condensateur}

Une tension pulsatoire n'est pas encore du continu. Le remède est un gros \cle{condensateur de lissage} $C$ placé en parallèle sur la charge. Le mécanisme est simple :
\begin{itemize}
\item \textbf{Charge :} au voisinage de chaque pic de la tension redressée, les diodes conduisent et rechargent rapidement le condensateur jusqu'à (presque) la valeur de crête $V_p$.
\item \textbf{Décharge :} lorsque la tension redressée tombe en dessous de la tension du condensateur, les diodes se bloquent ; le condensateur se décharge \emph{lentement} à travers la charge $R$, assurant lui-même le courant de charge.
\end{itemize}
La sortie oscille donc entre $V_p$ et $V_p - \Delta V$ : le ripple $\Delta V$ est petit si la décharge est lente, c'est-à-dire si $RC \gg T$.

Pendant la décharge, $u_C(t) = V_p\,e^{-t/(RC)}$. Si le temps de décharge $\approx T$ est très inférieur à $RC$, l'exponentielle est quasi-linéaire et la charge perdue par le condensateur est $\Delta Q = C\,\Delta V = I_{\text{charge}}\,T$, ce qui donne la formule pratique :

\begin{propriete}[Ripple d'une alimentation lissée par condensateur]
\[
\Delta V \approx \frac{I_{\text{charge}}}{f_r\,C},
\]
où $f_r$ est la fréquence de l'ondulation : $f_r = f$ pour le redressement demi-onde et $f_r = 2f$ pour le redressement double alternance. D'où la capacitance requise pour un ripple maximal imposé :
\[
C \geq \frac{I_{\text{charge}}}{f_r\,\Delta V_{\max}}.
\]
\end{propriete}

\begin{methode}[Facteur de ripple avec condensateur]
Le \cle{facteur de ripple} mesure le contenu alternatif résiduel :
\[
r = \frac{V_{\text{ripple (r.m.s.)}}}{V_{dc}}.
\]
Pour un redresseur en pont avec condensateur de lissage, en modélisant le ripple par une onde triangulaire, on obtient
\[
r \approx \frac{1}{4\sqrt{3}\,f\,R_{\text{charge}}\,C} \quad (\text{redressement double alternance, } f \text{ fréquence secteur}).
\]
Une bonne alimentation lissée atteint $r < 0.05$ ; après régulation, $r$ peut descendre sous $10^{-3}$.
\end{methode}

\begin{attention}[Tension de service du condensateur]
Le condensateur se charge jusqu'à la valeur de \emph{crête} $V_p = \sqrt{2}\,V_{\text{eff}}$, et non la valeur efficace. Sa tension nominale (de service) doit satisfaire $V_{\text{nom}} \geq 1.5\,V_p$ pour garantir une marge de sécurité. Brancher un condensateur \SI{16}{V} sur un secondaire \SI{12}{V} efficace ($V_p \approx \SI{17}{V}$) le détruit, parfois de façon explosive.
\end{attention}

\courssub{Régulation de tension}

Même lissée, la tension varie encore avec le courant de charge et les fluctuations du secteur. Un \cle{régulateur} délivre une tension de sortie constante :
\begin{itemize}
\item \textbf{Régulateur à diode Zener :} une diode Zener en avalanche maintient une tension quasi constante $V_Z$ aux bornes de la charge ; une résistance série absorbe l'excès de tension. Simple, mais limité aux faibles courants.
\item \textbf{Régulateur linéaire (famille 78xx) :} circuit intégré trois broches (ex. 7805 pour \SI{+5}{V}, 7812 pour \SI{+12}{V}). L'entrée doit dépasser la sortie d'environ \SI{2}{V} à \SI{3}{V}. Sortie très propre, mais la puissance excédentaire est dissipée en chaleur.
\item \textbf{Alimentation à découpage (SMPS) :} la tension est hachée à haute fréquence et re-filtrée, avec des rendements supérieurs à $80\,\%$. C'est la technologie des chargeurs de téléphone et adaptateurs d'ordinateurs portables modernes.
\end{itemize}

\begin{savaistu}[Pourquoi les chargeurs de téléphone sont-ils si légers aujourd'hui ?]
Les anciens chargeurs contenaient un gros transformateur fer \SI{50}{Hz} lourd. Les chargeurs SMPS modernes redressent d'abord le secteur, puis le hachent à plusieurs dizaines de kHz : à haute fréquence le transformateur peut être minuscule, d'où le format « allumette » d'un chargeur \SI{20}{W} actuel.
\end{savaistu}

\begin{methode}[Choix du transformateur pour une alimentation à pont avec lissage]
Pour une tension continue de sortie $V_{dc}$ requise, courant $I_{dc}$ :
\begin{enumerate}
\item Tension de crête nécessaire au condensateur : $V_p \approx V_{dc} + \frac{\Delta V}{2} + 2\times\SI{0.7}{V}$ (deux chutes de diode).
\item Tension efficace secondaire : $V_{\text{sec}} = \dfrac{V_p}{\sqrt{2}}$. (Sans condensateur, on utiliserait $V_{dc} = \dfrac{2V_p}{\pi}$, soit $V_{\text{sec}} = \dfrac{\pi\,V_{dc}}{2\sqrt{2}}$.)
\item Puissance apparente : $S \geq \dfrac{P_{\text{charge}}}{\eta_{\text{transfo}}}$, avec $\eta \approx 0.8$ à $0.95$ pour un bon transformateur.
\item Diodes : courant nominal $> I_{\text{charge}}$ (avec marge pour les pics de recharge du condensateur) et $\mathrm{PIC} > V_p$.
\end{enumerate}
\end{methode}

\begin{exoresolu}[Dimensionnement d'une alimentation 12 V, 0,5 A]
Une alimentation de laboratoire doit délivrer $V_{dc} = \SI{12}{V}$ sous $I = \SI{0.5}{A}$ à partir du secteur \SI{220}{V}, \SI{50}{Hz}, en utilisant un pont redresseur et un condensateur de lissage, avec un ripple $\Delta V \leq \SI{1.0}{V}$. Déterminer : (a) la tension secondaire du transformateur ; (b) la capacitance $C$ ; (c) la PIC des diodes ; (d) la tension de service du condensateur.
\tcblower
\textbf{(a) Tension secondaire.} Le condensateur se charge à la crête moins deux chutes de diode. La tension moyenne lissée vaut environ $V_p - 1.4 - \Delta V/2$. Donc :
\[
V_p = V_{dc} + \frac{\Delta V}{2} + 1.4 = 12 + 0.5 + 1.4 = \SI{13.9}{V},
\qquad V_{\text{sec}} = \frac{V_p}{\sqrt{2}} \approx \SI{9.8}{V}\ \text{eff.}
\]
Un secondaire standard \SI{12}{V} donne une marge confortable (nécessaire de toute façon si un régulateur 7812 suit).

\textbf{(b) Capacitance.} Redressement double alternance : $f_r = 2f = \SI{100}{Hz}$.
\[
C \geq \frac{I}{f_r\,\Delta V} = \frac{0.5}{100 \times 1.0} = 5.0\times10^{-3}\ \mathrm{F} = \SI{5000}{\micro F}.
\]
On choisira la valeur normalisée \SI{4700}{\micro F} ou \SI{6800}{\micro F}.

\textbf{(c) PIC.} Avec un secondaire \SI{12}{V} eff., $V_p = 12\sqrt{2} \approx \SI{17.0}{V}$. Pour un pont, $\mathrm{PIC} = V_p \approx \SI{17}{V}$ ; on choisira des diodes supportant au moins \SI{50}{V}, \SI{1}{A} (ex. 1N4007).

\textbf{(d) Tension du condensateur.} $V_{\text{nom}} \geq 1.5\,V_p = 1.5 \times 17.0 \approx \SI{25.5}{V}$ : on choisira un condensateur électrolytique \SI{25}{V} ou, plus sûr, \SI{35}{V}.
\end{exoresolu}

\begin{methode}[Réduire le ripple sans augmenter $C$]
Les très gros condensateurs sont encombrants et coûteux. Alternatives :
\begin{enumerate}
\item Utiliser un \textbf{filtre LC} (self en série, condensateur en parallèle) : la self s'oppose aux variations de courant.
\item Ajouter un \textbf{régulateur} (78xx ou SMPS) après un lissage modéré : le régulateur rejette électroniquement le ripple résiduel.
\item Augmenter la fréquence de l'ondulation (SMPS à haute fréquence) : comme $\Delta V \propto 1/f$, un petit condensateur suffit alors.
\end{enumerate}
\end{methode}

\begin{exercice}[Demi-onde versus pont]
Une tension sinusoïdale \SI{50}{Hz} de valeur de crête $V_p = \SI{20}{V}$ alimente (i) un redresseur à demi-onde, (ii) un redresseur en pont, tous deux chargés par $R = \SI{100}{\ohm}$ avec un condensateur de lissage $C = \SI{1000}{\micro F}$. Pour chaque cas, calculer $V_{dc}$ (non lissé), le ripple $\Delta V$ (lissé), et la PIC des diodes. Commenter les résultats.
\end{exercice}

\begin{corrige}[Exercice 1]
\textbf{Demi-onde :}
\begin{itemize}
\item $V_{dc} = V_p/\pi = 20/\pi \approx \SI{6.37}{V}$ (non lissé, diodes idéales).
\item Courant de charge moyen estimé $I_{dc} \approx V_{dc}/R = \SI{63.7}{mA}$.
\item Fréquence de ripple $f_r = \SI{50}{Hz}$.
\item $\Delta V \approx \dfrac{I_{dc}}{f_r C} = \dfrac{0.0637}{50 \times 10^{-3}} \approx \SI{1.27}{V}$.
\item $\mathrm{PIC} = V_p = \SI{20}{V}$.
\end{itemize}

\textbf{Pont (double alternance) :}
\begin{itemize}
\item $V_{dc} = 2V_p/\pi \approx \SI{12.73}{V}$ (non lissé, diodes idéales). En tenant compte des chutes de diodes, la crête de sortie est $V_{p,\text{out}} \approx 20 - 1.4 = \SI{18.6}{V}$.
\item Courant de charge moyen estimé $I_{dc} \approx V_{dc}/R = \SI{127.3}{mA}$ (basé sur la valeur moyenne non lissée, estimation conservatrice pour le dimensionnement du condensateur).
\item Fréquence de ripple $f_r = \SI{100}{Hz}$.
\item $\Delta V \approx \dfrac{I_{dc}}{f_r C} = \dfrac{0.1273}{100 \times 10^{-3}} \approx \SI{1.27}{V}$.
\item $\mathrm{PIC} = V_p = \SI{20}{V}$ par diode.
\end{itemize}

\textbf{Commentaire :} Le pont double la tension moyenne (non lissée) \emph{et} double la fréquence de l'ondulation. Pour le même condensateur, le ripple $\Delta V$ reste de même ordre de grandeur car le courant double aussi, mais le facteur de ripple $r = \Delta V/V_{dc}$ est divisé par deux. Le pont est supérieur sur tous les plans sauf le coût de deux diodes supplémentaires.
\end{corrige}

\begin{aretenir}[Points clés de la section]
\begin{itemize}
\item Le redressement convertit le c.a. en tension unilatérale pulsatoire à l'aide de diodes.
\item Demi-onde : $V_{dc} = V_p/\pi$, $r = 1.21$. Double alternance (pont ou prise médiane) : $V_{dc} = 2V_p/\pi$, $r = 0.48$.
\item PIC : $V_p$ pour demi-onde et pont ; $2V_p$ pour prise médiane.
\item Condensateur de lissage : se charge aux pics, se décharge dans la charge ; $\Delta V \approx I_{\text{charge}}/(f_r C)$ avec $f_r = 2f$ pour la double alternance.
\item Règles de dimensionnement : $C \geq I/(f_r\,\Delta V_{\max})$ ; condensateur noté $\geq 1.5\,V_p$ ; soustraire \SI{0.7}{V} par diode conductrice ; un régulateur (Zener, 78xx, SMPS) stabilise la sortie finale.
\end{itemize}
\end{aretenir}

\newpage

\coursec{Integration activity}

\courssub{Aims of the activity}

By the end of this integration activity, you should be able to:

\begin{itemize}
\item apply the ideal transformer relations $\dfrac{V_p}{V_s}=\dfrac{N_p}{N_s}=\dfrac{I_s}{I_p}$ to choose a suitable turns ratio;
\item select a full-wave bridge rectifier and justify the peak inverse voltage (PIV) rating of its diodes;
\item dimension a smoothing capacitor so that the ripple voltage stays below a given limit at a given load current;
\item analyse the thermal behaviour of a linear voltage regulator (7805) and decide whether a heat sink is required;
\item estimate the overall efficiency of a complete d.c.\ power supply and discuss voltage regulation;
\item present a technical report containing a circuit diagram, calculations and justified choices of commercial components.
\end{itemize}

This activity mobilises all the key competences of the chapter: \cle{rms values}, \cle{impedance}, \cle{power in a.c.\ circuits}, \cle{transformers}, \cle{rectification} and \cle{smoothing}.

\courssub{Situation}

\textbf{Design of a 5 V / 1 A phone-charger power supply.}\par

A young technician in Bamenda wants to set up a small workshop that repairs mobile phones. Because power cuts are frequent, customers often bring chargers that have been damaged by voltage fluctuations. He decides to build, from scratch, a simple and robust mains-powered d.c.\ supply delivering a \textbf{regulated 5 V at up to 1 A}, which can serve as a universal phone charger and as a bench supply for testing small circuits.

He buys components at the electronics market and collects the following data.

\textbf{Available mains supply (ENEO network):}
\begin{itemize}
\item rms voltage $V_p = 230\ \mathrm{V}$, frequency $f = 50\ \mathrm{Hz}$;
\item the voltage may vary by $\pm 10\ \%$ around its nominal value.
\end{itemize}

\textbf{Block diagram of the supply to be designed:}
\begin{center}
mains $\rightarrow$ \textbf{transformer} $\rightarrow$ \textbf{bridge rectifier} $\rightarrow$ \textbf{smoothing capacitor} $\rightarrow$ \textbf{regulator 7805} $\rightarrow$ output $5\ \mathrm{V}$.
\end{center}

\textbf{Technical data:}
\begin{itemize}
\item The transformer is assumed \emph{ideal} for the turns-ratio calculation; its real efficiency is about $\eta_T \approx 90\ \%$.
\item Each silicon diode of the bridge has a forward voltage drop $V_D = 0.7\ \mathrm{V}$ when conducting. In a bridge rectifier, two diodes conduct at a time.
\item The smoothing capacitor $C$ discharges approximately at constant current $I$ during about one half-period of the ripple, $\Delta t \approx \dfrac{1}{2f}$, between two recharging peaks (full-wave rectification, ripple frequency $f_{\text{ripple}} = 2f = 100\ \mathrm{Hz}$).
\item The 7805 regulator requires an input voltage at least $V_{i,\min} = 7.5\ \mathrm{V}$ to regulate correctly (dropout condition). Its quiescent current is negligible, so $I_{\text{in}} \approx I_{\text{out}} = 1\ \mathrm{A}$.
\item Target specifications: output $V_{\text{out}} = 5.0\ \mathrm{V}$, $I_{\text{out}} = 1.0\ \mathrm{A}$, ripple at the capacitor $\Delta V \leq 0.5\ \mathrm{V}$.
\end{itemize}

Working in groups, you must produce a written report (with circuit diagram, calculations and choice of real commercial components) answering the tasks below.

\begin{popfigure}
\begin{tikzpicture}[european, scale=0.9, transform shape]
\draw (0,0) to[sV=$230\ \mathrm{V}$] (0,2.4);
\draw (0,2.4) -- (1.0,2.4);
\draw (0,0) -- (1.0,0);
\draw[thick] (1.0,0) -- (1.0,2.4);
\draw[thick] (1.15,0) -- (1.15,2.4);
\draw[thick] (1.3,0) -- (1.3,2.4);
\draw (1.3,2.4) -- (2.3,2.4);
\draw (1.3,0) -- (2.3,0);
\node at (1.05,2.7) {\scriptsize $N_p$};
\node at (1.45,2.7) {\scriptsize $N_s$};
\draw (2.3,2.4) to[D=$D_1$] (4.3,2.4);
\draw (2.3,0) to[D,l_={$D_3$}] (4.3,0);
\draw (4.3,2.4) -- (4.3,1.2);
\draw (4.3,0) -- (4.3,1.2);
\draw (4.3,1.2) to[C=$C$] (6.3,1.2);
\draw (6.3,1.2) -- (7.3,1.2);
\node[draw, rectangle, minimum width=1.4cm, minimum height=0.8cm] at (8.4,1.2) {\scriptsize 7805};
\draw (9.1,1.2) -- (10.3,1.2);
\node at (10.3,1.55) {\scriptsize $5\ \mathrm{V}$ out};
\node at (3.3,2.75) {\scriptsize bridge};
\node at (5.3,1.55) {\scriptsize smoothing};
\end{tikzpicture}
\legende{Simplified block-level circuit of the linear 5 V / 1 A power supply: step-down transformer, diode bridge, smoothing capacitor and 7805 regulator.}
\end{popfigure}

\courssub{Tasks}

\begin{enumerate}
\item \textbf{Choice of the transformer.} The technician decides that the secondary should deliver about $V_s = 9\ \mathrm{V}$ (rms). Calculate the turns ratio $N_p/N_s$ of the ideal transformer. If the primary has $N_p = 1150$ turns, find $N_s$.
\item \textbf{Rectifier.} Draw the circuit diagram of a bridge rectifier connected to the secondary. Calculate the peak value of the secondary voltage, then the peak voltage available at the output of the bridge (take the two diode drops into account). Determine the peak inverse voltage (PIV) that each diode must withstand and propose a suitable commercial diode (1N4007: PIV $= 1000\ \mathrm{V}$, $I_F = 1\ \mathrm{A}$ — is it acceptable?).
\item \textbf{Smoothing capacitor.} Using $\Delta V \approx \dfrac{I\,\Delta t}{C}$ with $\Delta t \approx \dfrac{1}{2f} = 10\ \mathrm{ms}$, $I = 1\ \mathrm{A}$ and $\Delta V = 0.5\ \mathrm{V}$, calculate the minimum capacitance $C$. Choose a standard value. Then verify that the minimum voltage at the regulator input, $V_{C,\min} = V_{\text{peak,bridge}} - \Delta V$, satisfies the dropout condition of the 7805.
\item \textbf{Regulator and thermal check.} Calculate the average power dissipated by the 7805, $P_{7805} \approx (V_{C,\text{avg}} - 5)\,I_{\text{out}}$, where $V_{C,\text{avg}} = V_{\text{peak,bridge}} - \Delta V/2$. The 7805 in a TO-220 package without heat sink can dissipate safely about $1\ \mathrm{W}$; conclude on the need for a heat sink.
\item \textbf{Efficiency and regulation.} Estimate the useful output power, the power lost in the regulator, and the overall efficiency $\eta$ of the supply (include the transformer efficiency $\eta_T = 90\ \%$; neglect losses in the diodes and capacitor). Explain briefly why the output voltage remains at $5\ \mathrm{V}$ even if the mains drops by $10\ \%$ (check the dropout condition with $V_p = 207\ \mathrm{V}$).
\item \textbf{Report.} Write the final report: circuit diagram with all component values, table of chosen commercial components, and a short discussion of two possible improvements (e.g.\ fuse, LED indicator, larger heat sink, switching supply).
\end{enumerate}

\courssub{Model answer}

\textbf{1. Turns ratio.}\par
For an ideal transformer, the ratio of the rms voltages equals the ratio of the numbers of turns:
\[
\frac{N_p}{N_s} = \frac{V_p}{V_s} = \frac{230}{9} \approx 25.6.
\]
With $N_p = 1150$ turns:
\[
N_s = \frac{N_p}{25.6} = \frac{1150}{25.6} \approx 45\ \text{turns}.
\]
A commercial $230\ \mathrm{V}/9\ \mathrm{V}$, $12\ \mathrm{VA}$ transformer is suitable: it must supply at least $9\ \mathrm{V} \times 1.1\ \mathrm{A} \approx 10\ \mathrm{VA}$, with a small safety margin.

\textbf{2. Bridge rectifier.}\par
The peak secondary voltage is obtained from the rms value:
\[
\hat{V}_s = V_s\sqrt{2} = 9 \times 1.414 \approx 12.7\ \mathrm{V}.
\]
At any instant, two diodes of the bridge conduct in series with the load, so the peak voltage available across the capacitor is
\[
V_{\text{peak,bridge}} = \hat{V}_s - 2V_D = 12.7 - 1.4 = 11.3\ \mathrm{V}.
\]
In a bridge rectifier, each non-conducting diode must block a reverse voltage approximately equal to the peak of the secondary voltage:
\[
\text{PIV} \approx \hat{V}_s \approx 12.7\ \mathrm{V}.
\]
The 1N4007 (PIV $= 1000\ \mathrm{V}$, $I_F = 1\ \mathrm{A}$) is therefore \emph{largely acceptable} as far as reverse voltage is concerned: its PIV rating is nearly 80 times the required value. Its forward current rating, however, is exactly equal to the load current, which leaves no margin for the large repetitive charging current peaks of the capacitor; a $2\ \mathrm{A}$ bridge such as the KBP206 gives an extra safety margin and is preferred for continuous 1 A service.

\textbf{3. Smoothing capacitor.}\par
Between two charging peaks, the capacitor alone supplies the load. For a \emph{full-wave} rectifier the ripple frequency is $2f = 100\ \mathrm{Hz}$, so the discharge lasts about
\[
\Delta t = \frac{1}{2f} = \frac{1}{100} = 1.0 \times 10^{-2}\ \mathrm{s} = 10\ \mathrm{ms}.
\]
Assuming discharge at constant current $I$, the voltage falls by $\Delta V = \dfrac{I\,\Delta t}{C}$, hence
\[
C \geq \frac{I\,\Delta t}{\Delta V} = \frac{1.0 \times 1.0 \times 10^{-2}}{0.5} = 2.0 \times 10^{-2}\ \mathrm{F} = 20\,000\ \mu\mathrm{F}.
\]
Choose the standard value $C = 22\,000\ \mu\mathrm{F}$, rated at $25\ \mathrm{V}$ (working voltage well above the $11.3\ \mathrm{V}$ peak).\par
Dropout check — the \emph{minimum} of the ripple must stay above $7.5\ \mathrm{V}$:
\[
V_{C,\min} = V_{\text{peak,bridge}} - \Delta V = 11.3 - 0.5 = 10.8\ \mathrm{V} > 7.5\ \mathrm{V}.
\]
The 7805 regulates correctly. \checkmark

\begin{attention}[A frequent numerical trap]
A very common error here is to write $C = 2000\ \mu\mathrm{F}$. Check the powers of ten: $I\,\Delta t = 1.0 \times 10^{-2}\ \mathrm{C}$, and dividing by $\Delta V = 0.5\ \mathrm{V}$ \emph{doubles} the result: $C = 2.0 \times 10^{-2}\ \mathrm{F} = 20\,000\ \mu\mathrm{F}$. Large reservoir capacitors are normal in $50\ \mathrm{Hz}$ linear supplies — this is precisely why they are bulky.
\end{attention}

\textbf{4. Thermal check of the 7805.}\par
The average capacitor voltage is
\[
V_{C,\text{avg}} = V_{\text{peak,bridge}} - \frac{\Delta V}{2} = 11.3 - 0.25 = 11.05\ \mathrm{V} \approx 11.1\ \mathrm{V}.
\]
Since the regulator carries the full output current, the power it dissipates as heat is
\[
P_{7805} = (V_{C,\text{avg}} - V_{\text{out}})\,I_{\text{out}} = (11.05 - 5.0) \times 1.0 \approx 6.1\ \mathrm{W}.
\]
This is far above the $\approx 1\ \mathrm{W}$ limit of a bare TO-220 package: \textbf{a heat sink is compulsory}. A heat sink of thermal resistance below about $10\ \mathrm{K\,W^{-1}}$ keeps the junction temperature within safe limits; the regulator's internal thermal shutdown provides extra protection.

\textbf{5. Efficiency and regulation.}\par
Useful output power:
\[
P_{\text{out}} = V_{\text{out}} I_{\text{out}} = 5.0 \times 1.0 = 5.0\ \mathrm{W}.
\]
Power drawn from the transformer secondary (approximately $V_{C,\text{avg}} \times I_{\text{out}}$):
\[
P_{\text{sec}} \approx 11.05 \times 1.0 \approx 11.1\ \mathrm{W}.
\]
Including the transformer efficiency, the power drawn from the mains is
\[
P_{\text{mains}} = \frac{P_{\text{sec}}}{\eta_T} = \frac{11.1}{0.90} \approx 12.3\ \mathrm{W}.
\]
Overall efficiency:
\[
\eta = \frac{P_{\text{out}}}{P_{\text{mains}}} = \frac{5.0}{12.3} \approx 0.41 \approx 41\ \%.
\]
More than half of the energy is lost, mainly as heat in the 7805 ($6.1\ \mathrm{W}$ out of the $7.3\ \mathrm{W}$ lost) — the well-known weakness of linear regulators.\par
\emph{Regulation under a $-10\ \%$ mains sag:} $V_p = 207\ \mathrm{V}$ gives $V_s = 8.1\ \mathrm{V}$ (rms), so
\[
V_{\text{peak,bridge}} = 8.1\sqrt{2} - 1.4 \approx 10.1\ \mathrm{V}, \qquad V_{C,\min} = 10.1 - 0.5 = 9.6\ \mathrm{V} > 7.5\ \mathrm{V}.
\]
The dropout condition is still satisfied, so the output stays at $5.0\ \mathrm{V}$: the supply is correctly regulated even on a weak mains line — a valuable feature on the Cameroonian network.

\textbf{6. Report (summary of expected content).}\par
\begin{itemize}
\item Circuit diagram: mains $\rightarrow$ fuse $\rightarrow$ transformer $230/9\ \mathrm{V}$ $\rightarrow$ bridge (4 $\times$ 1N4007 or KBP206) $\rightarrow$ $C = 22\,000\ \mu\mathrm{F}/25\ \mathrm{V}$ $\rightarrow$ 7805 on heat sink (with $100\ \mathrm{nF}$ decoupling capacitors at input and output) $\rightarrow$ $5\ \mathrm{V}/1\ \mathrm{A}$ output with LED indicator.
\item Component table with ratings justified by the calculations above:
\end{itemize}
\begin{center}
\begin{tabularx}{\linewidth}{|l|X|X|}
\hline
\rowcolor{popblueL} \textbf{Component} & \textbf{Chosen value / type} & \textbf{Justification} \\
\hline
Transformer & $230/9\ \mathrm{V}$, $12\ \mathrm{VA}$ & ratio $25.6$; power $\geq 10\ \mathrm{VA}$ \\
\hline
Bridge & KBP206 ($2\ \mathrm{A}$, $600\ \mathrm{V}$) & PIV needed $\approx 12.7\ \mathrm{V}$; $I_F > 1\ \mathrm{A}$ with margin \\
\hline
Capacitor & $22\,000\ \mu\mathrm{F} / 25\ \mathrm{V}$ & $C \geq 20\,000\ \mu\mathrm{F}$; voltage rating $> 11.3\ \mathrm{V}$ \\
\hline
Regulator & 7805 + heat sink & $P_{7805} \approx 6.1\ \mathrm{W} \gg 1\ \mathrm{W}$ \\
\hline
\end{tabularx}
\end{center}
\begin{itemize}
\item Improvements suggested: mains fuse for safety, LED power indicator, larger heat sink, or a switch-mode (SMPS) design to raise the efficiency above $80\ \%$.
\end{itemize}

\begin{attention}[Common mistakes to avoid]
\begin{itemize}
\item Forgetting that $9\ \mathrm{V}$ is an \emph{rms} value: the capacitor charges to the \emph{peak} value $\hat{V}_s = V_s\sqrt{2}$.
\item Forgetting the two diode drops of the bridge ($2 \times 0.7\ \mathrm{V}$).
\item Using $\Delta t = T = 20\ \mathrm{ms}$ instead of $T/2 = 10\ \mathrm{ms}$ for a \emph{full-wave} rectifier (ripple frequency is $100\ \mathrm{Hz}$, not $50\ \mathrm{Hz}$).
\item Comparing the PIV with the rms voltage instead of the peak voltage.
\item Neglecting the regulator's dropout voltage: the \emph{minimum} of the ripple, not its average, must stay above $7.5\ \mathrm{V}$.
\item Mishandling powers of ten in $C = I\,\Delta t/\Delta V$: always convert the result into farads first, then into microfarads.
\end{itemize}
\end{attention}

\begin{savaistu}[Why do real phone chargers look different?]
Modern USB chargers use \emph{switch-mode power supplies} (SMPS): the mains is rectified first, then chopped at high frequency through a very small transformer. Because the transformer operates at tens of kilohertz instead of $50\ \mathrm{Hz}$, it can be tiny and light, and efficiencies above $80\ \%$ are common. The linear supply you have just designed is heavier and less efficient, but it is robust, cheap, easy to repair — and an excellent training ground for understanding every stage of energy conversion, from the ENEO socket to the $5\ \mathrm{V}$ USB port.
\end{savaistu}

\newpage

\coursec{Methods and worked examples}

\coursec{Alternating Current}

\courssub{Grandeurs efficaces et impédance en régime sinusoïdal}

\textbf{Skill 1 --- Computing and using r.m.s. values}\par

\begin{methode}[Handling r.m.s. values of sinusoidal quantities]
\begin{enumerate}
\item \textbf{Identify what is given.} In a.c. problems a voltage or current may be quoted as a \emph{peak} value $V_0$ (or $I_0$), a \emph{peak-to-peak} value $V_{pp}$, or an \emph{r.m.s.} value $V$ (or $I$). Underline which one it is.
\item \textbf{Convert with the sinusoidal relations} (valid \emph{only} for sine waves):
\[
V = \frac{V_0}{\sqrt{2}}, \qquad I = \frac{I_0}{\sqrt{2}}, \qquad V_{pp} = 2V_0.
\]
\item \textbf{Use r.m.s. values in all power and heating formulas}: $P = VI = RI^2 = V^2/R$. The r.m.s. value is precisely the value that makes the d.c. formulas work for a.c.
\item \textbf{Remember the measuring instruments}: a.c. voltmeters and ammeters (moving-iron or rectifier types) read \cle{r.m.s. values}; a cathode-ray oscilloscope displays the waveform, so you read \emph{peak} or \emph{peak-to-peak} values from the screen.
\item \textbf{For non-sinusoidal periodic signals}, go back to the definition:
\[
I_{\mathrm{rms}} = \sqrt{\langle i^2 \rangle} = \sqrt{\frac{1}{T}\int_0^T i^2\,dt},
\]
i.e. square the signal, take the mean over one period, then the square root.
\end{enumerate}
\end{methode}

\begin{exoresolu}[Mains supply and a heating element]
The domestic mains in Cameroon supplies a sinusoidal voltage of r.m.s. value $V = \SI{220}{V}$ at frequency $f = \SI{50}{Hz}$. An electric kettle used in a Yaound\'e household has a heating element of resistance $R = \SI{48.4}{\Omega}$.
\begin{enumerate}
\item Write the expression of the instantaneous voltage $v(t)$, giving its numerical peak value.
\item Calculate the r.m.s. current and the peak current through the element.
\item Calculate the mean power dissipated by the kettle.
\item The same kettle is connected to a \SI{220}{V} \emph{d.c.} battery. Compare the heating effects.
\end{enumerate}
\tcblower
\textbf{1. Instantaneous voltage.} The peak value is
\[
V_0 = V\sqrt{2} = 220 \times 1.414 \approx \SI{311}{V}.
\]
The angular frequency is $\omega = 2\pi f = 2\pi \times 50 = \SI{314}{rad.s^{-1}}$. Hence
\[
v(t) = 311\sin(314\,t)\ \text{(volts, } t \text{ in seconds)}.
\]

\textbf{2. Currents.} For a pure resistor, Ohm's law applies at every instant, so
\[
I = \frac{V}{R} = \frac{220}{48.4} \approx \SI{4.55}{A}, \qquad
I_0 = I\sqrt{2} = 4.55 \times 1.414 \approx \SI{6.43}{A}.
\]

\textbf{3. Mean power.} Using r.m.s. values:
\[
P = VI = 220 \times 4.55 \approx \SI{1.00e3}{W} = \SI{1.0}{kW}.
\]
(Check: $P = RI^2 = 48.4 \times 4.55^2 \approx \SI{1000}{W}$. \checkmark)

\textbf{4. Comparison with d.c.} By the very \emph{definition} of the r.m.s. value, a sinusoidal voltage of r.m.s. value \SI{220}{V} produces in a resistor \textbf{exactly the same heating} as a steady \SI{220}{V} d.c. source. The kettle behaves identically: $P = V^2/R = 220^2/48.4 = \SI{1.0}{kW}$. This is why r.m.s. values are so useful: they let us treat a.c. heating problems with ordinary d.c. formulas.
\end{exoresolu}

\textbf{Skill 2 --- Computing the impedance of R, L, C and series RLC circuits}\par

\begin{methode}[Impedance calculations in sinusoidal steady state]
\begin{enumerate}
\item \textbf{Write the angular frequency first}: $\omega = 2\pi f$. Most errors come from mixing $f$ and $\omega$.
\item \textbf{Individual impedances} (in ohms):
\[
Z_R = R, \qquad Z_L = L\omega, \qquad Z_C = \frac{1}{C\omega}.
\]
\item \textbf{Series RLC combination}: resistance and reactances combine as
\[
Z = \sqrt{R^2 + \left(L\omega - \frac{1}{C\omega}\right)^2}.
\]
The quantity $X = L\omega - \frac{1}{C\omega}$ is the \cle{reactance}.
\item \textbf{Current}: $I = V/Z$ (r.m.s. values). The phase angle $\varphi$ of $v$ relative to $i$ satisfies $\tan\varphi = X/R$, $\cos\varphi = R/Z$.
\item \textbf{Check the limiting behaviour}: at high frequency $Z_L \to \infty$, $Z_C \to 0$; at low frequency the reverse. Use this to sanity-check your answer.
\end{enumerate}
\end{methode}

\begin{popfigure}
\begin{tikzpicture}[>=stealth, scale=0.9]
% Phasor diagram
\draw[->, thick, popblue] (0,0) -- (3,0) node[right] {$V_R$ (in phase with $I$)};
\draw[->, thick, poporange] (3,0) -- (3,2.5) node[above right] {$V_L$ (leads $I$ by $\pi/2$)};
\draw[->, thick, popgreen] (3,0) -- (3,-3.5) node[below right] {$V_C$ (lags $I$ by $\pi/2$)};
\draw[->, very thick, popdark] (0,0) -- (3,-1) node[below right] {$V = \sqrt{V_R^2 + (V_L-V_C)^2}$};
\draw[dashed, thin] (3,-1) -- (3,0);
\draw[<->] (3.2, -1) -- (3.2, 0) node[midway, right] {$V_L-V_C$};
\end{tikzpicture}
\legende{Phasor diagram for a series RLC circuit. The supply voltage $V$ is the vector sum of $V_R$, $V_L$ and $V_C$.}
\end{popfigure}

\begin{exoresolu}[Series RLC circuit on the mains]
A coil of resistance $R = \SI{20}{\Omega}$ and inductance $L = \SI{0.10}{H}$ is connected in series with a capacitor $C = \SI{50}{\micro F}$ across a $\SI{220}{V}$, $\SI{50}{Hz}$ a.c. supply.
\begin{enumerate}
\item Calculate the impedance of the circuit.
\item Calculate the r.m.s. current.
\item Calculate the r.m.s. voltage across each component.
\item Verify that the supply voltage equals the phasor sum of the component voltages, and explain why $V \neq V_R + V_L + V_C$ arithmetically.
\end{enumerate}
\tcblower
\textbf{Data conversion}: $\omega = 2\pi f = 314\ \mathrm{rad\,s^{-1}}$.

\textbf{1. Impedance.}
\[
Z_L = L\omega = 0.10 \times 314 = \SI{31.4}{\Omega}, \qquad
Z_C = \frac{1}{C\omega} = \frac{1}{50\times10^{-6}\times 314} = \SI{63.7}{\Omega}.
\]
\[
Z = \sqrt{R^2 + (Z_L - Z_C)^2} = \sqrt{20^2 + (31.4 - 63.7)^2} = \sqrt{400 + 1043} \approx \SI{38.0}{\Omega}.
\]

\textbf{2. Current.}
\[
I = \frac{V}{Z} = \frac{220}{38.0} \approx \SI{5.79}{A}.
\]

\textbf{3. Component voltages.}
\[
V_R = RI = 20 \times 5.79 = \SI{116}{V}, \quad
V_L = Z_L I = 31.4 \times 5.79 = \SI{182}{V}, \quad
V_C = Z_C I = 63.7 \times 5.79 = \SI{369}{V}.
\]

\textbf{4. Phasor check.} The voltages $V_L$ and $V_C$ are in \emph{opposite phases} (they differ by $\pi$), and $V_R$ is in quadrature with both. Hence
\[
V = \sqrt{V_R^2 + (V_L - V_C)^2} = \sqrt{116^2 + (182-369)^2} = \sqrt{13456 + 34969} \approx \SI{220}{V}. \ \checkmark
\]
The arithmetic sum $116 + 182 + 369 = \SI{667}{V}$ is meaningless because the three voltages do not reach their maxima at the same instant: instantaneous voltages add, r.m.s. magnitudes do not. Note also that $V_C = \SI{369}{V}$ \emph{exceeds} the supply voltage --- a classic feature of RLC circuits near resonance, and a favourite GCE trap.
\end{exoresolu}

\courssub{Résonance et puissance dans les circuits AC}

\textbf{Skill 3 --- Resonance in a series RLC circuit}\par

\begin{methode}[Locating and exploiting resonance]
\begin{enumerate}
\item \textbf{Condition for resonance}: the reactance vanishes,
\[
L\omega_0 = \frac{1}{C\omega_0} \quad\Longrightarrow\quad f_0 = \frac{1}{2\pi\sqrt{LC}}.
\]
\item \textbf{Consequences to quote}: at $f = f_0$, the impedance is \emph{minimum} and purely resistive ($Z = R$); the current is \emph{maximum}, $I_{\max} = V/R$; current and voltage are \emph{in phase} ($\cos\varphi = 1$).
\item \textbf{Sharpness}: the smaller $R$ (for given $L$, $C$), the sharper the resonance peak and the larger the voltage magnification across $L$ or $C$ (quality factor $Q = \dfrac{L\omega_0}{R} = \dfrac{1}{RC\omega_0}$).
\item \textbf{Experimental identification}: vary $f$, keep $V$ constant, record $I(f)$; the peak of the response curve gives $f_0$. Alternatively observe $v$ and $i$ on a double-beam CRO: at resonance the two traces are in phase.
\item \textbf{Exam reflex}: if asked to \emph{find} an unknown $L$ or $C$, measure $f_0$ and invert the formula: $LC = \dfrac{1}{4\pi^2 f_0^2}$.
\end{enumerate}
\end{methode}

\begin{popfigure}
\begin{tikzpicture}
\begin{axis}[
    width=\linewidth, height=5cm,
    xlabel={Frequency $f$ (kHz)}, ylabel={Current $I$ (mA)},
    xmin=700, xmax=900, ymin=0, ymax=0.12,
    axis lines=middle, tick label style={font=\small},
    legend style={at={(0.98,0.98)}, anchor=north east, font=\small}
]
\addplot[domain=700:900, samples=200, thick, popblue] {0.1/(sqrt(max(0,100 + (2*pi*x*1e3*2e-3 - 1/(2*pi*x*1e3*1.98e-11))^2/100)))};
\addplot[dashed, thick, poporange] coordinates {(800,0) (800,0.11)} node[above, pos=0.9, black] {$f_0$};
\legend{$I(f)$, Resonance}
\end{axis}
\end{tikzpicture}
\legende{Resonance curve for the radio tuning circuit ($R=10\,\Omega$, $L=2\,\mathrm{mH}$, $C\approx20\,\mathrm{pF}$). The sharp peak shows high selectivity ($Q\approx1000$).}
\end{popfigure}

\begin{exoresolu}[Tuning a radio receiver]
A radio tuning circuit consists of a coil of inductance $L = \SI{2.0}{mH}$ and resistance $R = \SI{10}{\Omega}$ in series with a variable capacitor. The circuit is driven by the signal picked up by the antenna, modelled as a source of constant r.m.s. e.m.f. $V = \SI{1.0}{mV}$ at variable frequency.
\begin{enumerate}
\item Calculate the capacitance needed to tune the circuit to a station broadcasting at $f_0 = \SI{800}{kHz}$.
\item Calculate the r.m.s. current at resonance.
\item Calculate the r.m.s. voltage across the capacitor at resonance and comment.
\item The capacitance is now set to twice the value found in (1). Without full calculation, state whether the resonant frequency increases or decreases, and by what factor.
\end{enumerate}
\tcblower
\textbf{1. Capacitance.} From $f_0 = \dfrac{1}{2\pi\sqrt{LC}}$:
\[
C = \frac{1}{4\pi^2 f_0^2 L} = \frac{1}{4\pi^2 \times (8.0\times10^{5})^2 \times 2.0\times10^{-3}}
\approx \SI{1.98e-11}{F} \approx \SI{20}{pF}.
\]

\textbf{2. Current at resonance.} At resonance $Z = R$:
\[
I_{\max} = \frac{V}{R} = \frac{1.0\times10^{-3}}{10} = \SI{1.0e-4}{A} = \SI{0.10}{mA}.
\]

\textbf{3. Capacitor voltage.} First, $\omega_0 = 2\pi f_0 = 5.03\times10^{6}\ \mathrm{rad\,s^{-1}}$.
\[
Z_C = \frac{1}{C\omega_0} = \frac{1}{1.98\times10^{-11}\times 5.03\times10^{6}} \approx \SI{1.0e4}{\Omega}.
\]
\[
V_C = Z_C I_{\max} = 1.0\times10^{4} \times 1.0\times10^{-4} = \SI{1.0}{V}.
\]
The voltage across the capacitor is \SI{1.0}{V}, i.e. \textbf{1000 times larger} than the supply e.m.f. of \SI{1.0}{mV}! This \emph{voltage magnification} ($Q = L\omega_0/R \approx 1000$) is exactly how a tiny antenna signal is turned into a usable voltage: resonance amplifies selectively the chosen frequency while rejecting others.

\textbf{4. Doubling C.} Since $f_0 \propto 1/\sqrt{C}$, doubling $C$ \emph{decreases} $f_0$ by a factor $\sqrt{2}$: $f_0' = 800/\sqrt{2} \approx \SI{566}{kHz}$.
\end{exoresolu}

\textbf{Skill 4 --- Power in a.c. circuits and the power factor}\par

\begin{methode}[Mean power and power factor]
\begin{enumerate}
\item \textbf{Only resistance dissipates.} A pure inductor or capacitor absorbs energy during one quarter-cycle and returns it during the next: its mean power is \emph{zero}.
\item \textbf{Master formula}:
\[
P = VI\cos\varphi = RI^2,
\]
where $\cos\varphi = R/Z$ is the \cle{power factor}. Both forms must give the same answer --- use this as a check.
\item \textbf{Apparent power} $S = VI$ (in VA) is what the supply must deliver; the \emph{useful} power is $P = S\cos\varphi$. A low power factor means large currents (and large $RI^2$ line losses) for little useful power.
\item \textbf{Improving the power factor} of an inductive load (motor, fluorescent lamp): connect a capacitor in parallel so that the reactive currents of $L$ and $C$ cancel locally.
\item \textbf{Exam reflex}: never write $P = VI$ for a circuit containing $L$ or $C$ without the factor $\cos\varphi$.
\end{enumerate}
\end{methode}

\begin{exoresolu}[Workshop motor on the mains]
A small workshop in Douala runs a single-phase motor which, connected to the $\SI{220}{V}$, $\SI{50}{Hz}$ mains, draws an r.m.s. current of $\SI{8.0}{A}$. The motor is modelled as a resistance $R = \SI{15}{\Omega}$ in series with an inductance $L$.
\begin{enumerate}
\item Calculate the impedance of the motor and deduce $L$.
\item Calculate the power factor of the motor.
\item Calculate the mean power consumed by the motor in two different ways.
\item Calculate the energy consumed (in kWh) if the motor runs 6 hours per day for 26 days, and the cost at 65 FCFA per kWh.
\end{enumerate}
\tcblower
\textbf{1. Impedance and inductance.}
\[
Z = \frac{V}{I} = \frac{220}{8.0} = \SI{27.5}{\Omega}.
\]
\[
L\omega = \sqrt{Z^2 - R^2} = \sqrt{27.5^2 - 15^2} = \sqrt{756.25 - 225} = \sqrt{531.25} \approx \SI{23.0}{\Omega},
\]
\[
L = \frac{23.0}{2\pi \times 50} \approx \SI{7.3e-2}{H} = \SI{73}{mH}.
\]

\textbf{2. Power factor.}
\[
\cos\varphi = \frac{R}{Z} = \frac{15}{27.5} \approx 0.55.
\]

\textbf{3. Mean power.} Method A:
\[
P = VI\cos\varphi = 220 \times 8.0 \times 0.55 \approx \SI{960}{W}.
\]
Method B:
\[
P = RI^2 = 15 \times 8.0^2 = \SI{960}{W}. \ \checkmark
\]
The two methods agree, as they must. Note that the apparent power $S = VI = \SI{1760}{VA}$ is nearly twice the useful power --- the supply and cables must carry \SI{8.0}{A} although only \SI{960}{W} is useful.

\textbf{4. Energy and cost.}
\[
E = Pt = 0.960\ \mathrm{kW} \times (6 \times 26)\ \mathrm{h} = 0.960 \times 156 \approx \SI{150}{kWh}.
\]
\[
\text{Cost} = 150 \times 65 \approx \SI{9740}{FCFA} \approx \SI{9.7e3}{FCFA} \text{ per month.}
\]
A better power factor (capacitor compensation) would reduce the line current and the heating of the installation, even though the energy billed to the workshop stays the same.
\end{exoresolu}

\courssub{Transformateurs : modèle idéal et pertes réelles}

\textbf{Skill 5 --- The ideal transformer}\par

\begin{popfigure}
\begin{tikzpicture}[>=stealth, scale=1.2]
% Core
\draw[thick, fill=popblueL!30] (0,0) rectangle (1,3);
\draw[thick, fill=popblueL!30] (3,0) rectangle (4,3);
\draw[thick, fill=popblueL!30] (0,3) -- (4,3) -- (4,3.5) -- (0,3.5) -- cycle;
\draw[thick, fill=popblueL!30] (0,-0.5) -- (4,-0.5) -- (4,0) -- (0,0) -- cycle;
% Windings
\foreach \y in {0.5, 1.0, 1.5, 2.0, 2.5} {
    \draw[popblue, thick] (-0.3,\y) arc (180:0:0.5) -- (1.3,\y);
    \draw[poporange, thick] (3.7,\y) arc (0:-180:0.5) -- (2.7,\y);
}
% Labels
\node[left] at (-0.5,1.75) {Primary $N_p$};
\node[right] at (4.5,1.75) {Secondary $N_s$};
\node[above] at (2,3.7) {Laminated soft-iron core};
\draw[->, thick] (-1,1.75) -- (-0.3,1.75) node[midway, above] {$V_p, I_p$};
\draw[->, thick] (4.3,1.75) -- (5,1.75) node[midway, above] {$V_s, I_s$};
\end{tikzpicture}
\legende{Schematic of an ideal transformer: primary and secondary windings on a common closed core.}
\end{popfigure}

\begin{methode}[Applying the ideal transformer relations]
\begin{enumerate}
\item \textbf{Identify primary and secondary}: the primary is connected to the source; the secondary feeds the load. Count or note the turns $N_p$, $N_s$.
\item \textbf{Write the ideal relations} (no derivation required at GCE, but quote them correctly):
\[
\frac{V_p}{V_s} = \frac{N_p}{N_s} = \frac{I_s}{I_p}.
\]
All voltages and currents are r.m.s. values.
\item \textbf{Step-up or step-down?} $N_s > N_p$: step-up (voltage rises, current falls). $N_s < N_p$: step-down.
\item \textbf{Power conservation}: for the ideal transformer $V_p I_p = V_s I_s$ (100\% efficiency). Use this as a cross-check.
\item \textbf{Chain problems} (power transmission): treat each transformer and the transmission line separately; the line has resistance, so apply $P_{\text{loss}} = R_{\text{line}} I_{\text{line}}^2$ \emph{only} to the line section.
\end{enumerate}
\end{methode}

\begin{exoresolu}[Why transmit at high voltage?]
A small hydroelectric plant delivers $\SI{100}{kW}$ to a village through a transmission line of total resistance $\SI{4.0}{\Omega}$.
\begin{enumerate}
\item The power is first transmitted at an r.m.s. voltage of $\SI{5.0}{kV}$. Calculate the line current and the power lost as heat in the line.
\item A step-up transformer raises the transmission voltage to $\SI{50}{kV}$. Its primary has $N_p = 400$ turns; calculate $N_s$.
\item Recalculate the line current, the line losses, and the percentage of the generated power that is lost in each case. Comment.
\item At the village, a step-down transformer reduces the voltage to $\SI{220}{V}$. Assuming it is ideal and delivers the full $\SI{100}{kW}$ (neglect line losses), calculate the current available to the villagers.
\end{enumerate}
\tcblower
\textbf{1. Transmission at \SI{5.0}{kV}.}
\[
I = \frac{P}{V} = \frac{1.0\times10^{5}}{5.0\times10^{3}} = \SI{20}{A}, \qquad
P_{\text{loss}} = RI^2 = 4.0 \times 20^2 = \SI{1.6e3}{W} = \SI{1.6}{kW}.
\]

\textbf{2. Step-up transformer.}
\[
N_s = N_p \frac{V_s}{V_p} = 400 \times \frac{50\,000}{5\,000} = 4000\ \text{turns}.
\]

\textbf{3. Transmission at \SI{50}{kV}.}
\[
I' = \frac{1.0\times10^{5}}{5.0\times10^{4}} = \SI{2.0}{A}, \qquad
P'_{\text{loss}} = 4.0 \times 2.0^2 = \SI{16}{W}.
\]
Percentages lost:
\[
\text{at } 5\ \mathrm{kV}:\ \frac{1.6}{100} = 1.6\,\%; \qquad
\text{at } 50\ \mathrm{kV}:\ \frac{0.016}{100} = 0.016\,\%.
\]
Multiplying the voltage by 10 divides the current by 10 and the $RI^2$ losses by \textbf{100}. This is the fundamental reason why ENEO transmits energy at high voltage over long distances and steps it down near consumers.

\textbf{4. Village supply.}
\[
I_{\text{village}} = \frac{P}{V} = \frac{1.0\times10^{5}}{220} \approx \SI{455}{A}.
\]
The step-down transformer delivers a large current at low voltage --- exactly what domestic and workshop loads require.
\end{exoresolu}

\textbf{Skill 6 --- Losses in a real transformer and efficiency}\par

\begin{methode}[Identifying and minimising transformer losses]
\begin{enumerate}
\item \textbf{List the four classical sources of loss} and the remedy for each:
\begin{itemize}
\item \emph{Joule (copper) losses} $RI^2$ in the windings $\rightarrow$ use thick low-resistivity copper wire.
\item \emph{Eddy currents} induced in the core $\rightarrow$ build the core from thin \textbf{laminations} insulated from one another.
\item \emph{Hysteresis} (repeated magnetisation of the core) $\rightarrow$ use a \textbf{soft magnetic material} (soft iron, silicon steel) with a narrow hysteresis loop.
\item \emph{Flux leakage} $\rightarrow$ wind the coils on a closed core, one over the other, so nearly all the flux links both windings.
\end{itemize}
\item \textbf{Efficiency}:
\[
\eta = \frac{P_{\text{out}}}{P_{\text{in}}} = \frac{V_s I_s \cos\varphi_s}{V_p I_p \cos\varphi_p}, \qquad
P_{\text{loss}} = P_{\text{in}} - P_{\text{out}}.
\]
\item \textbf{Exam reflex}: if a transformer ``heats up'', the lost power reappears as internal energy; always express $\eta$ as a percentage and check $\eta < 100\,\%$.
\end{enumerate}
\end{methode}

\begin{exoresolu}[Efficiency of a distribution transformer]
A step-down transformer supplies a neighbourhood. Its primary is connected to a $\SI{11}{kV}$ line and draws an r.m.s. current of $\SI{5.0}{A}$ with a power factor of $0.95$. The secondary delivers $\SI{220}{V}$ to consumers drawing a total r.m.s. current of $\SI{230}{A}$ with power factor $0.90$.
\begin{enumerate}
\item Calculate the input and output powers.
\item Deduce the efficiency of the transformer and the power dissipated inside it.
\item State two physical origins of this lost power and how the construction of the transformer reduces each.
\end{enumerate}
\tcblower
\textbf{1. Powers.}
\[
P_{\text{in}} = V_p I_p \cos\varphi_p = 11\,000 \times 5.0 \times 0.95 = \SI{5.225e4}{W} \approx \SI{52.3}{kW}.
\]
\[
P_{\text{out}} = V_s I_s \cos\varphi_s = 220 \times 230 \times 0.90 = \SI{4.554e4}{W} \approx \SI{45.5}{kW}.
\]

\textbf{2. Efficiency and losses.}
\[
\eta = \frac{P_{\text{out}}}{P_{\text{in}}} = \frac{45.5}{52.3} \approx 0.871 \approx 87\,\%.
\]
\[
P_{\text{loss}} = 52.3 - 45.5 = \SI{6.7}{kW}.
\]
Nearly \SI{7}{kW} is wasted as heat, which is why large transformers need cooling (oil baths, radiating fins).

\textbf{3. Origins and remedies.} (Any two of the following.)
\begin{itemize}
\item \textbf{Eddy currents} in the core: the alternating flux induces circulating currents in the iron, which dissipate $RI^2$ heat. Remedy: the core is made of thin \emph{laminated} sheets insulated by varnish, forcing the currents into tiny loops of high resistance.
\item \textbf{Hysteresis}: energy is lost each cycle in repeatedly reversing the magnetisation of the core. Remedy: use \emph{soft} iron (or silicon steel) whose narrow hysteresis loop encloses little area, hence little energy loss per cycle.
\item \textbf{Copper losses}: the windings have resistance. Remedy: thick copper conductors.
\end{itemize}
\end{exoresolu}

\courssub{Redressement et lissage des signaux alternatifs}

\textbf{Skill 7 --- Rectification and smoothing of a.c. signals}\par

\begin{popfigure}
\begin{tikzpicture}[european, scale=1.1]
% Bridge rectifier with smoothing capacitor
\draw (0,0) to[sV=$v_s(t)$] (0,-3);
\draw (0,0) -- (1,0);
\draw (0,-3) -- (1,-3);
% Diodes
\draw (1,0) to[D=$D_1$] (3,0);
\draw (1,-3) to[D=$D_3$] (3,0);
\draw (3,-3) to[D=$D_2$] (1,-3);
\draw (1,0) to[D=$D_4$] (3,-3);
% Load and Capacitor
\draw (3,0) to[R=$R_L$, *-*] (3,-1.5);
\draw (3,-1.5) to[C=$C$, *-*] (3,-3);
% Output labels
\node[right] at (3.3,-0.75) {$+$};
\node[right] at (3.3,-2.25) {$-$};
\node[above] at (3,0) {$v_{out}$};
\end{tikzpicture}
\legende{Bridge rectifier (four diodes) with smoothing capacitor $C$ in parallel with load $R_L$. During each half-cycle, a diagonal pair of diodes conducts, giving unidirectional current through $R_L$.}
\end{popfigure}

\begin{methode}[Analysing rectifier circuits]
\begin{enumerate}
\item \textbf{Identify the configuration}:
\begin{itemize}
\item \emph{Half-wave}: one diode; only one half-cycle passes. Output frequency = supply frequency.
\item \emph{Full-wave}: centre-tapped transformer with two diodes, or a \textbf{bridge} of four diodes; both half-cycles pass, inverted when negative. Output ripple frequency = $2f$.
\end{itemize}
\item \textbf{Diode conduction rule}: an ideal diode conducts when forward biased ($V_A > V_K$) and blocks otherwise. Trace the current path for each half-cycle.
\item \textbf{Sketch the output}: half-wave gives separated humps; full-wave gives $|\sin|$ humps touching zero.
\item \textbf{Smoothing}: connect a large capacitor in \emph{parallel} with the load. It charges to the peak $V_0$ and discharges slowly through the load $R$ between peaks:
\[
\Delta V \approx \frac{I}{f_r C} \quad \text{with } f_r = f \text{ (half-wave) or } 2f \text{ (full-wave)}.
\]
Large $C$, large $R$ (small $I$) and full-wave rectification all reduce the ripple.
\item \textbf{Exam reflex}: state clearly that the smoothed output is \emph{unidirectional but not perfectly steady}; a regulator is needed for a truly constant voltage.
\end{enumerate}
\end{methode}

\begin{exoresolu}[Phone-charger style power supply]
A bridge rectifier (four ideal diodes) is fed from the $\SI{50}{Hz}$ mains through a step-down transformer whose secondary delivers a sinusoidal voltage of peak value $V_0 = \SI{12}{V}$. The rectified output feeds a load $R = \SI{24}{\Omega}$, with a smoothing capacitor $C = \SI{2200}{\micro F}$ connected in parallel with the load.
\begin{enumerate}
\item Describe, with the help of the diode conduction rule, how the bridge produces a full-wave rectified voltage. State the ripple frequency.
\item Sketch (describe precisely) the voltage across $R$ without the capacitor.
\item With the capacitor connected, estimate the peak-to-peak ripple $\Delta V$ when the load current is about $\SI{0.5}{A}$.
\item The capacitor is removed and replaced by one of $\SI{470}{\micro F}$. What happens to the ripple? Give the new value.
\item State one advantage of full-wave over half-wave rectification for smoothing.
\end{enumerate}
\tcblower
\textbf{1. Bridge operation.} During the positive half-cycle, two diagonally opposite diodes ($D_1$ and $D_2$) are forward biased and conduct, sending current through $R$ in one direction; the other two diodes ($D_3$, $D_4$) are reverse biased and block. During the negative half-cycle the roles interchange: the \emph{other} diagonal pair ($D_3$, $D_4$) conducts, but the current through $R$ keeps the \emph{same} direction. The output is therefore full-wave rectified, and the ripple frequency is
\[
f_r = 2f = \SI{100}{Hz}.
\]

\textbf{2. Output without capacitor.} The voltage across $R$ is $v_R(t) = |12\sin(314t)|$: a succession of identical positive half-sine humps of amplitude \SI{12}{V}, repeating every $T/2 = \SI{10}{ms}$, falling to zero between humps. The current is unidirectional but strongly varying.

\textbf{3. Ripple with $C = \SI{2200}{\micro F}$.} Between peaks the capacitor discharges approximately at constant current $I \approx V_0/R = 12/24 = \SI{0.5}{A}$ for a time $\Delta t \approx 1/f_r = \SI{10}{ms}$:
\[
\Delta V \approx \frac{I}{f_r C} = \frac{0.5}{100 \times 2200\times10^{-6}} \approx \SI{2.3}{V}.
\]
The output now stays between about \SI{9.7}{V} and \SI{12}{V}: a much steadier d.c. with roughly $20\,\%$ ripple.

\textbf{4. Smaller capacitor.} With $C = \SI{470}{\micro F}$:
\[
\Delta V \approx \frac{0.5}{100 \times 470\times10^{-6}} \approx \SI{10.6}{V}.
\]
The ripple becomes enormous (the voltage nearly collapses between peaks): the capacitor is too small to ``hold'' the charge for \SI{10}{ms}. Ripple is inversely proportional to $C$.

\textbf{5. Advantage of full-wave.} The capacitor is recharged \emph{twice per cycle} (every \SI{10}{ms} instead of every \SI{20}{ms}), so for the same $C$ and load the discharge time --- and hence the ripple --- is halved compared with half-wave rectification. Full-wave rectification therefore gives a smoother output and a higher mean voltage.
\end{exoresolu}

\begin{attention}[Frequent traps in the Alternating Current chapter]
\begin{itemize}
\item Using peak values in power formulas: $P = V_0 I_0$ is \emph{wrong}; use r.m.s. values, $P = VI\cos\varphi$.
\item Applying $V_0 = V\sqrt{2}$ to a non-sinusoidal signal (square, triangular): the factor $\sqrt{2}$ is for sine waves only.
\item Adding r.m.s. voltages arithmetically in a series RLC circuit: they must be added as phasors, $V = \sqrt{V_R^2 + (V_L - V_C)^2}$.
\item Writing $f_0 = 1/(2\pi\sqrt{LC})$ with $L$ in mH or $C$ in $\mu$F without converting to SI units.
\item Inverting the transformer ratio: $V_p/V_s = N_p/N_s = I_s/I_p$ --- note that the \emph{current} ratio is the inverse of the voltage ratio.
\item Forgetting that an ideal transformer works only on a.c.: a steady d.c. primary produces no changing flux, hence no secondary e.m.f.
\item Believing the smoothed rectified voltage is perfectly constant: there is always a residual ripple $\Delta V \approx I/(f_r C)$.
\end{itemize}
\end{attention}

\newpage

\coursec{Exercises}

\courssub{I check my knowledge}

\begin{exercice}[MCQ: RMS value of a composite signal]
A voltage signal is given by $v(t)=10\sin(314t)+5\sin(942t)$ volts. What is its RMS value?
\begin{enumerate}
\item $7.5\ \mathrm{V}$
\item $\sqrt{10^2+5^2}\ \mathrm{V}$
\item $\sqrt{\frac{10^2}{2}+\frac{5^2}{2}}\ \mathrm{V}$
\item $15\ \mathrm{V}$
\end{enumerate}
\end{exercice}

\begin{exercice}[True/False with justification]
State whether each statement is true or false. Justify your answer briefly.
\begin{enumerate}
\item At resonance in a series RLC circuit, the impedance is purely resistive and equal to $R$.
\item The turns ratio of an ideal transformer equals the ratio of primary to secondary currents.
\item In a half-wave rectifier with a smoothing capacitor, the ripple frequency equals the supply frequency.
\item The power factor of a purely inductive load is zero (lagging).
\end{enumerate}
\end{exercice}

\begin{exercice}[Definitions and direct calculations]
\begin{enumerate}
\item Define the \cle{root-mean-square (RMS) value} of an alternating current.
\item Define \cle{impedance} and state its unit.
\item A coil has a resistance of $5\ \Omega$ and an inductance of $20\ \mathrm{mH}$. Calculate its impedance at $50\ \mathrm{Hz}$.
\item An ideal transformer has $800$ primary turns and $200$ secondary turns. If the primary voltage is $240\ \mathrm{V}$, what is the secondary voltage?
\end{enumerate}
\end{exercice}

\begin{exercice}[Quick numerical applications]
\begin{enumerate}
\item Calculate the resonance frequency of a series LC circuit with $L=50\ \mathrm{mH}$ and $C=10\ \mu\mathrm{F}$.
\item A $230\ \mathrm{V}$ (RMS), $50\ \mathrm{Hz}$ supply is connected to a $100\ \Omega$ resistor in series with a $0.1\ \mathrm{H}$ inductor. Find the RMS current.
\item A transformer is $90\%$ efficient. If the input power is $500\ \mathrm{W}$, what is the output power?
\item For a bridge rectifier supplied from a $12\ \mathrm{V}$ (RMS) secondary, calculate the peak inverse voltage (PIV) across each diode.
\end{enumerate}
\end{exercice}

\courssub{I practise}

\begin{exercice}[Series RLC circuit analysis]
A series RLC circuit has $R=10\ \Omega$, $L=50\ \mathrm{mH}$, $C=10\ \mu\mathrm{F}$ and is connected to a $230\ \mathrm{V}$ (RMS), $50\ \mathrm{Hz}$ supply.
Calculate:
\begin{enumerate}
\item The impedance $Z$ of the circuit.
\item The RMS current $I$.
\item The active power $P$, reactive power $Q$, and apparent power $S$.
\item The power factor and state whether it is leading or lagging.
\end{enumerate}
\end{exercice}

\begin{exercice}[Parallel RLC circuit and resonance]
A parallel RLC circuit consists of a resistor $R=100\ \Omega$, an inductor $L=20\ \mathrm{mH}$, and a capacitor $C=5\ \mu\mathrm{F}$ connected in parallel across an AC source.
\begin{enumerate}
\item Determine the resonance frequency $f_0$.
\item Calculate the quality factor $Q$ at resonance.
\item Sketch the curve of the magnitude of admittance $|Y|$ versus frequency $f$, indicating the resonance point and the half-power frequencies.
\item Explain why the impedance is maximum at resonance in a parallel RLC circuit.
\end{enumerate}
\end{exercice}

\begin{exercice}[Transformer calculations]
A single-phase transformer has $460$ turns on the primary and is designed to step down $230\ \mathrm{V}$ to $12\ \mathrm{V}$ (RMS). The secondary supplies a resistive load of $24\ \mathrm{W}$. The iron losses are $2\ \mathrm{W}$ and the copper losses at this load are $1.5\ \mathrm{W}$.
\begin{enumerate}
\item Calculate the number of turns on the secondary.
\item Find the secondary current (RMS) when the load is connected.
\item Determine the primary current (RMS) assuming an ideal transformer for the current ratio.
\item Calculate the efficiency of the transformer at this load.
\end{enumerate}
\end{exercice}

\begin{exercice}[Bridge rectifier with smoothing capacitor]
A bridge rectifier is fed from a transformer secondary rated at $12\ \mathrm{V}$ RMS, $50\ \mathrm{Hz}$. The load draws a constant current of $500\ \mathrm{mA}$. A smoothing capacitor is connected across the load.
\begin{enumerate}
\item Calculate the peak voltage across the secondary winding.
\item Determine the peak inverse voltage (PIV) rating required for each diode.
\item Find the minimum capacitance needed to keep the peak-to-peak ripple voltage below $1\ \mathrm{V}$.
\item Estimate the average DC output voltage across the load.
\end{enumerate}
\end{exercice}

\courssub{I prepare for the GCE}

\begin{exercice}[Rectification comparison – Essay question (15 marks)]
A transformer has a secondary winding rated at $12\ \mathrm{V}$ RMS, $50\ \mathrm{Hz}$. It is used to supply a $10\ \Omega$ resistive load via two different rectifier circuits:
\begin{enumerate}
\item A half-wave rectifier with a $1000\ \mu\mathrm{F}$ smoothing capacitor.
\item A bridge rectifier with the same $1000\ \mu\mathrm{F}$ smoothing capacitor.
\end{enumerate}
For each circuit:
\begin{enumerate}
\item Draw the circuit diagram showing the transformer secondary, rectifier(s), capacitor, and load. [3 marks]
\item Calculate the peak secondary voltage and the PIV across the diode(s). [3 marks]
\item Estimate the average DC output voltage and the peak-to-peak ripple voltage. [6 marks]
\end{enumerate}
Compare the two circuits in terms of output voltage, ripple, transformer utilisation, and diode PIV requirements. Justify why the bridge rectifier is preferred in modern power supplies. [3 marks]
\end{exercice}

\begin{exercice}[Power in AC circuits and power factor correction – Essay question (15 marks)]
A factory receives a $400\ \mathrm{V}$ (RMS), $50\ \mathrm{Hz}$ three-phase supply (you may treat it as single-phase equivalent for calculations). The total load consists of:
\begin{itemize}
\item $20\ \mathrm{kW}$ of resistive heating (unity power factor).
\item $15\ \mathrm{kW}$ of induction motors operating at $0.7$ lagging power factor.
\item $10\ \mathrm{kVAR}$ of capacitive reactive power from existing power factor correction capacitors.
\end{itemize}
\begin{enumerate}
\item Calculate the total active power $P$, total reactive power $Q$, and the overall apparent power $S$ of the factory. [4 marks]
\item Determine the overall power factor of the factory. State whether it is leading or lagging. [3 marks]
\item The supply authority imposes a penalty if the power factor falls below $0.95$ lagging. Calculate the additional capacitive reactive power (in kVAR) required to raise the overall power factor to exactly $0.95$ lagging. [4 marks]
\item Draw the power triangle before and after correction, labelling all sides and angles. [2 marks]
\item Explain why power factor correction reduces the current drawn from the supply and the benefits for the consumer and the utility. [2 marks]
\end{enumerate}
\end{exercice}

\begin{exercice}[Transformer design and losses – Essay question (15 marks)]
A $5\ \mathrm{kVA}$, $230\ \mathrm{V}/115\ \mathrm{V}$, $50\ \mathrm{Hz}$ single-phase transformer has the following parameters:
\begin{itemize}
\item Primary resistance $R_1 = 0.5\ \Omega$, secondary resistance $R_2 = 0.125\ \Omega$.
\item Iron loss (core loss) $P_{\text{iron}} = 80\ \mathrm{W}$ (constant at rated voltage).
\item Leakage reactances referred to primary: $X_1 = 1.2\ \Omega$, $X_2' = 0.3\ \Omega$.
\end{itemize}
The transformer supplies a load drawing rated current at $0.8$ power factor lagging.
\begin{enumerate}
\item Calculate the equivalent resistance and reactance referred to the primary side. [3 marks]
\item Determine the percentage voltage regulation at this load. [4 marks]
\item Compute the copper losses at rated load. [2 marks]
\item Find the efficiency at rated load and $0.8$ power factor. [3 marks]
\item The transformer is now operated at $50\%$ of rated load with the same power factor. Estimate the new efficiency and explain why it differs from the full-load efficiency. [3 marks]
\end{enumerate}
\end{exercice}

\newpage

\coursec{Solutions to the exercises}

\coursec{Grandeurs efficaces et impédance en régime sinusoïdal}

\begin{definition}[Valeur efficace (RMS)]
La \cle{valeur efficace} (ou valeur quadratique moyenne, RMS \emph{Root Mean Square}) d'une grandeur alternative périodique $x(t)$ de période $T$ est la valeur d'une grandeur continue qui dissiperait la même puissance moyenne dans une charge résistive identique. Elle se note $X_{\mathrm{rms}}$ ou $X$ et se calcule par :
\[
X_{\mathrm{rms}} = \sqrt{\frac{1}{T}\int_0^T x^2(t)\,dt}.
\]
Pour un signal sinusoïdal pur $x(t) = X_m \sin(\omega t + \varphi)$, la valeur efficace est $X_{\mathrm{rms}} = \dfrac{X_m}{\sqrt{2}}$.
\end{definition}

\begin{propriete}[Valeur efficace d'une somme de sinusoïdes de fréquences différentes]
Si une tension (ou un courant) est la somme de plusieurs sinusoïdes de fréquences distinctes (non harmoniquement liées ou harmoniques d'ordre différent), la valeur efficace totale est la racine carrée de la somme des carrés des valeurs efficaces individuelles :
\[
V_{\mathrm{rms}} = \sqrt{V_{\mathrm{rms},1}^2 + V_{\mathrm{rms},2}^2 + \dots}
\]
Ceci découle de l'orthogonalité des fonctions sinusoïdales de fréquences différentes sur une période commune (ou sur un temps infini).
\end{propriete}

\begin{exemplebox}[Valeur efficace d'un signal composite]
Soit $v(t) = 10\sin(314t) + 5\sin(942t)\ \mathrm{V}$.
$\omega_1 = 314\ \mathrm{rad/s}\ (f_1=50\ \mathrm{Hz})$, $\omega_2 = 942\ \mathrm{rad/s}\ (f_2=150\ \mathrm{Hz})$.
$V_{\mathrm{rms},1} = 10/\sqrt{2}\ \mathrm{V}$, $V_{\mathrm{rms},1}^2 = 50\ \mathrm{V^2}$.
$V_{\mathrm{rms},2} = 5/\sqrt{2}\ \mathrm{V}$, $V_{\mathrm{rms},2}^2 = 12.5\ \mathrm{V^2}$.
$V_{\mathrm{rms}} = \sqrt{50 + 12.5} = \sqrt{62.5} \approx 7.91\ \mathrm{V}$.
\end{exemplebox}

\courssub{Impédance et réactance}

\begin{definition}[Impédance complexe]
En régime sinusoïdal forcé, l'\cle{impédance complexe} $Z$ d'un dipôle est le rapport du fasor tension $\vec{V}$ au fasor courant $\vec{I}$ :
\[
Z = \frac{\vec{V}}{\vec{I}} = R + jX \quad (\mathrm{en}\ \Omega).
\]
$R$ est la \cle{résistance} (partie réelle), $X$ la \cle{réactance} (partie imaginaire). Le module de l'impédance est $|Z| = V_{\mathrm{rms}}/I_{\mathrm{rms}}$.
\end{definition}

\begin{propriete}[Impédances des éléments de base]
\begin{itemize}
\item Résistance $R$ : $Z_R = R$, $X_R = 0$. Tension et courant en phase.
\item Inductance $L$ : $Z_L = j\omega L = jX_L$, $X_L = \omega L = 2\pi f L$. Le courant \emph{retarde} la tension de $90^\circ$.
\item Capacité $C$ : $Z_C = \dfrac{1}{j\omega C} = -jX_C$, $X_C = \dfrac{1}{\omega C} = \dfrac{1}{2\pi f C}$. Le courant \emph{avance} sur la tension de $90^\circ$.
\end{itemize}
\end{propriete}

\begin{propriete}[Association en série et en parallèle]
\begin{itemize}
\item \textbf{Série :} $Z_{\mathrm{eq}} = Z_1 + Z_2 + \dots$, $R_{\mathrm{eq}} = \sum R_k$, $X_{\mathrm{eq}} = \sum X_k$.
\item \textbf{Parallèle :} Il est plus commode d'utiliser l'\cle{admittance} $Y = 1/Z = G + jB$ (conductance $G$, susceptance $B$). $Y_{\mathrm{eq}} = Y_1 + Y_2 + \dots$.
\end{itemize}
\end{propriete}

\begin{exemplebox}[Bobine réelle en série]
Une bobine a $R = 5\ \Omega$, $L = 20\ \mathrm{mH}$, alimentée sous $f = 50\ \mathrm{Hz}$.
$X_L = 2\pi \times 50 \times 0.020 = 6.283\ \Omega$.
$Z = \sqrt{R^2 + X_L^2} = \sqrt{25 + 39.48} = 8.03\ \Omega$.
Angle de phase $\varphi = \arctan(X_L/R) = \arctan(1.257) = 51.5^\circ$ (courant retardant).
\end{exemplebox}

\begin{methode}[Analyse d'un circuit série RLC]
\begin{enumerate}
\item Calculer $X_L = 2\pi f L$ et $X_C = 1/(2\pi f C)$.
\item Déterminer la réactance nette $X = X_L - X_C$.
\item Calculer l'impédance $Z = \sqrt{R^2 + X^2}$ et l'angle $\varphi = \arctan(X/R)$.
\item Courant efficace $I = V_{\mathrm{rms}}/Z$.
\item Tensions aux bornes de chaque élément : $V_R = IR$, $V_L = IX_L$, $V_C = IX_C$.
\end{enumerate}
\end{methode}

\begin{exoresolu}[Circuit série RLC — Analyse complète]
Données : $R = 10\ \Omega$, $L = 50\ \mathrm{mH}$, $C = 10\ \mu\mathrm{F}$, $V_{\mathrm{rms}} = 230\ \mathrm{V}$, $f = 50\ \mathrm{Hz}$.
\begin{enumerate}
\item $X_L = 2\pi \times 50 \times 0.050 = 15.71\ \Omega$.
$X_C = \dfrac{1}{2\pi \times 50 \times 10^{-5}} = 318.3\ \Omega$.
$X = 15.71 - 318.3 = -302.6\ \Omega$ (circuit capacitif).
\item $Z = \sqrt{10^2 + (-302.6)^2} = 302.8\ \Omega$.
$\varphi = \arctan(-302.6/10) = -88.1^\circ$ (courant avance).
\item $I = 230 / 302.8 = 0.760\ \mathrm{A}$.
\item $V_R = 7.60\ \mathrm{V}$, $V_L = 11.9\ \mathrm{V}$, $V_C = 241\ \mathrm{V}$.
Note : $V_C > V_{\mathrm{source}}$ (phénomène de surtenseur en série).
\end{enumerate}
\end{exoresolu}

\begin{exercice}[Applications directes — Impédance et RMS]
\begin{enumerate}
\item Calculer la valeur efficace de $i(t) = 3\sin(100\pi t) + 4\sin(300\pi t)\ \mathrm{A}$.
\item Une bobine de $R=12\ \Omega$, $L=50\ \mathrm{mH}$ est branchée sur $230\ \mathrm{V}$, $50\ \mathrm{Hz}$. Déterminer $Z$, $I$, $\varphi$.
\item Un circuit série $R=20\ \Omega$, $L=100\ \mathrm{mH}$, $C=20\ \mu\mathrm{F}$ est alimenté sous $120\ \mathrm{V}$, $60\ \mathrm{Hz}$. Trouver $I$ et la nature du circuit.
\end{enumerate}
\end{exercice}

\begin{corrige}[Exercise 1 — Applications directes]
\begin{enumerate}
\item $I_{\mathrm{rms}} = \sqrt{(3/\sqrt{2})^2 + (4/\sqrt{2})^2} = \sqrt{4.5 + 8} = \sqrt{12.5} = 3.54\ \mathrm{A}$.
\item $X_L = 2\pi \times 50 \times 0.05 = 15.71\ \Omega$. $Z = \sqrt{12^2 + 15.71^2} = 19.77\ \Omega$. $I = 230/19.77 = 11.6\ \mathrm{A}$. $\varphi = \arctan(15.71/12) = 52.6^\circ$ (retard).
\item $X_L = 37.7\ \Omega$, $X_C = 132.6\ \Omega$. $X = -94.9\ \Omega$ (capacitif). $Z = \sqrt{20^2 + 94.9^2} = 97.0\ \Omega$. $I = 120/97.0 = 1.24\ \mathrm{A}$.
\end{enumerate}
\end{corrige}

\coursec{Résonance et puissance dans les circuits AC}

\courssub{Résonance série}

\begin{definition}[Résonance série]
Dans un circuit série RLC, la \cle{résonance} se produit lorsque la réactance nette est nulle : $X_L = X_C$. L'impédance est alors purement résistive et minimale : $Z_{\min} = R$. Le courant est maximal et en phase avec la tension.
La \cle{fréquence de résonance} (ou fréquence propre) est :
\[
f_0 = \frac{1}{2\pi\sqrt{LC}} \quad \text{ou} \quad \omega_0 = \frac{1}{\sqrt{LC}}.
\]
\end{definition}

\begin{propriete}[Facteur de qualité $Q$ en série]
Le \cle{facteur de qualité} caractérise la sélectivité du circuit :
\[
Q = \frac{\omega_0 L}{R} = \frac{1}{\omega_0 C R} = \frac{1}{R}\sqrt{\frac{L}{C}}.
\]
Aux bornes de $L$ ou $C$, la tension est amplifiée : $V_L = V_C = Q V_{\mathrm{source}}$.
La \cle{bande passante} (bande de fréquences à -3 dB, puissances demi) est $\Delta f = f_0 / Q$.
\end{propriete}

\courssub{Résonance parallèle}

\begin{definition}[Résonance parallèle]
Pour un circuit parallèle RLC (résistance $R$ en parallèle avec $L$ et $C$ idéaux), la résonance se produit lorsque les susceptances se compensent : $B_C = B_L \iff \omega_0 C = 1/(\omega_0 L)$. La fréquence est la même : $f_0 = 1/(2\pi\sqrt{LC})$.
L'admittance est minimale et purement conductive : $Y_{\min} = 1/R$. L'\cle{impédance est maximale} : $Z_{\max} = R$.
\end{definition}

\begin{propriete}[Facteur de qualité $Q$ en parallèle]
\[
Q = \frac{B_C}{G} = \omega_0 C R = \frac{R}{\omega_0 L} = R\sqrt{\frac{C}{L}}.
\]
Les courants dans $L$ et $C$ sont $Q$ fois le courant de ligne : $I_L = I_C = Q I_{\mathrm{source}}$ (ils sont en opposition de phase et circulent en boucle locale).
\end{propriete}

\begin{exemplebox}[Circuit parallèle RLC]
$R = 100\ \Omega$, $L = 20\ \mathrm{mH}$, $C = 5\ \mu\mathrm{F}$.
$f_0 = \dfrac{1}{2\pi\sqrt{0.020 \times 5\times10^{-6}}} = 503\ \mathrm{Hz}$.
$\omega_0 = 3162\ \mathrm{rad/s}$.
$Q = \omega_0 C R = 3162 \times 5\times10^{-6} \times 100 = 1.58$.
Bande passante $\Delta f = 503 / 1.58 = 318\ \mathrm{Hz}$.
Fréquences de coupure : $f_1 \approx 344\ \mathrm{Hz}$, $f_2 \approx 662\ \mathrm{Hz}$.
\end{exemplebox}

\begin{popfigure}
\begin{tikzpicture}
\begin{axis}[
    width=\linewidth, height=7cm,
    xlabel={Fréquence $f$ (Hz)}, ylabel={$|Z|$ ($\Omega$) / $|Y|$ (S)},
    xmin=0, xmax=1000, ymin=0, ymax=110,
    xtick={344,503,662}, xticklabels={$f_1$,$f_0$,$f_2$},
    ytick={100}, yticklabels={$R$},
    axis lines=middle, enlargelimits=0.1,
    clip=false, legend pos=north east
]
\addplot[domain=0:1000, samples=300, popblue, thick] {1/sqrt(max(0,(1/100)^2 + (2*pi*x*5e-6 - 1/(2*pi*x*0.02))^2))};
\addlegendentry{$|Z|$ parallèle (max à $f_0$)}
\addplot[domain=0:1000, samples=300, poporange, thick, dashed] {sqrt(max(0,10^2 + (2*pi*x*0.02 - 1/(2*pi*x*5e-6))^2))};
\addlegendentry{$|Z|$ série (min à $f_0$)}
\addplot[dashed, popmuted] coordinates {(503,0) (503,100)};
\addplot[dashed, popmuted] coordinates {(344,0) (344,70.7)};
\addplot[dashed, popmuted] coordinates {(662,0) (662,70.7)};
\node[above right, popblue] at (axis cs:503,100) {Résonance parallèle};
\node[below right, poporange] at (axis cs:503,10) {Résonance série};
\end{axis}
\end{tikzpicture}
\legende{Comparaison des courbes d'impédance : série (minimum) vs parallèle (maximum) à la résonance.}
\end{popfigure}

\courssub{Puissance en régime sinusoïdal}

\begin{definition}[Puissances instantanée, active, réactive, apparente]
Pour une tension $v(t) = V_m\sin(\omega t)$ et un courant $i(t) = I_m\sin(\omega t - \varphi)$ :
\begin{itemize}
\item \textbf{Puissance instantanée :} $p(t) = v(t)i(t) = V_{\mathrm{rms}}I_{\mathrm{rms}}\cos\varphi - V_{\mathrm{rms}}I_{\mathrm{rms}}\cos(2\omega t - \varphi)$.
\item \textbf{Puissance active (réelle) :} $P = V_{\mathrm{rms}}I_{\mathrm{rms}}\cos\varphi$ (Watts, W). Énergie réellement dissipée/convertie.
\item \textbf{Puissance réactive :} $Q = V_{\mathrm{rms}}I_{\mathrm{rms}}\sin\varphi$ (Volt-Ampères Réactifs, VAR). Énergie stockée/restituée par $L$ et $C$. $Q>0$ pour inductif (retard), $Q<0$ pour capacitif (avance).
\item \textbf{Puissance apparente :} $S = V_{\mathrm{rms}}I_{\mathrm{rms}}$ (Volt-Ampères, VA).
\end{itemize}
Relation : $S^2 = P^2 + Q^2$. \cle{Facteur de puissance} : $\mathrm{PF} = \cos\varphi = P/S$.
\end{definition}

\begin{aretenir}[Triangle des puissances]
\begin{center}
\begin{tikzpicture}[scale=0.8]
\draw[->, thick] (0,0) -- (5,0) node[right] {$P$ (W)};
\draw[->, thick] (0,0) -- (0,4) node[above] {$Q$ (VAR)};
\draw[thick, popblue] (0,0) -- (4,2) node[midway, above right] {$S$ (VA)};
\draw[dashed] (4,0) -- (4,2);
\draw[dashed] (0,2) -- (4,2);
\node at (2,-0.5) {$\varphi$};
\draw (0.5,0) arc (0:26.6:0.5);
\end{tikzpicture}
\end{center}
$\cos\varphi = P/S$, $\tan\varphi = Q/P$. $\varphi>0$ (inductif) $\Rightarrow$ PF \emph{retardant} (lagging). $\varphi<0$ (capacitif) $\Rightarrow$ PF \emph{avancé} (leading).
\end{aretenir}

\begin{methode}[Correction du facteur de puissance]
Pour améliorer le PF d'une charge inductive (usine, moteurs) vers une valeur cible $\cos\varphi_{\mathrm{target}}$ :
\begin{enumerate}
\item Calculer $P$ total et $Q$ total initial (somme algébrique des $Q$ de chaque charge : $Q_L>0$, $Q_C<0$).
\item Déterminer $Q_{\mathrm{target}} = P \tan(\arccos(\mathrm{PF}_{\mathrm{target}}))$.
\item Puissance réactive capacitive à ajouter : $Q_C = Q_{\mathrm{initial}} - Q_{\mathrm{target}}$ (si $Q_C > 0$, ajouter des condensateurs ; si $Q_C < 0$, le PF est déjà meilleur que la cible, il faudrait ajouter des inductances).
\item Capacité nécessaire (monophasé, tension $V$) : $C = \dfrac{Q_C}{\omega V^2}$.
\end{enumerate}
\end{methode}

\begin{exoresolu}[Correction de facteur de puissance — Usine]
Une usine consomme :
\begin{itemize}
\item Chauffage : $P_1 = 20\ \mathrm{kW}$, $Q_1 = 0$.
\item Moteurs : $P_2 = 15\ \mathrm{kW}$, $\mathrm{PF}_2 = 0.7$ retard $\Rightarrow \varphi_2 = 45.57^\circ$, $\tan\varphi_2 = 1.02$. $Q_2 = 15.3\ \mathrm{kVAR}$.
\item Condensateurs existants : $Q_3 = -10\ \mathrm{kVAR}$.
\end{itemize}
$P = 35\ \mathrm{kW}$. $Q = 0 + 15.3 - 10 = 5.3\ \mathrm{kVAR}$ (retard).
$S = \sqrt{35^2 + 5.3^2} = 35.4\ \mathrm{kVA}$.
$\mathrm{PF} = 35/35.4 = 0.989$ retard.
Cible $\mathrm{PF} = 0.95$ retard $\Rightarrow \varphi_{\mathrm{target}} = 18.19^\circ$, $\tan\varphi_{\mathrm{target}} = 0.329$.
$Q_{\mathrm{target}} = 35 \times 0.329 = 11.5\ \mathrm{kVAR}$ (retard).
$Q_{\mathrm{initial}} = 5.3\ \mathrm{kVAR} < Q_{\mathrm{target}}$.
\textbf{Conclusion :} Le PF actuel (0.989) est \emph{déjà meilleur} que la cible (0.95). Aucune capacité supplémentaire n'est nécessaire ; au contraire, il faudrait \emph{retirer} environ $6.2\ \mathrm{kVAR}$ de capacitif (ou ajouter de l'inductif) pour dégrader volontairement le PF à 0.95.
(Si l'usine n'avait pas les 10 kVAR initiaux : $Q_{\mathrm{initial}} = 15.3\ \mathrm{kVAR}$, il faudrait ajouter $Q_C = 15.3 - 11.5 = 3.8\ \mathrm{kVAR}$ capacitifs).
\end{exoresolu}

\begin{exercice}[Résonance et Puissance]
\begin{enumerate}
\item Un circuit série $R=5\ \Omega$, $L=10\ \mathrm{mH}$, $C=100\ \mu\mathrm{F}$. Calculer $f_0$, $Q$, $\Delta f$.
\item Un circuit parallèle $R=50\ \Omega$, $L=100\ \mathrm{mH}$, $C=10\ \mu\mathrm{F}$. Calculer $f_0$, $Q$, $Z_{\max}$.
\item Une charge absorbe $P=5\ \mathrm{kW}$ sous $230\ \mathrm{V}$ avec $\mathrm{PF}=0.6$ retard. Calculer $S$, $Q$, et la capacité à placer en parallèle pour obtenir $\mathrm{PF}=0.95$ retard ($f=50\ \mathrm{Hz}$).
\end{enumerate}
\end{exercice}

\begin{corrige}[Exercise 2 — Résonance et Puissance]
\begin{enumerate}
\item $f_0 = 1/(2\pi\sqrt{0.01 \times 10^{-4}}) = 159\ \mathrm{Hz}$. $\omega_0 = 1000\ \mathrm{rad/s}$. $Q = \omega_0 L/R = 1000 \times 0.01 / 5 = 2$. $\Delta f = 159/2 = 79.5\ \mathrm{Hz}$.
\item $f_0 = 1/(2\pi\sqrt{0.1 \times 10^{-5}}) = 159\ \mathrm{Hz}$. $Q = R\sqrt{C/L} = 50 \times \sqrt{10^{-5}/0.1} = 50 \times 0.01 = 0.5$. $Z_{\max} = R = 50\ \Omega$.
\item $S = P/\cos\varphi = 5/0.6 = 8.33\ \mathrm{kVA}$. $Q = S\sin\varphi = 8.33 \times 0.8 = 6.67\ \mathrm{kVAR}$.
Cible : $\cos\varphi' = 0.95 \Rightarrow \tan\varphi' = 0.329$. $Q' = 5 \times 0.329 = 1.64\ \mathrm{kVAR}$.
$Q_C = 6.67 - 1.64 = 5.03\ \mathrm{kVAR}$.
$C = Q_C / (\omega V^2) = 5030 / (314 \times 230^2) = 303\ \mu\mathrm{F}$.
\end{enumerate}
\end{corrige}

\coursec{Transformateurs : modèle idéal et pertes réelles}

\courssub{Transformateur idéal}

\begin{definition}[Transformateur idéal]
Un transformateur idéal est un dispositif statique à deux enroulements (primaire $N_p$ tours, secondaire $N_s$ tours) couplés par un flux magnétique parfait (pas de fuite, perméabilité infinie, pas de pertes). Les relations fondamentales sont :
\[
\frac{V_p}{V_s} = \frac{N_p}{N_s} = \frac{I_s}{I_p} = a \quad (\text{rapport de transformation}).
\]
Puissance conservée : $V_p I_p = V_s I_s$ (valeurs efficaces ou amplitudes).
\end{definition}

\begin{propriete}[Référence d'impédance]
Une impédance $Z_s$ connectée au secondaire est "vue" du primaire comme $Z_p = a^2 Z_s$. Réciproquement, $Z_s = Z_p / a^2$.
Ceci permet de ramener tout le circuit du côté primaire (ou secondaire) pour l'analyse.
\end{propriete}

\courssub{Transformateur réel : modèle équivalent et pertes}

\begin{definition}[Modèle équivalent référencé au primaire]
Le transformateur réel se modélise par :
\begin{itemize}
\item Résistances des enroulements : $R_1$ (primaire), $R_2'$ (secondaire ramené au primaire : $R_2' = a^2 R_2$).
\item Réactances de fuite : $X_1$, $X_2'$ ($X_2' = a^2 X_2$).
\item Branche d'excitation (noyau) : $R_m$ (pertes fer) en parallèle avec $X_m$ (magnétisation). Souvent simplifiée en $R_m$ seul pour les calculs de régulation/rendement.
\end{itemize}
\end{definition}

\begin{popfigure}
\begin{tikzpicture}[scale=0.9, european]
\draw (0,0) to[R=$R_1$] (2,0) to[L=$X_1$] (4,0) node[above] {$A$};
\draw (4,0) -- (4,2) to[R=$R_m$] (8,2) to[L=$X_m$] (8,0) -- (4,0);
\draw (4,0) to[R=$R_2'$] (6,0) to[L=$X_2'$] (8,0) node[above] {$B$};
\draw (0,0) -- (0,-1) node[below] {$V_1$};
\draw (8,0) -- (8,-1) node[below] {$V_2'$};
\draw (0,-0.5) to[battery1, l_={$V_1$}] (0,-1.5) -- (8,-1.5) -- (8,-0.5);
\end{tikzpicture}
\legende{Modèle équivalent du transformateur réel ramené au primaire.}
\end{popfigure}

\begin{propriete}[Pertes dans un transformateur réel]
\begin{enumerate}
\item \textbf{Pertes fer (ou pertes à vide) $P_{\mathrm{fer}}$ :} Dues au cycle d'hystérésis et aux courants de Foucault dans le noyau. Dépendent de la tension et de la fréquence (quasi constantes en charge). Minimisées par : tôles minces isolées (réduisent Foucault), acier au silicium / grains orientés (réduisent hystérésis).
\item \textbf{Pertes cuivre (ou pertes en charge) $P_{\mathrm{cu}}$ :} Effet Joule dans les enroulements : $P_{\mathrm{cu}} = I_1^2 R_1 + I_2^2 R_2 = I_1^2 R_{\mathrm{eq}}$. Dépendent du carré du courant (donc de la charge au carré). Minimisées par : conducteurs de section adéquate (cuivre/aluminium), bon refroidissement.
\end{enumerate}
Pertes totales : $P_{\mathrm{pertes}} = P_{\mathrm{fer}} + P_{\mathrm{cu}}$.
\end{propriete}

\begin{definition}[Rendement et Régulation]
\begin{itemize}
\item \textbf{Rendement :} $\eta = \dfrac{P_{\mathrm{sortie}}}{P_{\mathrm{entrée}}} \times 100\% = \dfrac{P_{\mathrm{sortie}}}{P_{\mathrm{sortie}} + P_{\mathrm{fer}} + P_{\mathrm{cu}}} \times 100\%$.
Le rendement maximal se produit lorsque $P_{\mathrm{cu}} = P_{\mathrm{fer}}$.
\item \textbf{Régulation de tension (pourcentage) :} Variation de la tension secondaire entre vide et charge nominale, rapportée à la tension nominale.
\[
\%\mathrm{Reg} = \frac{V_{2,\mathrm{vide}} - V_{2,\mathrm{charge}}}{V_{2,\mathrm{nominale}}} \times 100\% \approx \frac{I_1 (R_{\mathrm{eq}}\cos\varphi + X_{\mathrm{eq}}\sin\varphi)}{V_1} \times 100\%.
\]
$\varphi$ est l'angle du facteur de puissance de la charge ($+$ pour retard/inductif, $-$ pour avance/capacitif).
\end{itemize}
\end{definition}

\begin{exoresolu}[Transformateur 5 kVA — Régulation et Rendement]
Données : $S=5\ \mathrm{kVA}$, $V_1=230\ \mathrm{V}$, $V_2=115\ \mathrm{V}$, $f=50\ \mathrm{Hz}$.
$R_1=0.5\ \Omega$, $R_2=0.125\ \Omega$, $P_{\mathrm{fer}}=80\ \mathrm{W}$.
$X_1=1.2\ \Omega$, $X_2'=0.3\ \Omega$ (déjà ramenés). Charge nominale, $\mathrm{PF}=0.8$ retard.
$a = V_1/V_2 = 2$.
\begin{enumerate}
\item \textbf{Paramètres équivalents (primaire) :}
$R_{\mathrm{eq}} = R_1 + a^2 R_2 = 0.5 + 4 \times 0.125 = 1.00\ \Omega$.
$X_{\mathrm{eq}} = X_1 + X_2' = 1.2 + 0.3 = 1.50\ \Omega$.
\item \textbf{Régulation \% :}
$I_1 = S/V_1 = 5000/230 = 21.74\ \mathrm{A}$.
$\cos\varphi=0.8$, $\sin\varphi=0.6$.
$\Delta V = I_1(R_{\mathrm{eq}}\cos\varphi + X_{\mathrm{eq}}\sin\varphi) = 21.74(1.0\times0.8 + 1.5\times0.6) = 21.74 \times 1.7 = 36.96\ \mathrm{V}$.
$\%\mathrm{Reg} = (36.96/230) \times 100\% = 16.1\%$.
\item \textbf{Pertes cuivre charge nominale :}
$P_{\mathrm{cu}} = I_1^2 R_{\mathrm{eq}} = 21.74^2 \times 1.00 = 472.6\ \mathrm{W}$.
\item \textbf{Rendement charge nominale :}
$P_{\mathrm{out}} = S\cos\varphi = 5000 \times 0.8 = 4000\ \mathrm{W}$.
$P_{\mathrm{pertes}} = 80 + 472.6 = 552.6\ \mathrm{W}$.
$\eta = 4000 / (4000 + 552.6) = 87.9\%$.
\item \textbf{Rendement à 50\% charge (même PF) :}
$I_1' = 10.87\ \mathrm{A}$. $P_{\mathrm{cu}}' = (0.5)^2 \times 472.6 = 118.2\ \mathrm{W}$.
$P_{\mathrm{out}}' = 2000\ \mathrm{W}$. $P_{\mathrm{pertes}}' = 80 + 118.2 = 198.2\ \mathrm{W}$.
$\eta' = 2000 / 2198.2 = 91.0\%$.
\textit{Le rendement augmente car les pertes cuivre chutent au carré du courant (25\% de la valeur nominale) tandis que les pertes fer restent constantes. On se rapproche de la condition $P_{\mathrm{cu}} = P_{\mathrm{fer}}$ pour le rendement maximal.}
\end{enumerate}
\end{exoresolu}

\begin{exercice}[Transformateurs]
\begin{enumerate}
\item Un transformateur idéal a $N_p=800$, $N_s=200$, $V_p=240\ \mathrm{V}$. Trouver $V_s$.
\item Un transformateur 10 kVA, 2300/230 V a $R_1=2\ \Omega$, $R_2=0.02\ \Omega$, $P_{\mathrm{fer}}=120\ \mathrm{W}$. Calculer le rendement à pleine charge $\mathrm{PF}=1$ et à 50\% charge $\mathrm{PF}=0.8$ retard.
\item Un transformateur 460/12 V alimente une charge de 24 W (résistive). $N_p=460$, $P_{\mathrm{fer}}=2\ \mathrm{W}$, $P_{\mathrm{cu}}=1.5\ \mathrm{W}$. Trouver $N_s$, $I_s$, $I_p$ (idéal), $\eta$.
\end{enumerate}
\end{exercice}

\begin{corrige}[Exercise 3 — Transformateurs]
\begin{enumerate}
\item $V_s = V_p \times N_s/N_p = 240 \times 200/800 = 60\ \mathrm{V}$.
\item $a=10$. $R_{\mathrm{eq}} = 2 + 100 \times 0.02 = 4\ \Omega$.
Pleine charge PF=1 : $I_1 = 10000/2300 = 4.35\ \mathrm{A}$. $P_{\mathrm{cu}} = 4.35^2 \times 4 = 75.6\ \mathrm{W}$. $P_{\mathrm{out}}=10000\ \mathrm{W}$. $\eta = 10000/(10000+120+75.6) = 98.1\%$.
50\% charge PF=0.8 : $I_1' = 2.17\ \mathrm{A}$. $P_{\mathrm{cu}}' = 0.25 \times 75.6 = 18.9\ \mathrm{W}$. $P_{\mathrm{out}}' = 0.5 \times 10000 \times 0.8 = 4000\ \mathrm{W}$. $\eta' = 4000/(4000+120+18.9) = 96.6\%$.
\item $N_s = 460 \times 12/230 = 24$ tours. $I_s = 24/12 = 2.00\ \mathrm{A}$. $I_p = 2.00 \times 24/460 = 0.104\ \mathrm{A}$. $P_{\mathrm{out}}=24\ \mathrm{W}$, $P_{\mathrm{in}}=24+2+1.5=27.5\ \mathrm{W}$. $\eta = 24/27.5 = 87.3\%$.
\end{enumerate}
\end{corrige}

\coursec{Redressement et lissage des signaux alternatifs}

\courssub{Redresseurs : principes et caractéristiques}

\begin{definition}[Redressement]
Le \cle{redressement} consiste à convertir une tension alternative en tension unidirectionnelle (composante continue + ripple) à l'aide de diodes.
\begin{itemize}
\item \textbf{Demi-onde :} Une diode. Conduit 1 demi-période sur 2. Fréquence de ripple $f_r = f$. Tension continue moyenne (sans condensateur) $V_{\mathrm{dc}} = V_m/\pi = 0.318 V_m$.
\item \textbf{Double alternance (pont de diodes / bridge) :} 4 diodes. Conduit les 2 demi-périodes. Fréquence de ripple $f_r = 2f$. $V_{\mathrm{dc}} = 2V_m/\pi = 0.637 V_m$.
\item \textbf{Double alternance (centre-tap) :} 2 diodes + secondaire à prise médiane. $f_r = 2f$, $V_{\mathrm{dc}} = 0.637 V_m$. PIV diode = $2V_m$.
\end{itemize}
\end{definition}

\begin{propriete}[Tension de crête inverse (PIV / Peak Inverse Voltage)]
Tension maximale qu'une diode doit supporter en inverse sans claquage.
\begin{itemize}
\item Demi-onde : $\mathrm{PIV} = V_m$ (secondaire) ou $2V_m$ si condensateur de lissage charge à $V_m$ et secondaire va à $-V_m$.
\item Pont (Bridge) : $\mathrm{PIV} = V_m$ (tension de crête du secondaire).
\item Centre-tap : $\mathrm{PIV} = 2V_m$.
\end{itemize}
\end{propriete}

\courssub{Lissage par condensateur}

\begin{definition}[Condensateur de lissage]
Un condensateur $C$ placé en parallèle sur la charge $R_L$ se charge au pic de tension $V_m$ et se décharge à travers $R_L$ pendant la non-conduction des diodes.
\begin{itemize}
\item \textbf{Ripple (ondulation) crête-à-crête $V_{\mathrm{pp}}$ (approximation décharge linéaire, $V_{\mathrm{pp}} \ll V_m$) :}
\[
V_{\mathrm{pp}} \approx \frac{I_{\mathrm{dc}}}{f_r C} = \frac{V_{\mathrm{dc}}}{R_L f_r C}.
\]
\item \textbf{Tension continue moyenne :}
\[
V_{\mathrm{dc}} \approx V_m - \frac{V_{\mathrm{pp}}}{2}.
\]
\item \textbf{Facteur de ripple :} $r = V_{\mathrm{pp}} / V_{\mathrm{dc}}$ (ou $V_{\mathrm{rms,ripple}}/V_{\mathrm{dc}}$).
\end{itemize}
\end{definition}

\begin{methode}[Dimensionnement du condensateur de lissage]
Données : $V_s$ (RMS secondaire), $f$, $I_{\mathrm{dc}}$ (ou $R_L$), $V_{\mathrm{pp,max}}$.
\begin{enumerate}
\item $V_m = \sqrt{2} V_s$.
\item Fréquence ripple $f_r = f$ (demi-onde) ou $2f$ (pont/centre-tap).
\item $C_{\min} = \dfrac{I_{\mathrm{dc}}}{f_r V_{\mathrm{pp,max}}}$.
\item Vérifier $\mathrm{PIV}$ des diodes choisis.
\item Estimer $V_{\mathrm{dc}} \approx V_m - V_{\mathrm{pp}}/2$.
\end{enumerate}
\end{methode}

\begin{exemplebox}[Pont redresseur avec lissage]
$V_s = 12\ \mathrm{V}$, $f=50\ \mathrm{Hz}$, $I_{\mathrm{dc}} = 0.5\ \mathrm{A}$, $V_{\mathrm{pp}} < 1\ \mathrm{V}$.
$V_m = 16.97\ \mathrm{V}$. $\mathrm{PIV} = 17.0\ \mathrm{V}$.
$f_r = 100\ \mathrm{Hz}$.
$C \ge \dfrac{0.5}{100 \times 1} = 0.005\ \mathrm{F} = 5000\ \mu\mathrm{F}$.
$V_{\mathrm{dc}} \approx 16.97 - 0.5 = 16.47\ \mathrm{V}$.
\end{exemplebox}

\begin{exoresolu}[Comparaison Demi-onde vs Pont — Essai type GCE]
Secondaire $12\ \mathrm{V}$ RMS, $f=50\ \mathrm{Hz}$. Charge $R_L = 10\ \Omega$, $C = 1000\ \mu\mathrm{F}$.
$V_m = 17.0\ \mathrm{V}$.
\textbf{(a) Schémas :} Voir figure. (Demi-onde : 1 diode + C // R. Pont : 4 diodes en pont + C // R).
\textbf{(b) PIV :} Demi-onde : $2V_m = 34\ \mathrm{V}$. Pont : $V_m = 17\ \mathrm{V}$.
\textbf{(c) Tensions continues et ripple (approximation $V_{\mathrm{dc}} = V_m - V_{\mathrm{pp}}/2$, $V_{\mathrm{pp}} = V_{\mathrm{dc}}/(R_L f_r C)$) :}
\begin{itemize}
\item \textbf{Demi-onde} ($f_r = 50\ \mathrm{Hz}$) : $2R_L f C = 1$.
$V_{\mathrm{dc}} = V_m / (1 + 1/(2R_L f C)) = 17 / (1+1) = 8.5\ \mathrm{V}$.
$V_{\mathrm{pp}} = V_{\mathrm{dc}} / (R_L f C) = 8.5 / (10 \times 50 \times 0.001) = 17\ \mathrm{V}$.
\item \textbf{Pont} ($f_r = 100\ \mathrm{Hz}$) : $4R_L f C = 2$.
$V_{\mathrm{dc}} = V_m / (1 + 1/(4R_L f C)) = 17 / (1+0.5) = 11.3\ \mathrm{V}$.
$V_{\mathrm{pp}} = V_{\mathrm{dc}} / (2 R_L f C) = 11.3 / (2 \times 10 \times 50 \times 0.001) = 11.3\ \mathrm{V}$.
\end{itemize}
\textbf{(d) Comparaison :}
\begin{center}
\begin{tabularx}{\linewidth}{|l|X|X|}\hline
\textbf{Paramètre} & \textbf{Demi-onde} & \textbf{Pont (Bridge)} \\ \hline
$V_{\mathrm{dc}}$ moyenne & $8.5\ \mathrm{V}$ & $11.3\ \mathrm{V}$ (meilleur) \\ \hline
Ripple $V_{\mathrm{pp}}$ & $17\ \mathrm{V}$ & $11.3\ \mathrm{V}$ (meilleur) \\ \hline
PIV diode & $34\ \mathrm{V}$ & $17\ \mathrm{V}$ (meilleur) \\ \hline
Utilisation transfo & Mauvaise (courant continu dans secondaire $\to$ saturation) & Excellente (pas de composante DC dans secondaire) \\ \hline
\end{tabularx}
\end{center}
\textbf{Conclusion :} Le pont est supérieur : tension plus élevée, ripple plus faible, PIV plus faible, meilleure utilisation du transformateur. C'est le standard industriel.
\end{exoresolu}

\begin{popfigure}
\begin{tikzpicture}[scale=0.8, european]
% Half-wave rectifier
\begin{scope}[xshift=0cm]
\node[above] at (2,3) {\textbf{Demi-onde}};
\draw (0,0) to[sV, l_={$v_s$}] (0,2);
\draw (0,2) to[D] (2,2);
\draw (2,2) to[C=$C$, -*] (2,0) node[ground]{};
\draw (2,2) to[R=$R_L$, -*] (2,0);
\draw (0,0) -- (2,0);
\end{scope}
% Bridge rectifier
\begin{scope}[xshift=6cm]
\node[above] at (2,3) {\textbf{Pont (Bridge)}};
% AC Source
\draw (0,1) to[sV, l_={$v_s$}] (0,3);
% Top left diode (D1)
\draw (0,3) to[D] (2,3);
% Bottom left diode (D2)
\draw (0,1) to[D] (2,1);
% Top right diode (D3) - cathode at top
\draw (2,3) to[D] (2,1);
% Bottom right diode (D4) - anode at bottom? No, standard bridge:
% D1: anode left, cathode top
% D2: anode bottom, cathode left
% D3: anode bottom, cathode right
% D4: anode right, cathode top
% Let's draw standard bridge:
% Node A (left top) --D1--> Node B (top) --D4(cathode top)--> Node C (right top) --D3(anode bottom)--> Node D (bottom) --D2(cathode left)--> Node A
% Actually simpler:
\draw (0,3) to[D] (2,3); % D1
\draw (2,1) to[D] (2,3); % D4 (cathode up)
\draw (0,1) to[D] (2,1); % D2 (anode up? No, cathode left)
% Correct bridge:
% Left node (0,2) -- but source takes (0,1) to (0,3).
% Let's define nodes:
% Source between (0,1) and (0,3).
% D1 between (0,3) and (2,3) (Anode left, Cathode right)
% D2 between (2,1) and (0,1) (Anode right, Cathode left) -> to[D] (0,1) -- (2,1) draws anode left. Need anode right.
% Use [D, invert] or swap coordinates.
\draw (2,1) to[D] (0,1); % D2: Anode at (2,1), Cathode at (0,1)
\draw (2,3) to[D] (2,1); % D3: Anode at (2,1), Cathode at (2,3) -- wait.
% Standard Bridge:
% ~ (0,3) --D1--> + (2,3)
% ~ (0,1) <--D2-- + (2,1)  (D2 cathode at ~, anode at +) -> so (0,1) to[D, invert] (2,1) or (2,1) to[D] (0,1)
% + (2,3) <--D4-- - (2,1)  (D4 cathode at +, anode at -) -> (2,1) to[D] (2,3)
% Load between + (2,3) and - (2,1)

% Redrawing cleanly:
\draw (0,3) to[D] (2,3);          % D1
\draw (2,1) to[D] (0,1);          % D2 (anode at 2,1, cathode at 0,1)
\draw (2,1) to[D] (2,3);          % D3 (anode at 2,1, cathode at 2,3)
\draw (0,3) to[D, invert] (0,1);  % D4 (cathode at 0,3, anode at 0,1) -- wait, D4 connects + to - on right side? No.
% Bridge has 4 diodes. Two on left (source side), two on right (load side).
% Left leg: D1 (top, anode~ cathode+), D2 (bottom, anode+ cathode~)
% Right leg: D3 (top, cathode+ anode-), D4 (bottom, cathode- anode+)
% Nodes: ~1 (0,3), ~2 (0,1), + (2,3), - (2,1)
% D1: ~1 -> +  : (0,3) to[D] (2,3)
% D2: + -> ~2  : (2,1) to[D] (0,1)  (Anode +, Cathode ~2)
% D3: - -> +   : (2,1) to[D] (2,3)  (Anode -, Cathode +)
% D4: ~1 -> -  : (0,3) to[D, invert] (2,1)? No, D4 connects ~1 to -.
% D4: Cathode ~1, Anode -  => Current flows - to ~1. So (2,1) to[D] (0,3).
\draw (2,1) to[D] (0,3);          % D4 (Anode -, Cathode ~1)

% Load and Capacitor
\draw (2,3) to[C=$C$, -*] (2,1) node[ground]{};
\draw (2,3) to[R=$R_L$, -*] (2,1);
\end{scope}
\end{tikzpicture}
\legende{Schéma simplifié : Redresseur demi-onde (gauche) et pont de diodes (droite) avec condensateur de lissage $C$ et charge $R_L$.}
\end{popfigure}

\begin{savaistu}[Applications au Cameroun]
Au Cameroun, le réseau ENEO distribue du 230 V / 50 Hz monophasé (400 V triphasé). Les onduleurs (inverters) solaires très répandus (ex. quartiers non raccordés, backup coupures) utilisent des transformateurs haute fréquence (ferrite) pour l'isolation galvanique et la conversion DC/AC, suivis de filtres LC pour sinusoïde pure. Les chargeurs de téléphones (adaptateurs USB) sont des alimentations à découpage (SMPS) : redressement pont + condensateur 400 V + convertisseur DC/DC haute fréquence + petit transformateur + redressement sortie + lissage. Comprendre le redressement "linéaire" (50 Hz) reste la base pédagogique pour aborder l'électronique de puissance moderne.
\end{savaistu}

\begin{exercice}[Redressement et Lissage]
\begin{enumerate}
\item Un redresseur demi-onde est alimenté par un transformateur 230/12 V, $f=50\ \mathrm{Hz}$. Charge $R=100\ \Omega$. Sans condensateur : calculer $V_{\mathrm{dc}}$, $I_{\mathrm{dc}}$, PIV diode.
\item Même transformateur, pont de diodes, $C=470\ \mu\mathrm{F}$, $R=100\ \Omega$. Estimer $V_{\mathrm{dc}}$, $V_{\mathrm{pp}}$, PIV.
\item On veut obtenir $12\ \mathrm{V}$ DC stabilisé à partir du 230 V secteur avec un ripple $< 100\ \mathrm{mV}$ pour $I_{\max}=1\ \mathrm{A}$. Proposer un schéma (transformateur, redresseur, condensateur) et calculer $C$ et PIV. (Négliger la chute régulateur).
\end{enumerate}
\end{exercice}

\begin{corrige}[Exercise 4 — Redressement et Lissage]
\begin{enumerate}
\item $V_m = \sqrt{2}\times 12 = 16.97\ \mathrm{V}$. $V_{\mathrm{dc}} = V_m/\pi = 5.40\ \mathrm{V}$. $I_{\mathrm{dc}} = 54.0\ \mathrm{mA}$. PIV = $V_m = 17.0\ \mathrm{V}$ (sans C) ou $2V_m = 34\ \mathrm{V}$ (avec C gros).
\item Pont : $V_m = 16.97\ \mathrm{V}$. PIV = $17.0\ \mathrm{V}$.
$f_r = 100\ \mathrm{Hz}$. Itération : $V_{\mathrm{dc}} \approx V_m = 17\ \mathrm{V}$, $I \approx 0.17\ \mathrm{A}$.
$V_{\mathrm{pp}} = I/(f_r C) = 0.17/(100 \times 470\times10^{-6}) = 3.6\ \mathrm{V}$.
$V_{\mathrm{dc}} \approx 17 - 1.8 = 15.2\ \mathrm{V}$. Recalculer $I = 0.152\ \mathrm{A}$, $V_{\mathrm{pp}} = 3.2\ \mathrm{V}$, $V_{\mathrm{dc}} = 15.4\ \mathrm{V}$.
\item Besoin $V_{\mathrm{dc}} \approx 12\ \mathrm{V}$ $\Rightarrow$ $V_m \approx 12 + V_{\mathrm{pp}}/2 \approx 12.05\ \mathrm{V}$. $V_s \approx 8.5\ \mathrm{V}$ RMS. Transfo 230/9 V ou 12 V (régulateur 7812 absorbera l'excès).
Pont redresseur. $f_r=100\ \mathrm{Hz}$. $C = I/(f_r V_{\mathrm{pp}}) = 1/(100 \times 0.1) = 100\ \mu\mathrm{F}$ (prendre $220\ \mu\mathrm{F}$/25V standard). PIV = $V_m \approx 12.7\ \mathrm{V}$ (diodes 1N4001 50V OK).
\end{enumerate}
\end{corrige}

\begin{attention}[Pièges fréquents à l'examen GCE]
\begin{itemize}
\item Confondre $V_{\mathrm{rms}}$ et $V_m$ (pic) dans les formules de PIV ou de lissage. Toujours : $V_m = \sqrt{2} V_{\mathrm{rms}}$.
\item Oublier le facteur 2 dans la fréquence de ripple pour le pont ($f_r = 2f$) et le centre-tap.
\item Utiliser la formule $V_{\mathrm{pp}} = I/(fC)$ pour la demi-onde au lieu de $I/(2fC)$ (ou $I/(f_r C)$ avec $f_r$ correct).
\item Négliger le signe de $Q$ (inductif $+$, capacitif $-$) dans les bilans de puissance réactive.
\item Confondre régulation (chute de tension interne) et rendement (rapport puissances).
\item Pour le transformateur, ramener correctement les impédances secondaires au primaire par $a^2$ (et non $a$).
\end{itemize}
\end{attention}

\begin{experience}[TP suggéré : Caractéristiques d'un transformateur réel]
\textbf{Matériel :} Transformateur 230/12 V (ex. 30 VA), variac (ou secteur direct avec prudence), multimètres (V, A, W si wattmètre dispo), charges résistives (lampes, rhéostats).
\textbf{Manipulations :}
\begin{enumerate}
\item \textbf{Essai à vide :} Primaire au 230 V, secondaire ouvert. Mesurer $V_1$, $I_0$, $P_0$, $V_2$. En déduire $P_{\mathrm{fer}} \approx P_0$, $R_m$, $X_m$, rapport $a = V_1/V_2$.
\item \textbf{Essai en court-circuit :} Secondaire court-circuité (fil épais), variac sur primaire pour obtenir $I_{1,\mathrm{nom}}$. Mesurer $V_{cc}$, $I_{cc}$, $P_{cc}$. En déduire $R_{\mathrm{eq}} = P_{cc}/I_{cc}^2$, $Z_{\mathrm{eq}} = V_{cc}/I_{cc}$, $X_{\mathrm{eq}} = \sqrt{Z_{\mathrm{eq}}^2 - R_{\mathrm{eq}}^2}$.
\item \textbf{Essai en charge :} Charger le secondaire avec $R_L$ connu. Mesurer $V_2$, $I_2$, $V_1$, $I_1$, $P_1$. Calculer $\eta$, $\%\mathrm{Reg}$. Comparer aux valeurs théoriques du modèle équivalent.
\end{enumerate}
\textbf{Rapport :} Schéma de câblage, tableaux de mesures, calculs, courbe $\eta$ vs $I_2$, courbe $V_2$ vs $I_2$, commentaires sur les écarts (échauffement, hypothèses modèle).
\end{experience}

\newpage

\coursec{Summary sheet}

\begin{popfigure}
\centering
\begin{tikzpicture}[
    mindmap,
    concept color=popblue,
    level 1/.style={level distance=4.5cm, sibling angle=60},
    level 2/.style={level distance=3cm, sibling angle=45},
    every node/.style={font=\small, text width=2.8cm, align=center},
    branch1/.style={concept color=poporange},
    branch2/.style={concept color=popgreen},
    branch3/.style={concept color=poppurple},
    branch4/.style={concept color=poppink},
    branch5/.style={concept color=popgold},
    branch6/.style={concept color=popblue}
]
\node[concept, text width=3.5cm] {Alternating Current\\ (PHY09)}
    child[branch1] {node[concept, align=center]{RMS Values\\ \& Impedance}
        child {node[concept, scale=0.7] {$V_{\text{rms}}=\frac{V_0}{\sqrt{2}}$}}
        child {node[concept, scale=0.7] {$Z=\sqrt{R^2+(X_L-X_C)^2}$}}
        child {node[concept, scale=0.7] {Phasor Diagrams}}
    }
    child[branch2] {node[concept] {Resonance}
        child {node[concept, scale=0.7] {$f_0=\frac{1}{2\pi\sqrt{LC}}$}}
        child {node[concept, scale=0.7] {$X_L=X_C$, $Z_{\min}=R$}}
        child {node[concept, scale=0.7] {Q-factor, Bandwidth}}
    }
    child[branch3] {node[concept] {Power in AC}
        child {node[concept, scale=0.7] {$P=V_{\text{rms}}I_{\text{rms}}\cos\phi$}}
        child {node[concept, scale=0.7] {Real, Reactive, Apparent}}
        child {node[concept, scale=0.7] {Power Factor Correction}}
    }
    child[branch4] {node[concept] {Transformers}
        child {node[concept, scale=0.7] {Ideal: $\frac{V_p}{V_s}=\frac{N_p}{N_s}$}}
        child {node[concept, scale=0.7] {Losses: Core \& Copper}}
        child {node[concept, scale=0.7] {Efficiency $\eta$}}
    }
    child[branch5] {node[concept] {Rectification}
        child {node[concept, scale=0.7] {Half / Full Wave}}
        child {node[concept, scale=0.7] {Bridge Rectifier}}
        child {node[concept, scale=0.7] {Output Waveforms}}
    }
    child[branch6] {node[concept] {Smoothing}
        child {node[concept, scale=0.7] {Capacitor Filter}}
        child {node[concept, scale=0.7] {Ripple Voltage $V_r$}}
        child {node[concept, scale=0.7] {Ripple Factor $\gamma$}}
    };
\end{tikzpicture}
\legende{Mind map of Chapter PHY09: Alternating Current}
\end{popfigure}

\begin{bilan}

\courssub{Key Definitions}
\begin{itemize}
    \item \cle{RMS Value (Effective Value)}: Equivalent DC value producing same heating effect. $V_{\text{rms}} = V_0/\sqrt{2}$, $I_{\text{rms}} = I_0/\sqrt{2}$ (sinusoidal only).
    \item \cle{Reactance}: Opposition by $L$ or $C$. $X_L = \omega L = 2\pi f L$ (inductive), $X_C = 1/(\omega C)$ (capacitive).
    \item \cle{Impedance ($Z$)}: Total opposition in RLC series. $Z = \sqrt{R^2 + (X_L - X_C)^2}\ \Omega$.
    \item \cle{Phase Angle ($\phi$)}: $\tan\phi = (X_L - X_C)/R$. Current leads voltage if $X_C > X_L$ ($\phi<0$), lags if $X_L > X_C$ ($\phi>0$).
    \item \cle{Resonance}: $X_L = X_C \Rightarrow Z_{\min}=R$, $I_{\max}$, $\phi=0$. Resonant frequency $f_0 = 1/(2\pi\sqrt{LC})$.
    \item \cle{Power Factor}: $\cos\phi = R/Z = P/S$. Ratio of real to apparent power.
    \item \cle{Transformer Turns Ratio}: $k = N_p/N_s = V_p/V_s = I_s/I_p$ (ideal).
    \item \cle{Ripple Factor ($\gamma$)}: $\gamma = V_{r(\text{rms})}/V_{\text{dc}}$ (measure of smoothing quality).
\end{itemize}

\courssub{Laws, Formulas \& Theorems (Must Know)}
\begin{itemize}
    \item \textbf{Series RLC}: $V = V_0\sin(\omega t)$, $I = I_0\sin(\omega t - \phi)$, $I_0 = V_0/Z$.
    \item \textbf{Voltage across components}: $V_R = I_0 R$ (in phase), $V_L = I_0 X_L$ (leads by $\pi/2$), $V_C = I_0 X_C$ (lags by $\pi/2$). Phasor sum: $V_0 = \sqrt{V_R^2 + (V_L - V_C)^2}$.
    \item \textbf{Average Power}: $P = V_{\text{rms}} I_{\text{rms}} \cos\phi = I_{\text{rms}}^2 R$. $Q = V_{\text{rms}} I_{\text{rms}} \sin\phi$ (reactive), $S = V_{\text{rms}} I_{\text{rms}}$ (apparent).
    \item \textbf{Resonance}: $f_0 = \frac{1}{2\pi\sqrt{LC}}$, $\omega_0 = 1/\sqrt{LC}$. Quality factor $Q = \frac{\omega_0 L}{R} = \frac{1}{\omega_0 C R} = \frac{f_0}{\Delta f}$.
    \item \textbf{Ideal Transformer}: $\frac{V_p}{V_s} = \frac{N_p}{N_s} = \frac{I_s}{I_p}$, $P_p = P_s$.
    \item \textbf{Real Transformer Losses}:
        \begin{itemize}
            \item Copper losses ($I^2R$ in windings) $\to$ minimised by thick/low-resistivity wire.
            \item Iron/Core losses: Hysteresis (laminated soft iron core) + Eddy currents (laminations).
            \item Flux leakage (tight winding, core design).
        \end{itemize}
    \item \textbf{Rectification}:
        \begin{itemize}
            \item Half-wave: $V_{\text{dc}} = V_m/\pi$, $V_{\text{rms}} = V_m/2$.
            \item Full-wave (centre-tap/bridge): $V_{\text{dc}} = 2V_m/\pi$, $V_{\text{rms}} = V_m/\sqrt{2}$.
        \end{itemize}
    \item \textbf{Smoothing (Capacitor filter, full-wave)}: Ripple voltage $V_r \approx \frac{I_{\text{dc}}}{2fC} = \frac{V_{\text{dc}}}{2fCR_L}$. Ripple factor $\gamma \approx \frac{1}{4\sqrt{3}fCR_L}$.
\end{itemize}

\courssub{Methods \& Procedures}
\begin{enumerate}
    \item \textbf{Impedance/Phasor Analysis}: Draw phasor diagram ($V_R$ horizontal, $V_L$ up, $V_C$ down). Find $Z$, $\phi$, $I_{\text{rms}}$, $P$.
    \item \textbf{Resonance Problems}: Identify $f_0$. At resonance: $Z=R$, $V_L=V_C$ (can be $>V_{\text{source}}$), $P_{\max}$.
    \item \textbf{Transformer Calculations}: Use turns ratio for $V_s$, $I_s$. Efficiency $\eta = P_{\text{out}}/P_{\text{in}} \times 100\%$. Account for losses if data given.
    \item \textbf{Rectifier/Smoothing Design}: Choose diode PIV rating ($>V_m$ for half-wave, $>2V_m$ for bridge). Select $C$ for required $V_r$: $C \approx \frac{I_{\text{dc}}}{2f V_r}$.
\end{enumerate}

\courssub{Common Mistakes \& Traps}
\begin{itemize}
    \item Using peak values in power formula $P=VI$ instead of RMS.
    \item Confusing $\omega_0$ (rad/s) and $f_0$ (Hz) in resonance formula.
    \item Forgetting $\phi$ sign: $\tan\phi = (X_L-X_C)/R$; $\phi>0$ means inductive (current lags).
    \item Assuming $V_p/V_s = N_p/N_s$ holds for non-ideal transformers under load (voltage regulation drops $V_s$).
    \item In smoothing, using half-wave formula $V_r = I_{\text{dc}}/(fC)$ for full-wave circuits (factor of 2 error).
    \item Connecting smoothing capacitor with wrong polarity (electrolytic).
    \item Forgetting units: $Z$ in $\Omega$, $C$ in F, $L$ in H, $f$ in Hz.
\end{itemize}

\courssub{GCE Examination Tips}
\begin{itemize}
    \item \textbf{Phasor diagrams} are frequently asked: draw to scale, label axes, show phase angle clearly.
    \item \textbf{Derivations}: Know how to derive $Z = \sqrt{R^2+(X_L-X_C)^2}$ from phasors. Derivation of transformer equation \emph{not} required, but statement and application are.
    \item \textbf{Transformer losses}: Explain \emph{how} each is minimised (laminations $\to$ eddy currents, soft iron $\to$ hysteresis, thick wire $\to$ copper loss).
    \item \textbf{Rectification}: Sketch input/output waveforms for half-wave, full-wave (centre-tap and bridge). Label $V_m$, $T$, zero crossings.
    \item \textbf{Smoothing}: Explain role of capacitor (charges at peaks, discharges through load). Calculate $C$ for a given ripple spec.
    \item \textbf{Significant figures}: Use 2 or 3 SF matching data. Always write units in final answer box.
    \item \textbf{Keywords}: "RMS", "Impedance", "Resonance", "Power Factor", "Turns Ratio", "Ripple Factor".
\end{itemize}

\end{bilan}

\findecours

\end{document}
